% આ ભાગનો પૂર્ણ સંપાદનીય પાઠ છે. શૈલી, ફોન્ટ, ચિત્રો અને પ્રકરણસંદર્ભ માટે સંપૂર્ણ સ્રોત ZIP તથા reader/full-reader.tex જુઓ. આ એકલી સ્વતંત્ર કમ્પાઇલ થતી મુખ્ય ફાઇલ નથી.

% BEGIN OLP-0341
\olpart{lam}{લૅમ્બડા કલન}
\phantomsection\label{guunit:OLP-0341}


\begin{editorial}
આ ભાગ લૅમ્બડા કલન વિશે છે. પ્રસ્તાવનાનું પ્રકરણ Jeremy Avigadની
નોંધો પર આધારિત છે; તેનો કેટલોક ભાગ હવે પુનરાવર્તિત બને છે અને પછીનાં
પ્રકરણોમાં આવરી લેવાયો છે. વાક્યરચના, ચર્ચ--રોસર ગુણધર્મ અને
લૅમ્બડા-વ્યાખ્યેયતા વિશેનાં પ્રકરણો Zesen Qianએ તેમની Mitacsની
ઉનાળાની ઇન્ટર્નશિપ દરમિયાન તૈયાર કર્યાં હતાં. તેમની હજી સમીક્ષા
અને સુધારણા કરવાની બાકી છે.
\end{editorial}

\begingroup
\makeatletter\def\ol@sectioncs{section}\makeatother

% BEGIN OLP-0342
\olchapter{lam}{int}{પ્રસ્તાવના}
\olchapterid{lam}{int}
\phantomsection\label{guunit:OLP-0342}


\begin{editorial}
આ પ્રકરણમાં લૅમ્બડા કલન વિશેની Jeremyની મૂળ સંક્ષિપ્ત નોંધો છે.
વિભાગોને એકત્ર કરવાની અને લૅમ્બડા-વ્યાખ્યેયતા વિશેની સામગ્રીને
તેના અલગ, વધુ વિગતવાર પ્રકરણની સામગ્રી સાથે જોડવાની જરૂર છે.
\end{editorial}

\begingroup
\makeatletter\def\ol@sectioncs{section}\makeatother

% BEGIN OLP-0343
\olfileid{lam}{int}{ovr}
\olsection{ઝાંખી}
\phantomsection\label{guunit:OLP-0343}


Alonzo Churchએ 1930ના દાયકાની શરૂઆતમાં લૅમ્બડા કલનની રચના મૂળે
રચનાત્મક તર્કશાસ્ત્રના પાયા તરીકે કરી હતી; સંગણનીય વિધેયોના
નિદર્શન તરીકે \emph{નહીં}. પરંતુ ટૂંક સમયમાં બતાવવામાં આવ્યું કે
તે સંગણનીયતાની બીજી વ્યાખ્યાઓ, જેમ કે ટ્યુરિંગ-સંગણનીય વિધેયો
અને આંશિક પુનરાવર્તી વિધેયો, સાથે સમકક્ષ છે. શરૂઆતમાં આ વાતથી
થોડું આશ્ચર્ય થયું હતું, અને તેથી આ લાક્ષણિકીકરણ વધુ રસપ્રદ બને છે.

લૅમ્બડા સંકેતન વિધેયને નામ આપ્યા વિના, તેને વ્યાખ્યાયિત કરતી
સંકેતી અભિવ્યક્તિ દ્વારા સીધો ઉલ્લેખ કરવાની અનુકૂળ રીત છે.
``$f(x) = x + 3$થી વ્યાખ્યાયિત વિધેય~$f$ લો'' એમ કહેવાને બદલે
``$f$ એ વિધેય~$\lambd[x][(x + 3)]$ છે'' એમ કહી શકાય. બીજા શબ્દોમાં,
$\lambd[x][(x+3)]$ એ તેની દલીલમાં ત્રણ ઉમેરતા વિધેયનું માત્ર
\emph{નામ} છે. આ અભિવ્યક્તિમાં~$x$ અસ્થાયી ચલ અથવા સ્થાનધારક છે:
એ જ વિધેયને $\lambd[y][(y + 3)]$ વડે પણ દર્શાવી શકાય છે. બીજા
પ્રાચલો હોય ત્યારે પણ આ સંકેતન ચાલે છે. દાખલા તરીકે,
$g(x, y)$ બે ચલોનું વિધેય હોય અને~$k$ પ્રાકૃતિક સંખ્યા હોય,
તો~$\lambd[x][g(x,k)]$ એ દરેક~$x$ને~$g(x, k)$માં મોકલતું વિધેય છે.

સંકેતી અભિવ્યક્તિ પરથી વિધેય વ્યાખ્યાયિત કરવાની આ રીત
\emph{લૅમ્બડા અમૂર્તીકરણ} કહેવાય છે. તેનું બીજું પાસું
\emph{અનુપ્રયોગ} છે: વિધેય~$f$ (માનો કે પ્રાકૃતિક સંખ્યાઓ પર
વ્યાખ્યાયિત) હોય તો તેને 2 જેવા મૂલ્ય પર \emph{લાગુ} કરી શકાય.
પરંપરાગત સંકેતનમાં તેના પરિણામ માટે આપણે~$f(2)$ લખીએ છીએ.

લૅમ્બડા અમૂર્તીકરણ અને અનુપ્રયોગ જોડીએ તો શું થાય? અમૂર્ત કરેલા
ચલની જગ્યાએ લાગુ કરાયેલી દલીલ મૂકીને મળેલી અભિવ્યક્તિને સરળ બનાવી
શકાય. દાખલા તરીકે,
\[
(\lambd[x][(x + 3)])(2)
\]
ને~$2 + 3$ તરીકે સરળ બનાવી શકાય.

અહીં સુધી આપણે પરંપરાગત ખ્યાલો માટે નવાં સંકેતનો જ આપ્યાં છે.
પરંતુ લૅમ્બડા કલન ગણસિદ્ધાંતના દૃષ્ટિકોણથી વધુ મૂળભૂત ફેરફાર
કરે છે. આ માળખામાં:
\begin{enumerate}
\item દરેક વસ્તુ વિધેય દર્શાવે છે.
\item લૅમ્બડા અમૂર્તીકરણથી વિધેયો વ્યાખ્યાયિત કરી શકાય છે.
\item કોઈપણ વસ્તુને બીજી કોઈપણ વસ્તુ પર લાગુ કરી શકાય છે.
\end{enumerate}
દાખલા તરીકે, લૅમ્બડા કલનમાં~$F$ પદ હોય તો~$F(F)$ હંમેશાં અર્થપૂર્ણ
માનવામાં આવે છે. આ ઉદાર માળખું \emph{પ્રકારવિહીન} લૅમ્બડા કલન
કહેવાય છે. અહીં ``પ્રકારવિહીન''નો અર્થ છે કે કઈ વસ્તુને કઈ વસ્તુ
પર લાગુ કરી શકાય તેના પર કોઈ નિયંત્રણ નથી.

\begin{digress}
લૅમ્બડા કલનનું \emph{પ્રકારિત} સ્વરૂપ પણ છે; તે પ્રકારવિહીન
સ્વરૂપનો મહત્વનો પ્રકારાંતર છે. ઘણી બાબતોમાં પ્રકારિત લૅમ્બડા કલન
પ્રકારવિહીન સ્વરૂપ જેવું હોવા છતાં, તેને પરંપરાગત ગણસિદ્ધાંતના
માળખા સાથે સુસંગત કરવું ઘણું સહેલું છે અને તેના કેટલાક ગુણધર્મો
તદ્દન જુદા છે.

લૅમ્બડા કલન વિશેનું સંશોધન સૈદ્ધાંતિક કમ્પ્યુટર વિજ્ઞાન અને
પ્રોગ્રામિંગ ભાષાઓની રચનામાં કેન્દ્રસ્થાને રહ્યું છે. 1950ના
દાયકામાં John McCarthyએ રચેલી LISP એ આ વિચારોની અસર ધરાવતી
ભાષાનું પ્રારંભિક ઉદાહરણ છે.
\end{digress}
% END OLP-0343
\endgroup


\begingroup
\makeatletter\def\ol@sectioncs{section}\makeatother

% BEGIN OLP-0344
\olfileid{lam}{int}{syn}
\olsection{લૅમ્બડા કલનની વાક્યરચના}
\phantomsection\label{guunit:OLP-0344}


સૌપ્રથમ $x$, $y$, $z$,~\dots{} ચલોની શ્રેણી અને
$a$, $b$, $c$,~\dots{} જેવા કેટલાક અચલ સંકેતો લઈએ. પદોનો ગણ
નીચે પ્રમાણે આગમનથી વ્યાખ્યાયિત થાય છે:
\begin{enumerate}
\item દરેક ચલ પદ છે.
\item દરેક અચલ પદ છે.
\item $M$ અને $N$ પદો હોય તો $(MN)$ પણ પદ છે.
\item $M$ પદ અને $x$ ચલ હોય તો $(\lambd[x][M])$ પણ પદ છે.
\end{enumerate}
$(MN)$ સ્વરૂપનાં પદોને \emph{અનુપ્રયોગો} અને
$(\lambd[x][M])$ સ્વરૂપનાં પદોને \emph{અમૂર્તીકરણો} કહે છે.

કોઈપણ અચલ વિનાની પદ્ધતિને \emph{શુદ્ધ} લૅમ્બડા કલન કહે છે.
મુખ્યત્વે આપણે શુદ્ધ~$\lambd$-કલનમાં કામ કરીશું; તેથી બધાં નાના
અક્ષરો ચલો દર્શાવશે. મોટા અક્ષરો ($M$, $N$ વગેરે) $\lambd$-કલનનાં
પદો દર્શાવશે.

આપણે સંકેતન અંગેના કેટલાક નિયમો અનુસરીશું:

\begin{conv}
\begin{enumerate}
\item કૌંસ છોડી દેવાય ત્યારે અનુપ્રયોગ ડાબેથી જમણે થાય છે.
  દાખલા તરીકે, $M$, $N$, $P$ અને $Q$ પદો હોય તો~$MNPQ$ એ
  $(((MN)P)Q)$નું સંક્ષિપ્ત સ્વરૂપ છે.
\item ફરી, કૌંસ છોડી દેવાય ત્યારે લૅમ્બડા અમૂર્તીકરણનો વ્યાપ
  શક્ય તેટલો મોટો લેવાય છે. દાખલા તરીકે, $\lambd[x][MNP]$ને
  $(\lambd[x][((MN)P)])$ એમ વાંચવું.
\item એક લૅમ્બડાથી અનેક ચલોનું અમૂર્તીકરણ કરી શકાય છે.
  દાખલા તરીકે, $\lambd[xyz][M]$ એ
  $\lambd[x][\lambd[y][\lambd[z][M]]]$નું સંક્ષિપ્ત સ્વરૂપ છે.
\end{enumerate}
\end{conv}

દાખલા તરીકે,
\[
\lambd[xy][xxyx \lambd[z][xz]]
\]
એ નીચેનું સંક્ષિપ્ત સ્વરૂપ છે:
\[
\lambd[x][\lambd[y][((((xx)y)x)(\lambd[z][(xz)]))]].
\]
આ નિયમો યાદ રાખવા જોઈએ. શરૂઆતમાં એ માથું ખાઈ જશે, પરંતુ
ધીમે ધીમે ટેવ પડી જશે; પછી અનેક કૌંસોની ગૂંચવણ સંભાળવા કરતાં
આ નિયમો ઓછા કંટાળાજનક લાગશે.

માત્ર બંધ ચલોની નામાવલિમાં જુદાં પડતાં બે પદોને
$\alpha$-સામ્યવાળાં કહે છે; દાખલા તરીકે, $\lambd[x][x]$ અને
$\lambd[y][y]$. તેમને ``એક જ'' પદ માનવું અનુકૂળ રહેશે.
એટલે આપણે~$M$ અને~$N$ એક જ છે એમ કહીએ ત્યારે બંધ ચલોનું નામ
ફેરવવા સુધીની સમાનતા પણ તેમાં સામેલ છે. કોઈ~$\lambd$ના વ્યાપમાં
આવતા ચલોને ``બંધ'' અને બાકીનાને ``મુક્ત'' કહે છે. અગાઉના ઉદાહરણમાં
મુક્ત ચલો નથી; પરંતુ
\[
(\lambd[z][yz])x
\]
માં~$y$ અને~$x$ મુક્ત છે, જ્યારે~$z$ બંધ છે.
% END OLP-0344
\endgroup


\begingroup
\makeatletter\def\ol@sectioncs{section}\makeatother

% BEGIN OLP-0345
\olfileid{lam}{int}{red}
\olsection{લૅમ્બડા પદોનું લઘુકરણ}
\phantomsection\label{guunit:OLP-0345}


લૅમ્બડા પદો સાથે શું કરી શકાય? તેમને સરળ બનાવી શકાય.
$M$ અને $N$ ગમે તે લૅમ્બડા પદો તથા~$x$ ગમે તે ચલ હોય, તો~$M$માં
$x$ની જગ્યાએ~$N$ મૂકવાથી મળતું પદ~$\Subst{M}{N}{x}$ વડે
દર્શાવી શકીએ. મૂક્યા પછી~$N$ના મુક્ત ચલો સાથે અથડામણ થાય એવા
$M$ના બંધ ચલોનું નામ પહેલેથી બદલી લેવું પડે. દાખલા તરીકે,
\[
\Subst{(\lambd[w][xxw])}{yyz}{x} = \lambd[w][(yyz)(yyz)w].
\]

\begin{digress}
સ્થાનાપન્ન માટેના બીજા સંકેતો $[N/x]M$, $[x/N]M$ અને
$M[x/N]$ પણ છે. ગૂંચવાડાથી સાવધાન!
\end{digress}

સહજ અર્થમાં $(\lambd[x][M])N$ અને~$\Subst{M}{N}{x}$નો અર્થ
સમાન છે. પહેલું પદ બીજા પદથી બદલવાની ક્રિયા
\emph{$\beta$-સંકોચન} કહેવાય છે. $(\lambd[x][M])N$ને
\emph{સંકોચ્ય પદ (રેડેક્સ)} અને~$\Subst{M}{N}{x}$ને તેનું
\emph{સંકોચન-ફળ} કહે છે. સામાન્ય રીતે, પદ~$P$ના કોઈ ઉપપદનું
$\beta$-સંકોચન કરીને~$P'$ મળે તો~$P$નું~$P'$માં
\emph{એક પગલામાં $\beta$-લઘુકરણ} થાય છે એમ કહીએ અને
$P \redone P'$ લખીએ. $P$માંથી એક-પગલાના અમુક લઘુકરણો દ્વારા
(શૂન્ય પણ હોઈ શકે) $P'$ મળે તો~$P$નું~$P'$માં
\emph{$\beta$-લઘુકરણ} થાય છે; સંકેતમાં~$P \red P'$.
જે પદનું આગળ~$\beta$-લઘુકરણ થઈ ન શકે તેને
\emph{$\beta$-અલઘુકરણીય} અથવા \emph{$\beta$-સામાન્ય}
કહે છે. સંદર્ભ સ્પષ્ટ હોય ત્યારે ``$\beta$-લઘુકરણ''ને બદલે
માત્ર ``લઘુકરણ'' વગેરે કહીશું.

કેટલાંક ઉદાહરણો જોઈએ.
\begin{enumerate}
\item આપણને મળે છે:
\begin{align*}
(\lambd[x][xxy]) \lambd[z][z] & \redone (\lambd[z][z])(\lambd[z][z]) y \\
& \redone (\lambd[z][z]) y \\
& \redone y.
\end{align*}
\item પદને ``સરળ'' કરતાં તે વધુ જટિલ પણ બની શકે છે:
\begin{align*}
(\lambd[x][xxy])(\lambd[x][xxy]) & \redone (\lambd[x][xxy])(\lambd[x][xxy])y \\
& \redone (\lambd[x][xxy])(\lambd[x][xxy])yy \\
& \redone \dots
\end{align*}
\item લઘુકરણ પછી પદ યથાવત્ પણ રહી શકે છે:
\[
(\lambd[x][xx])(\lambd[x][xx]) \redone (\lambd[x][xx])(\lambd[x][xx]).
\]
\item વળી, કેટલાક પદોનું એક કરતાં વધુ રીતે લઘુકરણ થઈ શકે છે.
  દાખલા તરીકે,
\[
(\lambd[x][(\lambd[y][yx]) z)] v \redone (\lambd[y][yv]) z
\]
બહારના અનુપ્રયોગનું સંકોચન કરીને મળે છે; અને
\[
(\lambd[x][(\lambd[y][yx]) z)] v \redone (\lambd[x][zx]) v
\]
અંદરના અનુપ્રયોગનું સંકોચન કરીને મળે છે. જોકે આ કિસ્સામાં
બંને પદો આગળ જઈને એક જ પદ~$zv$માં લઘુકૃત થાય છે.
\end{enumerate}

છેલ્લા ઉદાહરણમાં મળેલું સમાન અંતિમ પરિણામ સંયોગ નથી. તે
લૅમ્બડા કલનનો ``ચર્ચ--રોસર ગુણધર્મ'' નામનો ઊંડો અને મહત્વનો
ગુણધર્મ બતાવે છે.
% END OLP-0345
\endgroup


\begingroup
\makeatletter\def\ol@sectioncs{section}\makeatother

% BEGIN OLP-0346
\olfileid{lam}{int}{cr}
\olsection{ચર્ચ--રોસર ગુણધર્મ}
\phantomsection\label{guunit:OLP-0346}


\begin{thm}
\ollabel{thm:church-rosser}
$M$, $N_1$ અને~$N_2$ એવાં પદો હોય કે $M \red N_1$ અને $M \red
N_2$. તો એવું પદ~$P$ છે કે $N_1 \red P$ અને $N_2 \red P$.
\end{thm}

\begin{cor}
માનો કે~$M$નું સામાન્ય સ્વરૂપમાં લઘુકરણ થઈ શકે છે. તો આ
સામાન્ય સ્વરૂપ અનન્ય છે.
\end{cor}

\begin{proof}
જો $M \red N_1$ અને $M \red N_2$, તો અગાઉના પ્રમેય મુજબ એવું
પદ~$P$ છે કે~$N_1$ અને~$N_2$ બંનેનું~$P$માં લઘુકરણ થાય છે.
$N_1$ અને~$N_2$ બંને સામાન્ય સ્વરૂપમાં હોય, તો આ ત્યારે જ બની
શકે જ્યારે $N_1 \ident P \ident N_2$.
\end{proof}

છેલ્લે, બે પદો~$M$ અને~$N$ કોઈ એક જ પદમાં લઘુકૃત થાય તો તેમને
\emph{$\beta$-સામ્યવાળાં}, અથવા ટૂંકમાં \emph{સામ્યવાળાં},
કહીશું. એટલે કે એવું~$P$ હોય કે $M \red P$ અને $N \red P$.
તેને $M \equal[\beta] N$ લખીએ છીએ. 
\olref{thm:church-rosser}નો ઉપયોગ કરીને ચકાસી શકાય કે
$\equal[\beta]$ સામ્ય-સંબંધ છે અને તેમાં વધારાનો ગુણધર્મ છે:
દરેક $M$ અને~$N$ માટે $M \red N$ અથવા $N \red M$ હોય તો
$M \equal[\beta] N$. (ખરેખર, આ ગુણધર્મ ધરાવતા સામ્ય-સંબંધોમાં
$\equal[\beta]$ \emph{સૌથી નાનો} છે એમ બતાવી શકાય.)
% END OLP-0346
\endgroup


\begingroup
\makeatletter\def\ol@sectioncs{section}\makeatother

% BEGIN OLP-0347
\olfileid{lam}{int}{cur}
\olsection{કરિંગ}
\phantomsection\label{guunit:OLP-0347}

\sourcecorrection{OLLAM-001}{આ વિભાગ પ્રસ્તાવના પ્રકરણમાં
આયાત થાય છે, પરંતુ સ્થિર અંગ્રેજી મૂળના ટિપ્પણી-મથાળામાં
\texttt{representability} અને ફાઇલ-ઓળખમાં~\texttt{rep} છે.
પ્રકરણની વાસ્તવિક જગ્યા પ્રમાણે અહીં~\texttt{int} મૂક્યું છે.}

$\lambd$-અમૂર્ત અભિવ્યક્તિ~$\lambd[x][M]$ એક દલીલવાળું
વિધેય દર્શાવે છે. અનેક દલીલો સ્વીકારતું વિધેય વ્યાખ્યાયિત કરવું
હોય ત્યારે આ મર્યાદા લાગે છે. એક રસ્તો એ છે કે~$\lambd$-કલન
વિસ્તારીને જોડી, ત્રિક વગેરે બનાવવાની છૂટ આપીએ; પછી, દાખલા
તરીકે, ત્રિસ્થાની વિધેય~$\lambd[x][M]$ની દલીલ ત્રિક હશે.
પરંતુ આ માટે \emph{કરિંગ} વધુ અનુકૂળ છે.

દાખલો જોઈએ. થોડા સમય માટે માની લઈએ કે~$\lambd$-કલનમાં
$+$ ક્રિયા છે. સરવાળાનું વિધેય~$2$-સ્થાની છે, એટલે બે દલીલો
સ્વીકારે છે. પરંતુ~$\lambd$-અમૂર્તીકરણ માત્ર એક દલીલવાળાં વિધેયો
આપે છે: વાક્યરચનામાં $\lambd[(x,y)][(x+y)]$ જેવી અભિવ્યક્તિ
મંજૂર નથી. તેમ છતાં, $\lambd[y][(x+y)]$ વડે અપાયેલું એકસ્થાની
વિધેય~$f_x(y)$ લઈ શકીએ; તે તેની એકમાત્ર દલીલ~$y$માં~$x$ ઉમેરે
છે. ખરેખર, આ એક વિધેય નથી; દરેક સંખ્યા~$x$ માટે ``$x$ ઉમેરો''
એવાં જુદાં વિધેયોનો પરિવાર છે. હવે~$\lambd[x][f_x]$ તરીકે બીજું
એકસ્થાની વિધેય~$g$ વ્યાખ્યાયિત કરીએ. $x$ પર લાગુ કરતાં~$g(x)$
વિધેય~$f_x$ આપે છે; એટલે એનાં મૂલ્યો પોતે વિધેયો છે. પહેલાં~$g$
ને~$x$ પર અને પછી મળેલા પરિણામને~$y$ પર લાગુ કરીએ તો
$(g(x))y = f_x(y) = x+y$ મળે છે. આ રીતે એકસ્થાની વિધેય~$g$
બે-દલીલવાળા સરવાળાના વિધેય જેવું જ કામ કરે છે. બે-સ્થાની વિધેયને
એકસ્થાની વિધેયમાં (જેના મૂલ્યો એકસ્થાની વિધેયો છે) ફેરવવાની
આ યુક્તિને જ ``કરિંગ'' કહે છે.

હવે~$\lambd$-કલનની વાક્યરચનામાં યોગ્ય એવો દાખલો જોઈએ.
વિધેય $f(x,y) = x$ કેવી રીતે દર્શાવવું? બે દલીલો સ્વીકારીને
પહેલી પરત કરતું વિધેય જોઈએ તો $\lambd[x][\lambd[y][x]]$ લખી
શકીએ. તે ખરેખર દલીલ~$x$ સ્વીકારીને વિધેય~$\lambd[y][x]$
આપતું વિધેય છે. $\lambd[y][x]$ બીજી દલીલ~$y$ સ્વીકારે છે,
પરંતુ તેને અવગણીને હંમેશાં~$x$ જ પરત કરે છે. હવે
$\lambd[x][\lambd[y][x]]$ને બે દલીલો પર લાગુ કરીએ:
\begin{align*}
  (\lambd[x][\lambd[y][x]])MN
  \bredone & (\lambd[y][M])N \\
  \bredone & M
\end{align*}

સામાન્ય રીતે, અમુક પદ~$N$ વડે વ્યાખ્યાયિત $x_1$, \dots,~$x_n$
પ્રાચલોવાળું વિધેય લખવા માટે
$\lambd[x_1][\lambd[x_2][\ldots\lambd[x_n][N]]]$ લખી શકીએ.
તેને~$n$ દલીલો પર લાગુ કરતાં મળે છે:
\begin{multline*}
  (\lambd[x_1][\lambd[x_2][\ldots\lambd[x_n][N]]]) M_1 \dots M_n \bredone\\
  \begin{aligned}
  \bredone {} & (\Subst{(\lambd[x_2][\ldots\lambd[x_n][N]])}{M_1}{x_1}) M_2
  \dots M_n\\
   \eqs {} & (\lambd[x_2][\ldots\lambd[x_n][\Subst{N}{M_1}{x_1}]]) M_2
            \dots M_n \\
  \vdots & \\
  \bredone {} &\Subst{\Subst{N}{M_1}{x_1}\ldots}{M_n}{x_n}
  \end{aligned}
\end{multline*}
છેલ્લી પંક્તિનો સીધો અર્થ છે કે વિધેયની વ્યાખ્યાના શરીરમાં
દરેક~$x_i$ની જગ્યાએ~$M_i$ મૂકવું. અનેક દલીલો વિધેયને આપીએ ત્યારે
આપણને બરાબર એ જ જોઈએ છે.
\sourcecorrection{OLLAM-002}{સ્થિર અંગ્રેજી મૂળની છેલ્લી પંક્તિમાં
\texttt{\textbackslash Subst}ની અંદર~$P$ આવે છે, પરંતુ આ સામાન્ય
વિધેયનું શરીર~$N$ છે અને~$P$ અહીં વ્યાખ્યાયિત જ નથી. તેથી
અહીં~$N$ મૂક્યું છે.}
% END OLP-0347
\endgroup


\begingroup
\makeatletter\def\ol@sectioncs{section}\makeatother

% BEGIN OLP-0348
\olfileid{lam}{int}{rep}
\olsection{લૅમ્બડા દ્વારા વ્યાખ્યાયિત અંકગણિતીય વિધેયો}
\phantomsection\label{guunit:OLP-0348}


લૅમ્બડા કલન સંગણનના નિદર્શન તરીકે કેવી રીતે કામ કરી શકે? શરૂઆતમાં
આ વિધાનનો અર્થ શું એવો પ્રશ્ન પણ થાય. પ્રાકૃતિક સંખ્યાઓ પરની
સંગણનીયતા વિશે વાત કરવા માટે આ સંખ્યાઓનું યોગ્ય નિરૂપણ શોધવું
જોઈએ. નીચેનું નિરૂપણ આશ્ચર્યજનક રીતે સારું કામ કરે છે.

\begin{defn}
દરેક પ્રાકૃતિક સંખ્યા~$n$ માટે \emph{ચર્ચ અંક}
$\num{n}$ એ લૅમ્બડા પદ $\lambd[x][\lambd[y][(x(x(x(\dots
x(y)))))]]$ હોય એમ વ્યાખ્યાયિત કરો; તેમાં કુલ~$n$ વાર~$x$ આવે છે.
\end{defn}

$\num{n}$ પદો ``પુનરાવર્તકો'' છે: દલીલ~$f$ આપીએ તો~$\num{n}$
$y$ને~$f^n(y)$માં મોકલતું વિધેય આપે છે. દરેક ચર્ચ અંક સામાન્ય
સ્વરૂપમાં છે એ ધ્યાનમાં લો. હવે પ્રાકૃતિક સંખ્યાઓ પરના વિધેયને
લૅમ્બડા પદ ``ગણે છે'' તેનો અર્થ કહી શકીએ.

\begin{defn}
$f(x_0, \dots, x_{k-1})$ એ~$\Nat^k$માંથી~$\Nat$માં જતું
$k$-સ્થાની આંશિક વિધેય હોય. $\lambd$-પદ~$F$ એ~$f$ને
\emph{લૅમ્બડા દ્વારા વ્યાખ્યાયિત} કરે છે એમ ત્યારે અને માત્ર ત્યારે જ
કહીએ જ્યારે પ્રાકૃતિક સંખ્યાઓની દરેક શ્રેણી $n_0$, \dots,~$n_{k-1}$
માટે,
\[
F\, \num{n_0}\, \num{n_1} \dots \num{n_{k-1}} \red \num{f(n_0, n_1, \dots,
  n_{k-1})}
\]
જો $f(n_0, \dots, n_{k-1})$ વ્યાખ્યાયિત હોય, અને ન હોય તો
$F\, \num{n_0}\, \num{n_1} \dots \num{n_{k-1}}$ને કોઈ સામાન્ય
સ્વરૂપ ન હોય.
\end{defn}
\sourcecorrection{OLLAM-003}{સ્થિર અંગ્રેજી મૂળ અહીં
$f(x_0,\dots,x_{k-1})$ને~$n$-સ્થાની કહે છે, પરંતુ દલીલોની
સંખ્યા~$k$ છે. આ કિસ્સામાં વિધેયનું ઇનપુટ ક્ષેત્ર~$\Nat^k$
છે. અહીં બંને બાબતો સ્પષ્ટ કરી છે.}
\sourcecorrection{OLLAM-004}{સ્થિર અંગ્રેજી મૂળમાં વિધેય
અવ્યાખ્યાયિત હોય ત્યારે~$F$ પછી અલ્પવિરામ મૂકીને દલીલોની યાદી
લખી છે. અહીં અલ્પવિરામ દૂર કરીને અગાઉની પંક્તિની જેમ જ
$F$નો ચર્ચ અંકો પરનો અનુપ્રયોગ લખ્યો છે.}

\begin{thm}
\ollabel{thm:lambda-def}
વિધેય~$f$ આંશિક સંગણનીય હોય ત્યારે અને માત્ર ત્યારે જ તે
લૅમ્બડા પદ વડે લૅમ્બડા દ્વારા વ્યાખ્યાયિત હોય.
\end{thm}

\begin{explain}
આ પ્રમેય કંઈક અંશે ચોંકાવનારું છે. સંગણનના નિદર્શન તરીકે
લૅમ્બડા કલન ઘણું સરળ કલન છે: તેમાં માત્ર લૅમ્બડા અમૂર્તીકરણ અને
અનુપ્રયોગ જ ક્રિયાઓ છે!{} છતાં આટલાં અલ્પ સાધનોથી કોઈપણ
સંગણકીય પ્રક્રિયા અમલમાં મૂકી શકાય છે.
\end{explain}
% END OLP-0348
\endgroup


\begingroup
\makeatletter\def\ol@sectioncs{section}\makeatother

% BEGIN OLP-0349
\olfileid{lam}{int}{cmp}
\olsection{લૅમ્બડા દ્વારા વ્યાખ્યાયિત વિધેયો સંગણનીય છે}
\phantomsection\label{guunit:OLP-0349}


\begin{thm}
\ollabel{thm:lambda-computable}
આંશિક વિધેય~$f$ લૅમ્બડા પદ વડે લૅમ્બડા દ્વારા વ્યાખ્યાયિત હોય, તો તે
સંગણનીય છે.
\end{thm}

\begin{proof}
માનો કે વિધેય~$f$ લૅમ્બડા પદ~$X$ વડે લૅમ્બડા દ્વારા વ્યાખ્યાયિત છે.
$f$ ગણવાની અનૌપચારિક પ્રક્રિયા વર્ણવીએ. ઇનપુટ $m_0$,
\dots,~$m_{n-1}$ માટે પદ $X \num{m_0} \ldots \num{m_{n-1}}$
લખો. ત્યાર પછી વૃક્ષ બનાવો: પહેલાં મૂળ પદનાં બધાં એક-પગલાનાં
લઘુકરણો લખો; તેમની નીચે તે પદોનાં બધાં એક-પગલાનાં લઘુકરણો
(અર્થાત્ મૂળ પદનાં બે-પગલાનાં લઘુકરણો) લખો; અને એમ આગળ વધો.
ક્યારેય ચર્ચ અંક મળે તો તેને ઉત્તર તરીકે પરત કરો; ન મળે તો
વિધેયનું મૂલ્ય અવ્યાખ્યાયિત રહે છે.
\sourcecorrection{OLLAM-005}{સ્થિર અંગ્રેજી મૂળમાં પ્રથમ
ઇનપુટ-પદમાં~\texttt{\textbackslash num m\_0} અને
\texttt{\textbackslash num m\_\{n-1\}} છે. આથી~$m_0$
અને~$m_{n-1}$ આખાં \texttt{\textbackslash num}ના દલીલરૂપે
જતા નથી. નીચેની જ પંક્તિમાં વપરાતી સાચી શૈલી મુજબ અહીં
\texttt{\textbackslash num\{m\_0\}} વગેરે લખ્યું છે.}

ચર્ચની માન્યતાનો આધાર લઈને કહી શકાય કે આ વિધેય સંગણનીય છે.
પ્રમેયની વધુ સારી સાબિતી માટે આ શોધપ્રક્રિયાનું પુનરાવર્તી વર્ણન
આપવું જોઈએ. દાખલા તરીકે,
``$\fn{IsASubterm}$'', ``$\fn{Substitute}$'',
``$\fn{ReducesToInOneStep}$'', ``$\fn{ReductionSequence}$'',
``$\fn{Numeral}$'' વગેરે આદિમ પુનરાવર્તી વિધેયો અને સંબંધોની
શ્રેણી વ્યાખ્યાયિત કરી શકાય. ત્યાર પછી $f(m_0, \dots, m_{n-1})$
ગણવાની આંશિક પુનરાવર્તી પ્રક્રિયા એ છે કે
$X \num{m_0} \dots \num{m_{n-1}}$થી શરૂ થઈ ચર્ચ અંક પર પૂરી થતી
એક-પગલાનાં લઘુકરણોની શ્રેણી શોધો, અને તે અંકને અનુરૂપ સંખ્યા
પરત કરો. વિગત લાંબી અને કંટાળાજનક છે, પરંતુ બાકી નિયમિત છે.
\end{proof}
% END OLP-0349
\endgroup


\begingroup
\makeatletter\def\ol@sectioncs{section}\makeatother

% BEGIN OLP-0350
\olfileid{lam}{int}{lrp}
\olsection{સંગણનીય વિધેયો લૅમ્બડા દ્વારા વ્યાખ્યાયિત છે}
\phantomsection\label{guunit:OLP-0350}


\begin{thm}
\ollabel{thm:computable-lambda}
દરેક સંગણનીય આંશિક વિધેય લૅમ્બડા દ્વારા વ્યાખ્યાયિત છે.
\end{thm}

\begin{proof}
દરેક આંશિક સંગણનીય વિધેય~$f$ કોઈ લૅમ્બડા પદ~$F$ વડે
લૅમ્બડા દ્વારા વ્યાખ્યાયિત છે એમ બતાવવું છે. ક્લીનીના સામાન્ય સ્વરૂપના
પ્રમેય મુજબ, દરેક આદિમ પુનરાવર્તી વિધેય કોઈ લૅમ્બડા પદ વડે
લૅમ્બડા દ્વારા વ્યાખ્યાયિત છે અને લૅમ્બડા દ્વારા વ્યાખ્યાયિત વિધેયો યોગ્ય
સંયોજનો તથા અમર્યાદિત શોધ હેઠળ બંધ છે એમ બતાવવું પૂરતું છે.
દરેક આદિમ પુનરાવર્તી વિધેય લૅમ્બડા પદ વડે લૅમ્બડા દ્વારા વ્યાખ્યાયિત
છે એમ બતાવવા માટે આરંભિક વિધેયો લૅમ્બડા દ્વારા વ્યાખ્યાયિત છે અને
લૅમ્બડા દ્વારા વ્યાખ્યાયિત આંશિક વિધેયો સંયોજન, આદિમ પુનરાવર્તન
અને અમર્યાદિત શોધ હેઠળ બંધ છે એમ બતાવવું પૂરતું છે.
\end{proof}

સાબિતીનો બાકીનો ભાગ વધુ વાંચનીય બને તે માટે આપણે વધુ પ્રચલિત
સંકેતન વાપરીશું. દાખલા તરીકે, $Mxyz$ને બદલે $M(x, y, z)$ લખીશું.
આ સંકેતન કંઈક સૂચવે છે, પરંતુ યાદ રાખવું કે પ્રકારવિહીન લૅમ્બડા
કલનના પદો સાથે નિશ્ચિત દલીલ-સંખ્યા જોડાયેલી નથી. તેથી એ જ
પદ~$M$ માટે $M(x,y)$ અને $M(x,y,z,w)$ બંને લખવું અર્થપૂર્ણ છે.
તેમ છતાં, આ સંકેતનથી આપણે~$M$ને ત્રણ ચલોનું વિધેય માનીને
વાપરીએ છીએ તે સ્પષ્ટ થાય છે અને વ્યાખ્યાઓ પાછળનો હેતુ સમજાય છે.
એ જ રીતે, ``$M = \lambd[x][\lambd[y][\lambd[z][\dots]]]$
થી~$M$ વ્યાખ્યાયિત કરો'' એમ કહેવાને બદલે ``$M(x,y,z) = \dots$
થી~$M$ વ્યાખ્યાયિત કરો'' એમ કહીશું.
% END OLP-0350
\endgroup


\begingroup
\makeatletter\def\ol@sectioncs{section}\makeatother

% BEGIN OLP-0351
\olfileid{lam}{int}{bas}
\olsection{મૂળભૂત આદિમ પુનરાવર્તી વિધેયો લૅમ્બડા દ્વારા વ્યાખ્યાયિત છે}
\phantomsection\label{guunit:OLP-0351}


\begin{lem}
વિધેયો $\Zero$, $\Succ$ અને $\Proj{n}{i}$ લૅમ્બડા દ્વારા વ્યાખ્યાયિત છે.
\end{lem}

\begin{proof}
$\Zero$ એટલે માત્ર $\lambd[x][\lambd[y][y]]$.

ઉત્તરગામી વિધેય~$\Succ$ને
$\fn{Succ}(u) = \lambd[x][\lambd[y][x(uxy)]]$ વડે વ્યાખ્યાયિત
કરીએ છીએ. આ કેવી રીતે કામ કરે છે તે વિચારો. દરેક ચર્ચ અંક~$\num{n}$ને
પુનરાવર્તક માનો અને કોઈપણ વિધેય~$f$ લો. ત્યારે
$\fn{Succ}(\num{n},f)$ એવું વિધેય છે જે ઇનપુટ~$y$ પર, $y$થી
શરૂ કરીને~$f$ને $n$ વખત લાગુ કરે છે અને પછી વધુ એક વખત લાગુ કરે છે.

પ્રક્ષેપો માટે કંઈ વિશેષ કહેવાનું નથી:
$\fn{Proj}^n_i(x_0, \dots, x_{n-1}) = x_i$. બીજા શબ્દોમાં,
આપણા સંકેતન મુજબ $\fn{Proj}^n_i$ એ લૅમ્બડા પદ
$\lambd[x_0][\dots \lambd[x_{n-1}][x_i]]$ છે.
\end{proof}
% END OLP-0351
\endgroup


\begingroup
\makeatletter\def\ol@sectioncs{section}\makeatother

% BEGIN OLP-0352
\olfileid{lam}{int}{com}
\olsection{લૅમ્બડા દ્વારા વ્યાખ્યાયિત વિધેયો સંયોજન હેઠળ બંધ છે}
\phantomsection\label{guunit:OLP-0352}


\begin{lem}
લૅમ્બડા દ્વારા વ્યાખ્યાયિત વિધેયો સંયોજન હેઠળ બંધ છે.
\end{lem}

\begin{proof}
ધારો કે $h$, $g_0,$ \dots,~$g_{k-1}$ના સંયોજનથી~$f$
વ્યાખ્યાયિત થાય છે. હવે ધારો કે $H$, $G_0$, \dots,~$G_{k-1}$ વડે
અનુક્રમે $h$, $g_0$, \dots,~$g_{k-1}$ લૅમ્બડા દ્વારા વ્યાખ્યાયિત છે. આપણે
$f$ને લૅમ્બડા દ્વારા વ્યાખ્યાયિત કરતું પદ~$F$ શોધવાનું છે. પરંતુ~$F$ને
સીધું આ રીતે વ્યાખ્યાયિત કરી શકીએ:
\[
F(x_0, \dots, x_{l-1}) = H(G_0(x_0, \dots, x_{l-1}),
\dots, G_{k-1}(x_0, \dots, x_{l-1})).
\]
બીજા શબ્દોમાં, લૅમ્બડા કલનની ભાષા સંયોજનનું નિરૂપણ કરવા માટે
અનુકૂળ છે.
\end{proof}
% END OLP-0352
\endgroup


\begingroup
\makeatletter\def\ol@sectioncs{section}\makeatother

% BEGIN OLP-0353
\olfileid{lam}{int}{pr}
\olsection{લૅમ્બડા દ્વારા વ્યાખ્યાયિત વિધેયો આદિમ પુનરાવર્તન હેઠળ બંધ છે}
\phantomsection\label{guunit:OLP-0353}


આદિમ પુનરાવર્તનની વાત આવે ત્યારે આખરે થોડું કામ કરવું પડે છે. આપણે
તબક્કાવાર આગળ વધશું. અગાઉની જેમ, ધારો કે આપણી પાસે એવા પદો~$G'$ અને
$H'$ છે જે અનુક્રમે વિધેયો~$g$ અને~$h$ને લૅમ્બડા દ્વારા વ્યાખ્યાયિત કરે છે.
આપણે નીચે પ્રમાણે વ્યાખ્યાયિત વિધેય~$f$ને લૅમ્બડા દ્વારા વ્યાખ્યાયિત કરતું
પદ~$F$ જોઈએ છે:
\begin{align*}
f(0, \vec z) & = g(\vec z) \\
f(x+1, \vec z) & = h(x, f(x,\vec z), \vec z).
\end{align*}
તેથી, સામાન્ય રીતે, આપેલાં લૅમ્બડા પદો~$G'$ અને~$H'$ માટે દરેક
પ્રાકૃતિક સંખ્યા~$n$ માટે નીચેની શરતો પૂરું કરતું પદ~$F$ શોધવું પૂરતું છે:
\begin{align*}
F(\num{0}, \vec z) & \equiv G'(\vec z) \\
F(\overline{n+1}, \vec z) & \equiv H'(\num{n}, F(\num{n}, \vec z), \vec z)
\end{align*}
કારણ કે~$G'$ અને~$H'$ અનુક્રમે~$g$ અને~$h$ને લૅમ્બડા દ્વારા વ્યાખ્યાયિત કરે
છે, $\vec z$ની જગ્યાએ અંકો~$\num{\vec m}$ મૂકતાં
$F(\num{n+1}, \num{\vec m})$ યોગ્ય ઉત્તરનાં સામાન્ય સ્વરૂપમાં
પહોંચશે.
\sourcecorrection{OLLAM-006}{સ્થિર અંગ્રેજી મૂળ શરૂઆતમાં~$G$
અને~$H$ આપે છે છતાં આગળ~$G'$ અને~$H'$ વાપરે છે, અને શોધવાનું પદ
ફરી~$H$ કહે છે. અહીં આપેલાં પદો~$G'$, $H'$ અને શોધવાનું પદ~$F$
રાખ્યાં છે; પહેલી બે પુનરાવર્તન-શરતોમાં પણ એ જ આપેલાં પદો છે.}
\sourcecorrection{OLLAM-007}{સ્થિર મૂળના પગથિયાના સમીકરણમાં
$h(z,f(x,\vec z),\vec z)$ લખાયું છે. અહીં પુનરાવર્તનનો સૂચક~$x$
છે, તેથી~$h(x,f(x,\vec z),\vec z)$ મૂક્યું છે.}

આ માટે, પરંતુ, નીચેની શરતો પૂરું કરતું પદ~$F$ શોધવું પૂરતું છે:
\begin{align*}
F(\num 0) & \equiv G \\
F(\overline {n+1}) & \equiv H(\num{n},F(\num{n}))
\intertext{દરેક પ્રાકૃતિક સંખ્યા $n$ માટે, જ્યાં}
G & = \lambd[\vec z][G'(\vec z)] \text{ અને}\\
H(u,v) & = \lambd[\vec z][H'(u,v(\vec z),\vec z)].
\end{align*}
બીજા શબ્દોમાં, લૅમ્બડા સંકેતનની યુક્તિથી વધારાના પ્રાચલો~$\vec z$
વિશે અલગથી ચિંતા કરવાની જરૂર રહેતી નથી---તે લૅમ્બડા પદમાં જ
સમાઈ જાય છે.
\sourcecorrection{OLLAM-008}{સ્થિર અંગ્રેજી મૂળમાં
$H(u,v)$ની વ્યાખ્યામાં~$v(u,\vec z)$ છે. અહીં~$v=F(\num{n})$
પહેલેથી~$n$ પર લાગુ કરેલું પદ છે અને હવે તે માત્ર~$\vec z$
સ્વીકારે છે; તેથી~$v(\vec z)$ મૂક્યું છે.}

પદ~$F$ વ્યાખ્યાયિત કરતાં પહેલાં ક્રમિત જોડીઓ સંભાળવાની એક રીત
જોઈએ. આગલું લેમ્મા તે આપે છે.

\begin{lem}
એવું લૅમ્બડા પદ~$D$ છે કે લૅમ્બડા પદોની દરેક જોડી~$M$ અને~$N$
માટે $D(M,N)(\num{0}) \red M$ અને $D(M,N)(\num{1}) \red N$.
\end{lem}

\begin{proof}
પહેલાં લૅમ્બડા પદ~$K$ને આ રીતે વ્યાખ્યાયિત કરો:
\[
K(y) = \lambd[x][y].
\]
બીજા શબ્દોમાં, $K$ એ પદ~$\lambd[y][\lambd[x][y]]$ છે. બીજી
દૃષ્ટિએ, દરેક~$M$ માટે $K(M)$ એવું અચલ વિધેય છે જે કોઈપણ ઇનપુટ
પર~$M$ પરત કરે છે.

હવે $D(x,y,z)$ને $D(x,y,z) = z (K(y))x$ વડે વ્યાખ્યાયિત કરો.
ત્યારે મળે છે:
\begin{align*}
D(M,N,\num 0) & \red \num 0 (K(N)) M \red M \text{ અને}\\
D(M,N,\num 1) & \red \num 1 (K(N)) M \red K(N) M \red N,
\end{align*}
જે જોઈએ હતું તે જ.
\end{proof}

વિચાર એ છે કે $D(M,N)$ જોડી~$\tuple{M, N}$નું નિરૂપણ કરે છે.
જો~$P$ આવી જોડીનું નિરૂપણ કરે એમ માનીએ, તો $P(\num 0)$ અને
$P(\num 1)$ અનુક્રમે ડાબા અને જમણા પ્રક્ષેપો, $(P)_0$ અને
$(P)_1$, દર્શાવે છે. આગળ આપણે આ છેલ્લાં સંકેતો વાપરીશું.

\begin{lem}
લૅમ્બડા દ્વારા વ્યાખ્યાયિત વિધેયો આદિમ પુનરાવર્તન હેઠળ બંધ છે.
\end{lem}

\begin{proof}
આપેલાં કોઈપણ પદો~$G$ અને~$H$ માટે નીચેની શરતો પૂરું કરતું
પદ~$F$ શોધી શકાય એમ બતાવવાનું છે:
\begin{align*}
F(\num{0}) & \equiv G \\
F(\num{n+1}) & \equiv H(\num{n}, F(\num{n}))
\end{align*}
દરેક પ્રાકૃતિક સંખ્યા~$n$ માટે. મૂળ વિચાર અંકોને પુનરાવર્તક તરીકે
વાપરી \emph{જોડીઓ}ની નીચેની શ્રેણી ગણવાનો છે:
\[
\tuple{\num 0, F(\num 0)}, \tuple{\num 1, F(\num 1)}, \dots,
\]
પ્રથમ જોડી માત્ર~$\tuple{\num 0, G}$ છે તે ધ્યાનમાં લો. જોડી
$\tuple{\num{n}, F(\num{n})}$ આપેલી હોય, તો આગલી જોડી
$\tuple{\num{n+1}, F(\num{n+1})}$ એ
$\tuple{\num{n+1}, H(\num{n}, F(\num{n}))}$ને સમકક્ષ હોવી
જોઈએ. આ એક પગથિયાનો ફેરફાર કરતું લૅમ્બડા પદ~$T$ રચીશું.

વિગતો આ પ્રમાણે છે. $T(u)$ને આ રીતે વ્યાખ્યાયિત કરો:
\[
T(u) = \tuple{S((u)_0), H((u)_0,(u)_1)}.
\]
અહીં પ્રથમ ઘટકમાં અગાઉ વ્યાખ્યાયિત ઉત્તરગામી લૅમ્બડા પદ વપરાય છે.
હવે ચકાસવું સરળ છે કે કોઈપણ સંખ્યા~$n$ માટે
\[
T(\tuple{\num{n}, M}) \red \tuple{\num{n+1}, H(\num{n}, M)}.
\]
ઉપર સૂચવ્યા મુજબ, $G$ અને~$H$ આપેલાં હોય ત્યારે~$F(u)$ને
આ રીતે વ્યાખ્યાયિત કરો:
\[
F(u) = (u(T,\tuple{\num 0, G}))_1.
\]
બીજા શબ્દોમાં, ઇનપુટ~$\num{n}$ પર~$F$, જોડી
$\tuple{\num 0, G}$ પરથી~$T$ને $n$ વખત લાગુ કરે છે અને પછી
બીજો ઘટક પરત કરે છે. શરૂઆતમાં મળે છે:
\begin{enumerate}
\item $\num{0} (T, \tuple{\num 0, G}) \equiv \tuple{\num{0}, G}$
\item $F(\num{0}) \equiv G$
\end{enumerate}
$n$ પર આગમનથી દરેક પ્રાકૃતિક સંખ્યા માટે નીચેની બે બાબતો
બતાવી શકીએ:
\begin{enumerate}
\item $\num{n+1}(T, \tuple{\num{0}, G}) \equiv \tuple{\num{n+1}, F(\num{n+1})}$
\item $F(\num{n+1}) \equiv H(\num{n}, F(\num{n}))$
\end{enumerate}
બીજી બાબત માટે,
\begin{align*}
F(\num{n+1}) & \red (\num{n+1}(T, \tuple{\num 0, G}))_1 \\
& \equiv (T(\num{n} (T, \tuple{\num 0, G})))_1 \\
& \equiv (T(\tuple{\num{n}, F(\num{n})}))_1 \\
& \equiv (\tuple{\num{n+1}, H(\num{n}, F(\num{n}))})_1 \\
& \equiv H(\num{n}, F(\num{n})).
\end{align*}
અહીં છેલ્લેથી બીજી પંક્તિમાં આગમનની ધારણા વાપરી છે. પહેલી
બાબત માટે,
\begin{align*}
\num{n+1} (T, \tuple{\num 0, G}) &
\equiv T(\num{n} (T, \tuple{\num 0, G})) \\
& \equiv T( \tuple{\num{n}, F(\num{n})}) \\
& \equiv \tuple{\num{n+1}, H(\num{n}, F(\num{n}))} \\
& \equiv \tuple{\num{n+1}, F(\num{n+1})}.
\end{align*}
અહીં છેલ્લી પંક્તિમાં બીજી બાબત વાપરી છે. આમ આપણે બતાવ્યું કે
$F(\num 0) \equiv G$ અને દરેક~$n$ માટે
$F(\num {n+1}) \equiv H(\num n, F(\num{n}))$; બરાબર એ જ
આપણને જોઈએ હતું.
\end{proof}
% END OLP-0353
\endgroup


\begingroup
\makeatletter\def\ol@sectioncs{section}\makeatother

% BEGIN OLP-0354
\olfileid{lam}{int}{fix}
\olsection{નિશ્ચિત-બિંદુ સંયોજકો}
\phantomsection\label{guunit:OLP-0354}


ધારો કે આપણી પાસે લૅમ્બડા પદ~$g$ છે અને એવું બીજું પદ~$k$ જોઈએ છે
કે~$k$, $gk$ સાથે $\beta$-સમકક્ષ હોય. આપણા સંકેતન પ્રમાણે પદો
વ્યાખ્યાયિત કરો:
\[
\fn{diag}(x) = xx
\]
અને
\[
l(x) = g(\fn{diag}(x))
\]
બીજા શબ્દોમાં, $l$ એ પદ~$\lambd[x][g(xx)]$ છે. $k$ એ પદ~$ll$
રાખો. ત્યારે મળે છે:
\begin{align*}
k & = (\lambd[x][g(xx)])(\lambd[x][g(xx)]) \\
& \red  g((\lambd[x][g(xx)])(\lambd[x][g(xx)])) \\
& = gk.
\end{align*}
જો
\[
Y = \lambd[g][((\lambd[x][g(xx)])(\lambd[x][g(xx)]))]
\]
રાખીએ, તો $Yg$ અને $g(Yg)$ બંને એક જ પદ સુધી લઘુકરણ પામે છે;
આથી $Yg \equiv_\beta g(Yg)$. તેને ``કરીનો સંયોજક'' કહે છે. તેના
બદલે જો
\[
Y = (\lambd[xg][g(xxg)])(\lambd[xg][g(xxg)])
\]
રાખીએ, તો ખરેખર $Yg$નું $g(Yg)$ સુધી લઘુકરણ થાય છે; આ વધુ
મજબૂત વિધાન છે. $Y$ના આ બીજા રૂપને ``ટ્યુરિંગનો સંયોજક'' કહે છે.
% END OLP-0354
\endgroup


\begingroup
\makeatletter\def\ol@sectioncs{section}\makeatother

% BEGIN OLP-0355
\olfileid{lam}{int}{min}
\olsection{લૅમ્બડા દ્વારા વ્યાખ્યાયિત વિધેયો લઘુત્તમીકરણ હેઠળ બંધ છે}
\phantomsection\label{guunit:OLP-0355}


\begin{lem}
ધારો કે $f(x,y)$ લૅમ્બડા દ્વારા વ્યાખ્યાયિત છે. વિધેય~$g$ને
આ રીતે વ્યાખ્યાયિત કરો:
\[
g(x) \simeq \umin{y}{f(x,y)}.
\]
ત્યારે~$g$ પણ લૅમ્બડા દ્વારા વ્યાખ્યાયિત છે.
\end{lem}

\begin{proof}
મુખ્ય વિચાર લગભગ આ છે. $x$ આપેલું હોય ત્યારે નિશ્ચિત-બિંદુ
લૅમ્બડા પદ~$Y$ વડે એવું વિધેય~$h_x(n)$ વ્યાખ્યાયિત કરીશું જે~$n$થી
શરૂ કરીને યોગ્ય~$y$ શોધે; પછી $g(x)$ એટલે માત્ર~$h_x(0)$.
વિધેય~$h_x$ને નિશ્ચિત-બિંદુ સમીકરણના ઉકેલ તરીકે લખી શકાય:
\[
h_x(n) \simeq
\begin{cases}
n & \text{જો $f(x,n) = 0$ હોય} \\
h_x(n+1) & \text{અન્યથા.}
\end{cases}
\]

હવે વિગતો જોઈએ. $f$ લૅમ્બડા-વ્યાખ્યેય હોવાથી કોઈ પદ~$F$
વડે તે લૅમ્બડા દ્વારા વ્યાખ્યાયિત છે. યાદ કરો કે આપણી પાસે એવું લૅમ્બડા
પદ~$D$ પણ છે કે $D(M, N, \bar{0}) \red M$ અને $D(M, N, \bar{1})
\red N$. હાલ પૂરતું~$x$ નિશ્ચિત રાખીએ. $h_x$ને
લૅમ્બડા દ્વારા વ્યાખ્યાયિત કરવા માટે~$x$ પર આધાર રાખતું એવું પદ~$H$
શોધવું છે જે નીચેની શરત પૂરી કરે:
\[
H(\num{n}) \equiv D(\num{n}, H(S(\num{n})), F(x, \num{n})).
\]
નિશ્ચિત-બિંદુ પદ~$Y$ વડે આ કરી શકાય છે. પહેલાં~$U$ને નીચેનું
પદ માનો:
\[
\lambd[h][\lambd[z][D(z,(h(Sz)),F(x,z))]],
\]
અને પછી~$H$ને પદ~$YU$ માનો. ધ્યાન આપો કે~$H$માં એકમાત્ર મુક્ત
ચલ~$x$ છે. હવે~$H$ ઉપરનું સમીકરણ પૂરે છે તે બતાવીએ.
\sourcecorrection{OLLAM-009}{લેમ્માની ધારણા માત્ર એટલી છે કે~$f$
લૅમ્બડા-વ્યાખ્યેય છે. સ્થિર અંગ્રેજી મૂળ અહીં અચાનક~$f$ને આદિમ
પુનરાવર્તી કહે છે, જે અનાવશ્યક મજબૂત ધારણા છે. આપેલી ધારણાથી
જ~$F$ મળે છે એમ ગુજરાતી પાઠમાં લખ્યું છે.}

$Y$ની વ્યાખ્યા મુજબ,
\[
H = YU \equiv U(YU) = U(H).
\]
ખાસ કરીને, દરેક પ્રાકૃતિક સંખ્યા~$n$ માટે,
\begin{align*}
H(\num{n}) & \equiv U(H, \num{n}) \\
& \red D(\num{n}, H(S(\num{n})), F(x, \num{n})),
\end{align*}
જે જોઈએ હતું તે જ. છેલ્લી પંક્તિમાં~$x$ની જગ્યાએ અંક~$\num{m}$
મૂકો. જો $F(\num{m}, \num{n})$નું લઘુકરણ~$\num{0}$ સુધી થાય,
તો આખી અભિવ્યક્તિ~$\num{n}$ સુધી ઘટે છે; અને જો
$F(\num{m}, \num{n})$નું લઘુકરણ કોઈ બીજા અંક સુધી થાય,
તો અભિવ્યક્તિ~$H(S(\num{n}))$ સુધી ઘટે છે.
આ બીજો કિસ્સો અંક એક પૂરતો મર્યાદિત નથી: જોડી-પસંદગીમાં શૂન્યથી
મોટા ચર્ચ અંક પર અચલ વિધેય વારંવાર લાગુ પડે છે અને બીજો ઘટક જ આપે છે.

સાબિતી પૂરી કરવા~$G$ને $\lambd[x][H(\num 0)]$ માનો.
ત્યારે~$G$, $g$ને લૅમ્બડા દ્વારા વ્યાખ્યાયિત કરે છે. બીજા શબ્દોમાં,
દરેક~$m$ માટે, જો~$g(m)$ વ્યાખ્યાયિત હોય તો
$G(\num m)$નું લઘુકરણ~$\overline {g(m)}$ સુધી થાય છે;
નહીં તો તેનું કોઈ સામાન્ય સ્વરૂપ હોતું નથી.
\end{proof}
% END OLP-0355
\endgroup


\OLEndChapterHook
% END OLP-0342
\endgroup


\begingroup
\makeatletter\def\ol@sectioncs{section}\makeatother

% BEGIN OLP-0356
\olchapter{lam}{syn}{વાક્યરચના}
\olchapterid{lam}{syn}
\phantomsection\label{guunit:OLP-0356}


\begingroup
\makeatletter\def\ol@sectioncs{section}\makeatother

% BEGIN OLP-0357
\olfileid{lam}{syn}{trm}
\olsection{પદો}
\phantomsection\label{guunit:OLP-0357}


લૅમ્બડા કલનનાં પદો ચલોનો અનંત પુરવઠો $\Obj{v_0}$, $\Obj{v_1}$,
\dots{}, સંકેત~``$\lambd$'' અને કૌંસોમાંથી આગમનથી બને છે. ચલો
દર્શાવવા આપણે $x$, $y$, $z$, \dots{} વાપરીશું અને પદો દર્શાવવા
$M$, $N$, $P$, \dots{} વાપરીશું.

\begin{defn}[પદો] \ollabel{defn:term}
લૅમ્બડા કલનનાં \emph{પદો}નો ગણ આગમનથી આ રીતે વ્યાખ્યાયિત થાય છે:
\begin{enumerate}
  \item \ollabel{defn:term-var} $x$ ચલ હોય તો $x$ પદ છે.
  \item \ollabel{defn:term-abs} $x$ ચલ અને $M$ પદ હોય તો
    $(\lambd[x][M])$ પદ છે.
  \item \ollabel{defn:term-app} $M$ અને $N$ બંને પદો હોય તો
    $(MN)$ પદ છે.
\end{enumerate}
\end{defn}

કોઈ પદ~$(\lambd[x][M])$ \olref{defn:term-abs} મુજબ રચાયું હોય
તો તેને \emph{અમૂર્તીકરણ}નું પરિણામ કહીએ છીએ અને તેમાં
$\lambd[x]$માંના~$x$ને \emph{પ્રાચલ} કહીએ છીએ. એ જ રીતે,
\olref{defn:term-app} મુજબ રચાયેલું પદ~$(MN)$
\emph{અનુપ્રયોગ}નું પરિણામ છે.

ઉપર વ્યાખ્યાયિત પદોમાં બધા જરૂરી કૌંસ છે. આ ક્યારેક ભારે પડે છે;
દાખલા તરીકે પદ
$(\lambd[x][((\lambd[x][x])(\lambd[x][(xx)]))])$ જુઓ. કૌંસ
છોડી શકાય તે માટે આપણે આગળ નિયમો આપીશું. જોકે, આ ઔપચારિક
વ્યાખ્યાથી \olref{defn:term} મુજબ પદ કઈ રીતે રચાયું છે તે નક્કી
કરવું સરળ બને છે. દાખલા તરીકે,
$(\lambd[x][((\lambd[x][x])(\lambd[x][(xx)]))])$ પદની રચનાનું
છેલ્લું પગલું અમૂર્તીકરણ જ હોવું જોઈએ, જેમાં પ્રાચલ~$x$ છે.
તે પદ~$((\lambd[x][x])(\lambd[x][(xx)]))$માંથી અમૂર્તીકરણ વડે
બને છે; આ બીજું પદ બે પદોનો અનુપ્રયોગ છે. તે બંને પદો પોતપોતે
અમૂર્તીકરણથી બને છે, અને આ રીતે આગળ જઈ શકાય છે.

\begin{prob}
$(\lambd[g][((\lambd[x][(g (x x))])
  (\lambd[x][(g (x x))]))])$ની રચના વર્ણવો.
\sourcecorrection{OLLAM-010}{સ્થિર અંગ્રેજી મૂળના આ અભ્યાસપ્રશ્નમાં
બે અમૂર્તીકરણોનો અનુપ્રયોગ બહારના કૌંસ વિના લખાયો છે. અહીં
વ્યાખ્યાના~$(MN)$ નિયમ અનુસાર તેમની આસપાસ કૌંસ ઉમેર્યા છે.}
\end{prob}
% END OLP-0357
\endgroup


\begingroup
\makeatletter\def\ol@sectioncs{section}\makeatother

% BEGIN OLP-0358
\olfileid{lam}{syn}{unq}
\olsection{અનન્ય વાંચનીયતા}
\phantomsection\label{guunit:OLP-0358}


દરેક પદને રચવાની રીત અનન્ય છે કે કેમ એવો પ્રશ્ન થાય; તેનો
જવાબ હા છે. દરેક લૅમ્બડા પદને રચવાની અને તેનો અર્થ સમજવાની
રીત એક જ છે. આગળની ચર્ચામાં, \emph{રચનાક્રમ} એટલે
\olref[trm]{defn:term}ના રચના-નિયમો એક કે અનેક વાર વાપરીને
પદ રચવાની પ્રક્રિયા.

\begin{lem}\ollabel{lem:term-start}
પદની શરૂઆત ચલથી અથવા કૌંસથી થાય છે.
\end{lem}

\begin{proof}
\olref[trm]{defn:term} મુજબ રચાયું હોય તો જ કંઈક પદ ગણાય.
તે \olref[trm]{defn:term-var}નું પરિણામ હોય તો ચલ જ હોવું જોઈએ.
તે \olref[trm]{defn:term-abs} અથવા
\olref[trm]{defn:term-app}નું પરિણામ હોય તો તેની શરૂઆત કૌંસથી
થાય છે.
\end{proof}

\begin{lem}\ollabel{lem:app-start}
અનુપ્રયોગના પરિણામની શરૂઆત બે ખુલતા કૌંસથી અથવા ખુલતા કૌંસ
પછી ચલથી થાય છે.
\end{lem}

\begin{proof}
$M$ અનુપ્રયોગનું પરિણામ હોય તો તે~$(PQ)$ સ્વરૂપનું છે;
તેથી તેની શરૂઆત કૌંસથી થાય છે. $P$ પદ હોવાથી,
\olref{lem:term-start} મુજબ તેની શરૂઆત કૌંસથી અથવા ચલથી
થાય છે.
\end{proof}

\begin{lem}\ollabel{lem:initial}
પદનો કોઈ યોગ્ય આરંભિક ખંડ પોતે પદ નથી.
\end{lem}

\begin{prob}
પદોની લંબાઈ પર આગમનથી
\olref[lam][syn][unq]{lem:initial} સાબિત કરો.
\end{prob}

\begin{prop}[અનન્ય વાંચનીયતા] \ollabel{prop:unq}
દરેક પદનો રચનાક્રમ અનન્ય છે. બીજા શબ્દોમાં, કોઈ રચનાક્રમથી
પદ~$M$ રચાયું હોય તો એ પદ રચી શકે એવો બીજો રચનાક્રમ નથી.
\end{prop}

\begin{proof}
  પદોના રચનાક્રમ પર આગમનથી આ સાબિત કરીએ.
  \begin{enumerate}
    \item $M$ એ~$x$ સ્વરૂપનું છે, જ્યાં~$x$ કોઈ ચલ છે.
      અમૂર્તીકરણ અને અનુપ્રયોગના પરિણામો હંમેશા કૌંસથી શરૂ
      થાય છે; તેથી~$M$ રચવામાં તે વપરાયા હોઈ શકે નહીં.
      આમ~$M$નો રચનાક્રમ
      \olref[trm]{defn:term}\olref[trm]{defn:term-var}નું
      એક જ પગલું હોવો જોઈએ.
    \item $M$ એ~$(\lambd[x][N])$ સ્વરૂપનું છે, જ્યાં~$x$
      કોઈ ચલ અને~$N$ પદ છે. તે એકલો ચલ નથી; તેથી
      \olref[trm]{defn:term}\olref[trm]{defn:term-var}
      મુજબ રચાયું હોઈ શકે નહીં. \olref{lem:app-start} મુજબ
      તે અનુપ્રયોગનું પરિણામ પણ નથી. આમ~$M$ માત્ર~$N$ના
      અમૂર્તીકરણથી જ રચાઈ શકે છે. આગમનની ધારણાથી~$N$નો
      રચનાક્રમ પોતે અનન્ય છે.
    \item $M$ એ~$(PQ)$ સ્વરૂપનું છે, જ્યાં~$P$ અને~$Q$
      પદો છે. તેની શરૂઆત કૌંસથી થાય છે; તેથી તે
      \olref[trm]{defn:term}\olref[trm]{defn:term-var}
      મુજબ પણ રચાઈ શકે નહીં. \olref{lem:term-start} મુજબ
      $P$ની શરૂઆત~$\lambd$થી થઈ શકતી નથી; તેથી~$(PQ)$
      અમૂર્તીકરણનું પરિણામ પણ હોઈ શકે નહીં. હવે ધારો કે
      $M$ને અનુપ્રયોગથી રચવાની બીજી રીત હોય, એટલે કે તે
      $(P'Q')$ સ્વરૂપનું પણ હોય. તો~$P$ એ~$P'$નો યોગ્ય
      આરંભિક ખંડ હશે (અથવા તેનો ઊલટો); પરંતુ
      \olref{lem:initial} મુજબ આ અશક્ય છે. તેથી~$P$ અને
      $Q$ અનન્ય રીતે નક્કી થાય છે, અને આગમનની ધારણાથી
      $P$ અને~$Q$ બંનેના રચનાક્રમ પણ અનન્ય છે.
  \end{enumerate}
\end{proof}

ઉપરના પ્રસ્તાવને વધુ સહેલાઈથી વાંચી શકાય એમ આ રીતે કહી શકાય:
\begin{prop}
  પદ~$M$ નીચેના સ્વરૂપોમાંથી માત્ર એકનું હોઈ શકે છે:
  \begin{enumerate}
    \item $x$, જ્યાં~$x$ એ~$M$થી અનન્ય રીતે નક્કી થતો ચલ છે.
    \item $(\lambd[x][N])$, જ્યાં~$x$ ચલ અને~$N$ બીજું
      પદ છે; બંને~$M$થી અનન્ય રીતે નક્કી થાય છે.
    \item $(PQ)$, જ્યાં~$P$ અને~$Q$ એ~$M$થી અનન્ય રીતે
      નક્કી થતાં બે પદો છે.
  \end{enumerate}
\end{prop}
% END OLP-0358
\endgroup


\begingroup
\makeatletter\def\ol@sectioncs{section}\makeatother

% BEGIN OLP-0359
\olfileid{lam}{syn}{abb}
\olsection{સંક્ષિપ્ત વાક્યરચના}
\phantomsection\label{guunit:OLP-0359}


\olref[trm]{defn:term} મુજબ વ્યાખ્યાયિત પદો લખવામાં ક્યારેક
ભારે પડે છે; તેથી વધુ સંક્ષિપ્ત વાક્યરચના ઉપયોગી છે. આ
સંક્ષિપ્ત સંકેતનમાં પણ પદોને અનન્ય રીતે વાંચી શકાય તેની
કાળજી લેવી જોઈએ. વ્યાપકપણે વપરાતું એક સ્વરૂપ
\emph{સંક્ષિપ્ત પદો} કહેવાય છે અને તે આ પ્રમાણે છે.

\begin{enumerate}
\item કૌંસ છોડવામાં આવે ત્યારે અનુપ્રયોગ ડાબેથી જમણે થાય છે.
  દાખલા તરીકે, $M$, $N$, $P$ અને~$Q$ પદો હોય તો
  $MNPQ$ એ~$(((MN)P)Q)$નું સંક્ષિપ્ત સ્વરૂપ છે.
\item ફરી, કૌંસ છોડવામાં આવે ત્યારે લૅમ્બડા અમૂર્તીકરણનો
  વ્યાપ શક્ય તેટલો મોટો લેવાય છે. દાખલા તરીકે,
  $\lambd[x][MNP]$ને~$(\lambd[x][MNP])$ એમ વાંચવું.
\item એક લૅમ્બડાથી અનેક ચલોનું અમૂર્તીકરણ કરી શકાય છે.
  દાખલા તરીકે, $\lambd[xyz][M]$ એ
  $\lambd[x][\lambd[y][\lambd[z][M]]]$નું સંક્ષિપ્ત સ્વરૂપ છે.
\end{enumerate}

દાખલા તરીકે,
\[
\lambd[xy][xxyx \lambd[z][xz]]
\]
એ નીચેનું સંક્ષિપ્ત સ્વરૂપ છે:
\[
(\lambd[x][(\lambd[y][((((xx)y)x)(\lambd[z][(xz)]))])]).
\]

\begin{prob}
સંક્ષિપ્ત પદ~$\lambd[g][(\lambd[x][g (x x)])
  \lambd[x][g (x x)]]$ને પૂર્ણ સ્વરૂપે લખો.
\end{prob}
% END OLP-0359
\endgroup


\begingroup
\makeatletter\def\ol@sectioncs{section}\makeatother

% BEGIN OLP-0360
\olfileid{lam}{syn}{fv}
\olsection{મુક્ત ચલો}
\phantomsection\label{guunit:OLP-0360}


લૅમ્બડા કલન વિધેયો વિશે છે, અને લૅમ્બડા અમૂર્તીકરણથી વિધેયો
મળે છે. સાહજિક રીતે, $\lambd[x][M]$ એવું વિધેય છે જેમાં વિધેયનો
આર્ગ્યુમેન્ટ~$x$ને સોંપીએ ત્યારે મૂલ્યો~$M$ નક્કી કરે છે. પરંતુ
$M$માં~$x$ની દરેક ઘટના અહીં સંબંધિત નથી: જો~$M$માં બીજું
અમૂર્તીકરણ~$\lambd[x][N]$ હોય તો~$N$માંની~$x$ની ઘટનાઓ
$\lambd[x][N]$ માટે સંબંધિત છે, $\lambd[x][M]$ માટે નહીં. તેથી
$\lambd[x][M]$ની અંદરનું લૅમ્બડા અમૂર્તીકરણ~$\lambd[x]$, $M$માંની
$x$ની જે ઘટનાઓ બીજા કોઈ લૅમ્બડા અમૂર્તીકરણથી અગાઉથી બદ્ધ નથી
તેને \emph{બાંધે} છે---એ~$M$માંની~$x$ની \emph{મુક્ત} ઘટનાઓ છે.

\begin{defn}[વ્યાપ]
પદ~$N$ની અંદર~$\lambd[x][M]$ની ઘટના હોય તો તેમાંની~$M$ની
તેને અનુરૂપ ઘટનાને~$\lambd[x]$નો \emph{વ્યાપ} કહેવાય છે.
\sourcecorrection{OLLAM-011}{સ્થિર અંગ્રેજી મૂળમાં વ્યાપ તરીકે
પદ~$N$ની અનુરૂપ ઘટના લખાઈ છે; વ્યાપ તો અંદરના પદ~$M$ની
અનુરૂપ ઘટના છે. અહીં~$M$ મૂક્યું છે.}
\end{defn}

\begin{defn}[મુક્ત અને બદ્ધ ઘટના]
  પદ~$M$માં ચલ~$x$ની કોઈ ઘટના $\lambd[x]$ના વ્યાપમાં ન હોય તો
  તે \emph{મુક્ત} છે, અને હોય તો \emph{બદ્ધ} છે.
  $\lambd[x][M]$માં ચલ~$x$ની કોઈ ઘટના આરંભના~$\lambd[x]$થી
  ત્યારે અને માત્ર ત્યારે જ બદ્ધ થાય છે જ્યારે~$M$માંની~$x$ની
  તે ઘટના મુક્ત હોય.
\end{defn}

\begin{ex}
  $\lambd[x][x y]$માં $x$ અને~$y$ બંને~$\lambd[x]$ના વ્યાપમાં
  છે; તેથી~$x$ એ~$\lambd[x]$થી બદ્ધ છે. $y$ કોઈ
  $\lambd[y]$ના વ્યાપમાં નથી, તેથી તે મુક્ત છે.
  $\lambd[x][x x]$માં~$x$ની બંને ઘટનાઓ~$\lambd[x]$થી બદ્ધ છે,
  કારણ કે~$xx$માં બંને મુક્ત છે.
  $((\lambd[x][xx])x)$માં~$x$ની છેલ્લી ઘટના મુક્ત છે, કારણ કે
  તે કોઈ~$\lambd[x]$ના વ્યાપમાં નથી.
  $\lambd[x][(\lambd[x][x])x]$માં પહેલા~$\lambd[x]$નો વ્યાપ
  $(\lambd[x][x])x$ છે અને બીજા~$\lambd[x]$નો વ્યાપ
  $x$ની છેલ્લેથી બીજી ઘટના છે. $(\lambd[x][x])x$માં
  $x$ની છેલ્લી ઘટના મુક્ત અને છેલ્લેથી બીજી બદ્ધ છે.
  આમ, $\lambd[x][(\lambd[x][x])x]$માં~$x$ની છેલ્લેથી બીજી
  ઘટના બીજા~$\lambd[x]$થી અને છેલ્લી ઘટના પહેલા
  $\lambd[x]$થી બદ્ધ છે.
\end{ex}

પદ~$P$માં ચલોની બધી ઘટનાઓ તપાસીને તેના મુક્ત ચલોનો ગણ
મેળવી શકાય છે. આ ગણ~$\FV{P}$થી દર્શાવીએ છીએ; તેની સાહજિક
વ્યાખ્યા આ છે:

\begin{defn}[પદના મુક્ત ચલો] \ollabel{def:fv}
  પદના \emph{મુક્ત ચલો}નો ગણ આગમનથી આ રીતે વ્યાખ્યાયિત થાય છે:
  \begin{enumerate}
    \item $\FV{x} = \{x\}$ \ollabel{def:fv1}
    \item $\FV{\lambd[x][N]} = \FV{N} \setminus \{x\}$    \ollabel{def:fv2}
    \item $\FV{PQ} = \FV{P} \cup \FV{Q}$ \ollabel{def:fv3}
  \end{enumerate}
\end{defn}

\begin{prob}
  \begin{enumerate}
  \item આ પદમાં~$\lambd[g]$ અને બે~$\lambd[x]$ના વ્યાપ શોધો:
    $\lambd[g][(\lambd[x][g (x x)]) \lambd[x][g (x x)]]$.
  \item $\lambd[g][(\lambd[x][g (x x)]) \lambd[x][g (x x)]]$માં
    ચલોની બધી ઘટનાઓ બદ્ધ છે? તે અનુક્રમે કયા અમૂર્તીકરણોથી
    બદ્ધ થાય છે?
    \item $\FV{\lambd[x][(\lambd[y][(\lambd[z][x y]) z]) y]}$
    શોધો.
  \end{enumerate}
\end{prob}

\begin{explain}
મુક્ત ચલ બહારના જગત (અર્થાત્~\emph{પરિવેશ}) તરફના નિર્દેશ જેવો
છે. તેથી મુક્ત ચલો ધરાવતું પદ આંશિક રીતે નિર્ધારિત પદ તરીકે
જોઈ શકાય: તેનું વર્તન આપણે પરિવેશ કેવી રીતે ગોઠવીએ તેના પર
આધારિત છે. દાખલા તરીકે, પદ~$\lambd[x][f x]$ આર્ગ્યુમેન્ટ~$x$
સ્વીકારે છે અને તેના પર~$f$ લાગુ કરીને મળતું મૂલ્ય આપે છે.
અહીં ચલ~$f$ મુક્ત છે. પદનું મૂલ્ય તે જે પરિવેશમાં છે તેના પર,
ખાસ કરીને તે પરિવેશમાં~$f$ના મૂલ્ય પર, આધારિત છે.

આ પદ પર અમૂર્તીકરણ લાગુ કરીએ તો
$\lambd[f][\lambd[x][f
    x]]$ મળે છે. હવે આ પદ પરિવેશના ચલ~$f$ પર આધારિત નથી,
કારણ કે તે બે આર્ગ્યુમેન્ટ સ્વીકારીને પહેલાને બીજા પર લાગુ
કરવાનું પરિણામ આપતું વિધેય દર્શાવે છે. પરિવેશમાં~$f$ બદલીશું
તો પણ આ પદના વર્તન પર અસર નહીં થાય: પદ માત્ર આર્ગ્યુમેન્ટ
તરીકે આપેલું મૂલ્ય વાપરશે, પરિવેશમાંના~$f$નું મૂલ્ય નહીં.
\end{explain}

\begin{defn}[બંધ પદ, સંયોજક]
  મુક્ત ચલો વિનાના પદને \emph{બંધ પદ} અથવા
  \emph{સંયોજક} કહે છે.
\end{defn}

\begin{lem}\ollabel{lem:fv}
  \begin{enumerate}
    \item \ollabel{lem:fv-abs} $y \neq x$ હોય તો
      $y \in \FV{\lambd[x][N]}$ ત્યારે અને માત્ર ત્યારે જ
      જ્યારે~$y \in \FV{N}$.
    \item \ollabel{lem:fv-app} $y \in \FV{PQ}$ ત્યારે અને માત્ર
      ત્યારે જ જ્યારે~$y \in \FV{P}$ અથવા~$y
      \in \FV{Q}$.
    \end{enumerate}
\end{lem}

\begin{proof}
  અભ્યાસપ્રશ્ન.
\end{proof}

\begin{prob}
  \olref[lam][syn][fv]{lem:fv} સાબિત કરો.
\end{prob}
% END OLP-0360
\endgroup


\begingroup
\makeatletter\def\ol@sectioncs{section}\makeatother

% BEGIN OLP-0361
\olfileid{lam}{syn}{sub}
\olsection{સ્થાનાપન્ન}
\phantomsection\label{guunit:OLP-0361}


\begin{explain}
મુક્ત ચલો પરિવેશના ચલો તરફ નિર્દેશ કરે છે; તેથી મુક્ત ચલની
જગ્યાએ કોઈ નિશ્ચિત મૂલ્ય વાપરવું અર્થપૂર્ણ છે. દાખલા તરીકે,
$\lambd[x][f x]$માંના~$f$ની જગ્યાએ નિશ્ચિત પદ, જેમ કે
ઓળખ-વિધેય~$\lambd[y][y]$, મૂકવું હોય. પરિણામે
$\lambd[x][(\lambd[y][y]) x]$ મળે છે. મુક્ત ચલોની જગ્યાએ
લૅમ્બડા પદો મૂકવાની પ્રક્રિયા \emph{સ્થાનાપન્ન} કહેવાય છે.
\end{explain}

\begin{defn}[સ્થાનાપન્ન] \ollabel{defn:substitution}
  પદ~$M$માં ચલ~$x$ની જગ્યાએ પદ~$N$નું \emph{સ્થાનાપન્ન},
  $\Subst{M}{N}{x}$, આગમનથી આ રીતે વ્યાખ્યાયિત થાય છે:
  \begin{enumerate}
    \item $\Subst{x}{N}{x} = N$. \ollabel{defn:substitution-1}
    \item $\Subst{y}{N}{x}  = y$ જો $x \neq y$ હોય.
      \ollabel{defn:substitution-2}
    \item $\Subst{PQ}{N}{x}  = (\Subst{P}{N}{x}) (\Subst{Q}{N}{x})$.
      \ollabel{def:substitution-3}
    \item $\Subst{(\lambd[y][P])}{N}{x} = \lambd[y][\Subst{P}{N}{x}]$,
      જો $x \neq y$ અને $y \notin \FV{N}$ હોય તથા અંદરનું
      સ્થાનાપન્ન વ્યાખ્યાયિત હોય. જો $x=y$ હોય અથવા
      $x\notin\FV{P}$ હોય તો આખું અમૂર્તીકરણ યથાવત રહે છે;
      બાકીના કિસ્સામાં સ્થાનાપન્ન અવ્યાખ્યાયિત છે.
      \ollabel{defn:substitution-4}
  \end{enumerate}
\sourcecorrection{OLLAM-012}{સ્થિર અંગ્રેજી મૂળનો ચોથો નિયમ
બાકીના બધા કિસ્સાઓને અવ્યાખ્યાયિત કહે છે, જ્યારે તરત પછીનું
ઉદાહરણ બદ્ધ ચલને બદલી ન નાખતાં અમૂર્તીકરણ યથાવત રાખવાનું કહે
છે. અહીં બાંધનાર ચલ જ બદલવાનું હોય, અથવા બદલવાના ચલની
મુક્ત ઘટના અંદરના પદમાં ન હોય, તો અમૂર્તીકરણ યથાવત રાખ્યું
છે; આથી નિષ્પ્રયોજન સ્થાનાપન્નમાં પણ ખોટો બંધન-અવરોધ ઊભો
થતો નથી.}
\end{defn}

\begin{explain}
\olref{defn:substitution}\olref{defn:substitution-4}માં
$x \neq y$ની શરત એટલા માટે છે કે~$M$માં ચલ~$x$ની
\emph{બદ્ધ} ઘટનાઓને~$N$થી બદલવી નથી. દાખલા તરીકે,
$\Subst{\lambd[x][x]}{y}{x}$ ગણીએ તો પરિણામ
$\lambd[x][y]$ નહીં પરંતુ~$\lambd[x][x]$ જ હોવું જોઈએ.

$\lambd[y][P]$માં~$x$ની જગ્યાએ~$N$ મૂકતાં આપણે
$y \notin \FV{N}$ પણ માગીએ છીએ. દાખલા તરીકે,
$\lambd[y][x]$માં~$x$ની જગ્યાએ~$y$ મૂકી શકાતું નથી;
અર્થાત્~$\Subst{\lambd[y][x]}{y}{x}$ અવ્યાખ્યાયિત છે.
એમ કરવાથી~$\lambd[y][y]$ મળશે, જે આર્ગ્યુમેન્ટ સ્વીકારીને
તે જ પાછું આપતું વિધેય છે. પરંતુ~$\lambd[y][x]$ હંમેશા
પદ~$x$ (અથવા~$x$ જેનો નિર્દેશ કરે તે) પાછું આપતું વિધેય છે.
આથી આપણને ખરેખર એવું વિધેય જોઈએ જે આર્ગ્યુમેન્ટ સ્વીકારે,
તેને અવગણે અને પરિવેશના ચલ~$y$ને પાછું આપે. તે યોગ્ય રીતે
કરવા પહેલાં બદ્ધ ચલ~$y$નું ``નામ બદલવું'' પડે.
\end{explain}

\begin{prob}
  નીચેના સ્થાનાપન્નોનાં પરિણામો શું છે?
  \begin{enumerate}
  \item $\Subst{\lambd[y][x(\lambd[w][vwx])]}{(uv)}{x}$
  \item $\Subst{\lambd[y][x(\lambd[x][x])]}{(\lambd[y][xy])}{x}$
  \item $\Subst{y(\lambd[v][xv])}{(\lambd[y][vy])}{x}$
  \end{enumerate}
\end{prob}

\begin{thm} \ollabel{thm:notinfv}
  $x \notin \FV{M}$ હોય અને સ્થાનાપન્ન વ્યાખ્યાયિત હોય તો
  $\FV{\Subst{M}{N}{x}} = \FV{M}$.
\end{thm}

\begin{proof}
  $M$ના રચનાક્રમ પર આગમનથી.
  \begin{enumerate}
  \item $M$ ચલ હોય: અભ્યાસપ્રશ્ન.
  \item $M$ એ~$(PQ)$ સ્વરૂપનું હોય: અભ્યાસપ્રશ્ન.
  \item $M$ એ~$\lambd[y][P]$ સ્વરૂપનું હોય. જો બદલવાનો ચલ
    બાંધનાર ચલ જ હોય તો નિયમ અનુસાર પદ યથાવત રહે છે અને
    દાવો તરત મળે છે. અન્યથા,
    $\Subst{\lambd[y][P]}{N}{x}$ વ્યાખ્યાયિત હોવાથી તે
    $\lambd[y][\Subst{P}{N}{x}]$ છે. તેથી
    $\Subst{P}{N}{x}$ વ્યાખ્યાયિત છે; વળી $x \neq y$
    અને~$x \notin \FV{P}$. પછી:
    \begin{multline*}
      \FV{\Subst{\lambd[y][P]}{N}{x}} \\
    \begin{aligned}
      & = \FV{\lambd[y][\Subst{P}{N}{x}]} &&
      \text{\olref{defn:substitution-4} મુજબ}\\
      & = \FV{\Subst{P}{N}{x}} \setminus \{y\} &&
      \text{\olref[fv]{def:fv}\olref[fv]{def:fv2} મુજબ}\\
      & = \FV{P} \setminus \{y\} && \text{આગમનની ધારણાથી} \\
      & = \FV{\lambd[y][P]} && \text{\olref[fv]{def:fv}\olref[fv]{def:fv2} મુજબ}
    \end{aligned}
    \end{multline*}
  \end{enumerate}
\sourcecorrection{OLLAM-013}{સ્થિર અંગ્રેજી મૂળની આ સાબિતીમાં
$x\notin\FV{Q}$ લખાયું છે, પરંતુ આ કિસ્સામાં~$M$ એ
$\lambd[y][P]$ છે; જરૂરી ધારણા~$x\notin\FV{P}$ છે.
બાંધનાર ચલ જ બદલવાનું હોય તે યથાવત રહેતો કિસ્સો પણ
અલગ બતાવ્યો છે અને સમીકરણ-શ્રેણીમાંનું બેવડું સમાનચિહ્ન દૂર
કર્યું છે.}
\end{proof}

\begin{prob}
  \olref[lam][syn][sub]{thm:notinfv}ની સાબિતી પૂર્ણ કરો.
\end{prob}

\begin{thm} \ollabel{thm:infv}
  $x \in \FV{M}$ હોય અને સ્થાનાપન્ન વ્યાખ્યાયિત હોય તો
  $\FV{\Subst{M}{N}{x}} = (\FV{M} \setminus
  \{x\}) \cup \FV{N}$.
\sourcecorrection{OLLAM-014}{સ્થિર અંગ્રેજી મૂળની ધારણામાં
$x\in\FV{M})$ લખાયેલો વધારાનો બંધ કૌંસ દૂર કર્યો છે.}
\end{thm}

\begin{proof}
  $M$ના રચનાક્રમ પર આગમનથી.
  \begin{enumerate}
  \item $M$ ચલ હોય: અભ્યાસપ્રશ્ન.
  \item $M$ એ~$PQ$ સ્વરૂપનું હોય: $\Subst{(PQ)}{N}{x}$
    વ્યાખ્યાયિત હોવાથી તે
    $(\Subst{P}{N}{x})(\Subst{Q}{N}{x})$ છે, જ્યાં બંને
    સ્થાનાપન્ન વ્યાખ્યાયિત છે. વળી~$x \in \FV{PQ}$
    હોવાથી~$x \in \FV{P}$ અથવા~$x \in \FV{Q}$ અથવા
    બંને સાચાં છે. બાકીનો ભાગ અભ્યાસપ્રશ્ન છે.
    \sourcecorrection{OLLAM-015}{સ્થિર અંગ્રેજી મૂળમાં
    અનુપ્રયોગના આ કિસ્સામાં~$\Subst{(PQ)}{N}{y}$ છે,
    જ્યારે આખા પ્રમેયમાં બદલવાનો ચલ~$x$ છે. અહીં~$x$
    મૂક્યું છે.}
  \item $M$ એ~$\lambd[y][P]$ સ્વરૂપનું હોય.
    $\Subst{\lambd[y][P]}{N}{x}$ વ્યાખ્યાયિત હોવાથી તે
    $\lambd[y][\Subst{P}{N}{x}]$ છે; અહીં
    $\Subst{P}{N}{x}$ વ્યાખ્યાયિત છે, $x \neq y$ અને
    $y \notin \FV{N}$. વળી, $x \in
    \FV{\lambd[y][P]}$ હોવાથી~$x \in \FV{P}$ પણ છે.
    હવે:
    \sourcecorrection{OLLAM-016}{સ્થિર અંગ્રેજી મૂળમાં અહીં
    $y\in\FV{\lambd[x][P]}$ પરથી~$y\in\FV{P}$ કહે છે;
    આ કિસ્સાનું પદ તો~$\lambd[y][P]$ છે અને જરૂરી
    ધારણા~$x\in\FV{\lambd[y][P]}$, તેથી~$x\in\FV{P}$ છે.}
    \begin{multline*}
      \FV{\Subst{(\lambd[y][P])}{N}{x}} \\
      \begin{aligned}
      & = \FV{\lambd[y][\Subst{P}{N}{x}]} \\
      & = \FV{\Subst{P}{N}{x}} \setminus \{y\} \\
      & = ((\FV{P} \setminus \{x\}) \cup \FV{N}) \setminus \{y\}
       && \text{આગમનની ધારણાથી} \\
      & = (\FV{P} \setminus \{x, y\}) \cup \FV{N}
       && y \notin \FV{N} \\
      & = (\FV{\lambd[y][P]} \setminus \{x\}) \cup \FV{N}
      \end{aligned}
      \end{multline*}
    \sourcecorrection{OLLAM-017}{સ્થિર અંગ્રેજી મૂળની
    આગમન-શ્રેણીમાં~$\FV{P}\setminus\{y\}$ અને
    $\FV{N}\setminus\{x\}$ ખોટી જગ્યાએ છે તથા કૌંસ
    અધૂરા છે. અહીં પહેલા~$x$ દૂર કર્યા પછી સંઘમાંથી~$y$
    દૂર કર્યો છે; છેલ્લું સરળીકરણ~$y\notin\FV{N}$થી
    મળે છે, મૂળમાં લખેલા~$x\notin\FV{N}$થી નહીં.
    સમીકરણ-શ્રેણીનું બેવડું સમાનચિહ્ન પણ દૂર કર્યું છે.}
  \end{enumerate}
\end{proof}

\begin{prob}
  \olref[lam][syn][sub]{thm:infv}ની સાબિતી પૂર્ણ કરો.
\end{prob}

\begin{thm}\ollabel{thm:clr}
  સ્થાનાપન્ન વ્યાખ્યાયિત હોય અને~$x \notin \FV{N}$ હોય તો
  $x \notin \FV{\Subst{M}{N}{x}}$.
\end{thm}

\begin{proof}
  અભ્યાસપ્રશ્ન.
\end{proof}

\begin{prob}
  \olref[lam][syn][sub]{thm:clr} સાબિત કરો.
\end{prob}

\begin{thm}\ollabel{thm:inv}
  $\Subst{M}{y}{x}$ વ્યાખ્યાયિત હોય અને
  $y \notin \FV{M}$ હોય તો
  $\Subst{\Subst{M}{y}{x}}{x}{y} = M$.
\end{thm}

\begin{proof}
  $M$ના રચનાક્રમ પર આગમનથી.
  \begin{enumerate}
  \item $M$ ચલ~$z$ હોય: અભ્યાસપ્રશ્ન.
  \item $M$ એ~$(PQ)$ સ્વરૂપનું હોય. પછી:
    \begin{align*}
      \Subst{\Subst{(PQ)}{y}{x}}{x}{y}
      &=\Subst{((\Subst{P}{y}{x})(\Subst{Q}{y}{x}))}{x}{y} \\
      &= (\Subst{\Subst{P}{y}{x}}{x}{y})(\Subst{\Subst{Q}{y}{x}}{x}{y}) \\
      &= (PQ) \text{આગમનની ધારણાથી}
    \end{align*}
  \item $M$ એ~$\lambd[z][N]$ સ્વરૂપનું હોય. બદલવાના ચલની
    મુક્ત ઘટના જ ન હોય તો બંને સ્થાનાપન્ન પદને યથાવત રાખે છે.
    બાકીના કિસ્સામાં~$\Subst{\lambd[z][N]}{y}{x}$
    વ્યાખ્યાયિત હોવાથી~$z \neq y$ અને~$z\neq x$ છે.
    તેથી:
    \begin{align*}
      \Subst{\Subst{(\lambd[z][N])}{y}{x}}{x}{y}\\
      & = \Subst{(\lambd[z][\Subst{N}{y}{x}])}{x}{y} \\
      & = \lambd[z][\Subst{\Subst{N}{y}{x}}{x}{y}] \\
      &= \lambd[z][N] \text{આગમનની ધારણાથી}
    \end{align*}
    \sourcecorrection{OLLAM-018}{સ્થિર અંગ્રેજી મૂળ પહેલું
    સ્થાનાપન્ન વ્યાખ્યાયિત હોવાથી સીધું~$z\neq y$
    માને છે; બદલવાના ચલની મુક્ત ઘટના ન હોય તો આ જરૂરી
    નથી. તે યથાવત રહેતો કિસ્સો અલગ લીધો છે. બાકીનું
    પ્રદર્શન બંને બંધન-અથડામણો ટળે એવા~$z\neq x,y$
    કિસ્સામાં છે.}
  \end{enumerate}
\end{proof}

\begin{prob}
  \olref[lam][syn][sub]{thm:inv}ની સાબિતી પૂર્ણ કરો.
\end{prob}
% END OLP-0361
\endgroup


\begingroup
\makeatletter\def\ol@sectioncs{section}\makeatother

% BEGIN OLP-0362
\olfileid{lam}{syn}{alp}
\olsection{$\alpha$-રૂપાંતરણ}
\phantomsection\label{guunit:OLP-0362}


$\lambd[x][x]$ અને~$\lambd[y][y]$ વચ્ચે શું સંબંધ છે?
બંને ઓળખ-વિધેય દર્શાવે છે. અલબત્ત, વાક્યરચનાની દૃષ્ટિએ તે
જુદાં પદો છે. ફરક માત્ર બદ્ધ ચલના નામમાં છે; એકમાંના બદ્ધ
ચલનું ``નામ બદલીને'' બીજું મળે છે. તેને
\emph{$\alpha$-રૂપાંતરણ} કહે છે.

\begin{defn}[બદ્ધ ચલનું નામ બદલવું, $\aconvone$]
  પદ~$M$માં~$\lambd[x][N]$ની કોઈ ઘટના હોય, $x\neq y$,
  $y \notin
  \FV{N}$ હોય અને~$\Subst{N}{y}{x}$ વ્યાખ્યાયિત હોય તો
  તે ઘટનાને
  \begin{equation*}
    \lambd[y][\Subst{N}{y}{x}]
  \end{equation*}
  વડે બદલીને~$M'$ મળે તેને \emph{બદ્ધ ચલનું નામ બદલવું}
  કહીએ છીએ અને~$M \redone[\alpha] M'$ લખીએ છીએ.
\sourcecorrection{OLLAM-019}{સ્થિર અંગ્રેજી મૂળની પહેલી
વ્યાખ્યામાં~$x\neq y$ની શરત નથી, પરંતુ આગળની બંને સમકક્ષ
વ્યાખ્યાઓમાં તે છે. અહીં નામ ખરેખર બદલાય તે માટે એ શરત
સ્પષ્ટ મૂકી છે.}
\end{defn}

\begin{defn}[સંબંધની રચના-સુસંગતતા]
  પદો પરનો સંબંધ~$R$ નીચેની શરતો સંતોષે તો તેને
  \emph{રચના-સુસંગત} કહીએ છીએ:
  \begin{enumerate}
  \item જો $R N N'$ હોય તો $R \lambd[x][N] \lambd[x][N']$.
  \item જો $R P P'$ હોય તો $R (PQ) (P'Q)$.
  \item જો $R Q Q'$ હોય તો $R (PQ) (PQ')$.
  \end{enumerate}
\end{defn}

આથી વ્યાખ્યા ફરી આ રીતે કહી શકાય:
\begin{defn}[બદ્ધ ચલનું નામ બદલવું, $\aconvone$]
  \emph{બદ્ધ ચલનું નામ બદલવાનો} સંબંધ~($\redone[\alpha]$)
  પદો પરનો સૌથી નાનો રચના-સુસંગત સંબંધ છે જે નીચેની
  શરત સંતોષે છે:
  \begin{align*}
    &\lambd[x][N] \redone[\alpha] \lambd[y][\Subst{N}{y}{x}] && \text{જો
      $x \neq y$, $y \notin \FV{N}$ હોય} \\
    & &&\text{અને $\Subst{N}{y}{x}$ વ્યાખ્યાયિત હોય}
  \end{align*}
\end{defn}

અહીં ``સૌથી નાનો'' એટલે સંબંધમાં રચના-સુસંગતતા અને વધારાની
શરતથી જરૂરી બનેલી જોડીઓ જ હોય, બીજી કોઈ નહીં. તેથી સંબંધને
આગમનથી પણ વ્યાખ્યાયિત કરી શકાય:

\begin{defn}[બદ્ધ ચલનું નામ બદલવું, $\aconvone$] \ollabel{defn:aconvone}
  \emph{બદ્ધ ચલનું નામ બદલવાનો} સંબંધ~($\aconvone$)
  આગમનથી આ રીતે વ્યાખ્યાયિત થાય છે:
  \begin{enumerate}
  \item જો $N \aconvone N'$ હોય તો $\lambd[x][N] \aconvone
    \lambd[x][N']$. \ollabel{defn:aconvone1}
  \item જો $P \aconvone P'$ હોય તો $(PQ) \aconvone (P'Q)$.
    \ollabel{defn:aconvone2}
  \item જો $Q \aconvone Q'$ હોય તો $(PQ) \aconvone (PQ')$.
    \ollabel{defn:aconvone3}
  \item જો $x \neq y$, $y \notin \FV{N}$ અને~$\Subst{N}{y}{x}$
    વ્યાખ્યાયિત હોય તો
    $\lambd[x][N] \redone[\alpha] \lambd[y][\Subst{N}{y}{x}]$.
    \ollabel{defn:aconvone4}
  \end{enumerate}
\end{defn}

આ વ્યાખ્યાઓ સમકક્ષ છે, પરંતુ તેની સાબિતી અભ્યાસપ્રશ્ન તરીકે
છોડીએ છીએ. હવે પછી આગમનાત્મક વ્યાખ્યા વાપરીશું.

\begin{defn}[$\alpha$-રૂપાંતરણ, $\aconv$]
  \emph{$\alpha$-રૂપાંતરણ}~($\aconv$) એ પદો પરનો
  $\aconvone$ને સમાવતા સ્વવાચક અને પરંપરિત સંબંધોમાં
  સૌથી નાનો સંબંધ છે.
\end{defn}

અગાઉની જેમ, ``સૌથી નાનો'' એટલે સંબંધમાં સ્વવાચકતા,
પરંપરિતતા અને~$\aconvone$ને કારણે જરૂરી બનેલી જોડીઓ જ હોય.
તેમાંથી નીચેની સમકક્ષ વ્યાખ્યા મળે છે:
\sourcecorrection{OLLAM-020}{સ્થિર અંગ્રેજી મૂળની આ
સમજૂતીમાં ``સૌથી નાના'' સંબંધ માટે પરંપરિતતા અને
$\aconvone$નું નામ છે, પરંતુ વ્યાખ્યામાં જરૂરી સ્વવાચકતાનો
ઉલ્લેખ છૂટી ગયો છે. અહીં તે ઉમેર્યો છે.}

\begin{defn}[$\alpha$-રૂપાંતરણ, $\aconv$] \ollabel{defn:aconv}
  \emph{$\alpha$-રૂપાંતરણ}~($\aconv$) આગમનથી આ રીતે
  વ્યાખ્યાયિત થાય છે:
  \begin{enumerate}
  \item જો $P \aconv Q$ અને~$Q \aconv R$ હોય તો
    $P \aconv R$. \ollabel{defn:aconv1}
  \item જો $P \aconvone Q$ હોય તો $P \aconv Q$.
    \ollabel{defn:aconv2}
  \item $P \aconv P$. \ollabel{defn:aconv3}
  \end{enumerate}
\end{defn}

\begin{ex}
  $\lambd[x][f x]$નું~$\lambd[y][f
  y]$માં $\alpha$-રૂપાંતરણ થાય છે, અને ઊલટું પણ.
  અનૌપચારિક રીતે કહીએ તો બંને આર્ગ્યુમેન્ટ સ્વીકારીને તેના
  પર~$f$ લાગુ કરવાનું પરિણામ આપે છે; અહીં~$f$ પરિવેશનો
  ચલ છે.

  $\lambd[x][f x]$નું~$\lambd[x][g
    x]$માં $\alpha$-રૂપાંતરણ થતું નથી. પહેલું અને બીજું પદ
  અનુક્રમે પરિવેશના ચલો~$f$ અને~$g$ તરફ નિર્દેશ કરે છે.
  તેથી એ જુદાં વિધેયો છે: જ્યાં~$f$ અને~$g$ જુદાં હોય એવા
  પરિવેશોમાં તેમનું વર્તન જુદું હોય છે.
\end{ex}

\begin{prob}
  નીચેની પદ-જોડીઓ $\alpha$-રૂપાંતરણથી એકબીજામાં બદલી શકાય છે?
  \begin{enumerate}
  \item $\lambd[x][\lambd[y][x]]$ અને $\lambd[y][\lambd[x][y]]$
  \item $\lambd[x][\lambd[y][x]]$ અને $\lambd[c][\lambd[b][a]]$
  \item $\lambd[x][\lambd[y][x]]$ અને $\lambd[c][\lambd[b][a]]$
  \end{enumerate}
\sourcecorrection{OLLAM-021}{સ્થિર અંગ્રેજી મૂળમાં આ
અભ્યાસપ્રશ્નની બીજી અને ત્રીજી જોડી એકસરખી છે. અનુમાનથી
કોઈ નવી જોડી ન ઘડતાં બંને મૂળ મુજબ જાળવી છે.}
\end{prob}

\begin{lem}\ollabel{lem:fv-one}
  $P \aconvone Q$ હોય તો $\FV{P} = \FV{Q}$.
\end{lem}

\begin{proof}
  $P \aconvone Q$ની નિષ્પત્તિ પર આગમનથી.
  \begin{enumerate}
  \item છેલ્લો નિયમ~\olref{defn:aconvone4} હોય તો~$P$
    એ~$\lambd[x][N]$ અને~$Q$ એ
    $\lambd[y][\Subst{N}{y}{x}]$ સ્વરૂપના છે; અહીં
    $x \neq y$, $y \notin \FV{N}$ અને~$\Subst{N}{y}{x}$
    વ્યાખ્યાયિત છે. હવે~$x \in \FV{N}$ છે કે નહીં તે
    મુજબ કિસ્સા કરીએ:
    \begin{enumerate}
    \item $x \in \FV{N}$ હોય તો:
      \begin{align*}
        \FV{\lambd[y][\Subst{N}{y}{x}]} &=
        \FV{\Subst{N}{y}{x}} \setminus \{y\} \\
        & = ((\FV{N} \setminus \{x\}) \cup \{y\}) \setminus \{y\}
         && \text{\olref[sub]{thm:infv} મુજબ} \\
        & = \FV{N} \setminus \{x\} \\
        & = \FV{\lambd[x][N]}
      \end{align*}
    \item $x \notin \FV{N}$ હોય તો:
      \begin{align*}
        \FV{\lambd[y][\Subst{N}{y}{x}]}
        & = \FV{\Subst{N}{y}{x}} \setminus \{y\} \\
        & = \FV{N} \setminus \{x\}
         && \text{\olref[sub]{thm:notinfv} મુજબ} \\
        & = \FV{\lambd[x][N]}.
      \end{align*}
    \end{enumerate}
  \item બીજા ત્રણ કિસ્સા અભ્યાસપ્રશ્ન તરીકે છોડીએ છીએ.
  \end{enumerate}
\sourcecorrection{OLLAM-022}{સ્થિર અંગ્રેજી મૂળની આ
સાબિતીમાં~$x\in FV(N)$, $x\notin FV(N)$ અને
$FV\{\Subst{N}{y}{x}\}$ એમ ત્રણ જગ્યાએ મુક્ત-ચલ
સંકેતનું~$\backslash$ છૂટી ગયું છે. અહીં બધે~$\FV{}$
સંકેત વાપર્યો છે.}
\end{proof}

\begin{prob}
  \olref[lam][syn][alp]{lem:fv-one}ની સાબિતી પૂર્ણ કરો.
\end{prob}

\begin{lem}\ollabel{lem:inv}
  $P \aconvone Q$ હોય તો $Q \aconvone P$.
\end{lem}

\begin{proof}
  $P \aconvone Q$ની નિષ્પત્તિ પર આગમનથી.
  \begin{enumerate}
  \item છેલ્લો નિયમ~\olref{defn:aconvone4} હોય તો~$P$
    એ~$\lambd[x][N]$ અને~$Q$ એ
    $\lambd[y][\Subst{N}{y}{x}]$ સ્વરૂપનાં છે; અહીં
    $x \neq y$, $y \notin \FV{N}$ અને~$\Subst{N}{y}{x}$
    વ્યાખ્યાયિત છે. પ્રથમ, \olref[sub]{thm:clr} મુજબ
    $x \notin
    \FV{\Subst{N}{y}{x}}$. \olref[sub]{thm:inv} મુજબ
    $\Subst{\Subst{N}{y}{x}}{x}{y}$ વ્યાખ્યાયિત પણ છે અને
    $N$ બરાબર પણ છે. તેથી~\olref{defn:aconvone4} મુજબ
    $\lambd[y][\Subst{N}{y}{x}]
    \aconvone \lambd[x][\Subst{\Subst{N}{y}{x}}{x}{y}] =
    \lambd[x][N]$.
    \sourcecorrection{OLLAM-023}{સ્થિર અંગ્રેજી મૂળ અહીં
    સ્થાનાપન્ન પછી~$y$ મુક્ત નથી એમ કહે છે; દાખલા તરીકે
    $N=x$ હોય તો પરિણામ~$y$ જ છે. વિપરીત નામ-બદલી માટે
    જરૂરી શરત~$x\notin\FV{\Subst{N}{y}{x}}$ છે, જે
    અગાઉના પ્રમેયથી મળે છે.}
  \end{enumerate}
\end{proof}

\begin{prob}
  \olref[lam][syn][alp]{lem:inv}ની સાબિતી પૂર્ણ કરો.
\end{prob}

\begin{thm}
  $\alpha$-રૂપાંતરણ પદો પરનો સામ્ય સંબંધ છે; એટલે તે
  સ્વવાચક, સંમિત અને પરંપરિત છે.
\end{thm}

\begin{proof}
  \begin{enumerate}
  \item દરેક પદ~$M$ માટે બદ્ધ ચલોનું નામ \emph{શૂન્ય} વાર
    બદલીને~$M$માંથી~$M$ મળે છે.
  \item બદ્ધ ચલોનું નામ બદલવાની કોઈ શ્રેણીથી~$P$નું
    $Q$માં $\alpha$-રૂપાંતરણ થાય તો~$Q$માંથી એ ફેરફારો
    ઊલટા ક્રમે ઉલટાવીને (\olref{lem:inv} મુજબ)~$P$ મળે છે.
  \item એક શ્રેણીથી~$P$નું~$Q$માં અને બીજીથી~$Q$નું~$R$માં
    $\alpha$-રૂપાંતરણ થાય તો પહેલી પછી બીજી શ્રેણી
    લાગુ કરીને~$P$માંથી~$R$ મળે છે.
  \end{enumerate}
\end{proof}

હવે પછી~$M$ અને~$N$ને \emph{$\alpha$-સામ્યવાળાં},
$M
\aeq N$, ત્યારે અને માત્ર ત્યારે જ કહીએ છીએ જ્યારે~$M$નું
$N$માં $\alpha$-રૂપાંતરણ થાય. હમણાં સાબિત કર્યું તેમ,
આ ત્યારે અને માત્ર ત્યારે જ બને છે જ્યારે~$N$નું પણ
$M$માં $\alpha$-રૂપાંતરણ થાય.

\begin{thm}\ollabel{thm:fv}
  $M \aeq N$ હોય તો $\FV{M} = \FV{N}$.
\end{thm}

\begin{proof}
  \olref{lem:fv-one}માંથી સીધું મળે છે.
\end{proof}

\begin{lem}\ollabel{lem:sub:R}
  $R \aeq R'$ હોય અને~$\Subst{M}{R}{y}$ વ્યાખ્યાયિત હોય
  તો~$\Subst{M}{R'}{y}$ પણ વ્યાખ્યાયિત છે અને
  $\Subst{M}{R}{y}$ સાથે $\alpha$-સામ્યવાળું છે.
\end{lem}

\begin{proof}
  અભ્યાસપ્રશ્ન.
\end{proof}

\begin{prob}
  \olref[lam][syn][alp]{lem:sub:R} સાબિત કરો.
\end{prob}

યાદ કરો કે~\olref[sub]{sec}માં કેટલાક કિસ્સામાં સ્થાનાપન્ન
અવ્યાખ્યાયિત છે. પરંતુ પદોના બદ્ધ ચલોનું નામ
$\alpha$-રૂપાંતરણથી બદલીને સ્થાનાપન્ન વ્યાખ્યાયિત બનાવી
શકાય છે. નીચેના પ્રમેય મુજબ પરિણામમાં $\alpha$-સામ્ય
જળવાય છે:

\begin{thm}\ollabel{thm:sub}
  કોઈપણ~$M$, $R$ અને~$y$ માટે એવું~$M'$ મળે છે કે
  $M \aeq M'$ અને~$\Subst{M'}{R}{y}$ વ્યાખ્યાયિત છે.
  વધુમાં, બીજી જોડી~$M'' \aeq M$ અને~$R''$ હોય, જેમાં
  $\Subst{M''}{R''}{y}$ વ્યાખ્યાયિત તથા~$R'' \aeq R$
  હોય તો
  $\Subst{M'}{R}{y} \aeq \Subst{M''}{R''}{y}$.
\end{thm}

\begin{proof}
  $M$ના રચનાક્રમ પર આગમનથી:
  \begin{enumerate}
  \item $M$ ચલ~$z$ હોય: અભ્યાસપ્રશ્ન.
  \item $M$ એ~$\lambd[x][N]$ સ્વરૂપનું હોય. $x$ અને~$y$થી
    જુદો, $z \notin \FV{N}$ અને~$z
    \notin \FV{R}$ હોય એવો ચલ~$z$ પસંદ કરો. આગમનની
    ધારણાથી એવું~$N'$ મળે છે કે~$N' \aeq N$ અને
    $\Subst{N'}{z}{x}$ વ્યાખ્યાયિત છે. ત્યારબાદ
    \olref{defn:aconvone}\olref{defn:aconvone1} મુજબ
    $\lambd[x][N] \aeq \lambd[x][N']$ પણ છે. હવે
    \olref{defn:aconvone}\olref{defn:aconvone4} મુજબ
    $\lambd[x][N']
    \aeq \lambd[z][\Subst{N'}{z}{x}]$ છે. આ પગલું
    લઈ શકાય છે કારણ કે~$z \ne x$,
    $z \notin \FV{N'}$ અને~$\Subst{N'}{z}{x}$
    વ્યાખ્યાયિત છે. પરંતુ માત્ર~$z \neq y$ અને
    $z \notin \FV{R}$થી
    $\Subst{\lambd[z][\Subst{N'}{z}{x}]}{R}{y}$
    વ્યાખ્યાયિત છે એવું ફલિત થતું નથી: અંદરના બંધનોની
    શરતો પણ તપાસવી પડે.
    \sourcecorrection{OLLAM-024}{સ્થિર અંગ્રેજી મૂળ અહીં
    ફક્ત બહારના બાંધનાર~$z$ની અથડામણ ન હોવાથી આખું
    સ્થાનાપન્ન વ્યાખ્યાયિત છે એમ કહે છે. પરંતુ
    $N=\lambd[w][y]$, $R=w$ અને~$x,y,z,w$ જુદા હોય
    ત્યારે~$\Subst{N}{z}{x}=N$ વ્યાખ્યાયિત છે, છતાં
    $\Subst{\lambd[z][N]}{w}{y}$ અંદરના~$w$થી થતી
    અથડામણને કારણે અવ્યાખ્યાયિત છે. અગાઉથી બધાં
    જરૂરી અંદરનાં બંધનોનાં નામ તાજાં કરવાની વધારાની
    દલીલ જોઈએ; સ્થિર મૂળની આ સાબિતી અહીં અધૂરી છે.
    મૂળમાંની બે~$FV$ અભિવ્યક્તિઓને પણ~$\FV{}$ કર્યું છે.}

    વધુમાં, એ જ શરતો સંતોષતાં બીજાં~$N''$ અને~$R''$
    હોય તો સ્ત્રોત નીચેની સાંકળ આપે છે; મધ્યનાં બે
    પગલાં માત્ર આલ્ફા-સામ્ય સુધી જ માન્ય છે:
    \begin{multline*}
      \Subst{(\lambd[z][\Subst{N''}{z}{x}])}{R''}{y} \\
    \begin{aligned}
      &= \lambd[z][\Subst{\Subst{N''}{z}{x}}{R''}{y}] \\
      &\aeq \lambd[z][\Subst{\Subst{N''}{z}{x}}{R}{y}] && \text{
        \olref{lem:sub:R} મુજબ}\\
      &\aeq\lambd[z][\Subst{\Subst{N'}{z}{x}}{R}{y}]
      && \text{આગમનની ધારણાથી}\\
      &=\Subst{(\lambd[z][\Subst{N'}{z}{x}])}{R}{y}
    \end{aligned}
    \end{multline*}
    \sourcecorrection{OLLAM-025}{સ્થિર અંગ્રેજી મૂળની
    સમીકરણ-સાંકળમાં~\olref{lem:sub:R} અને આગમનનો ઉપયોગ
    કરતી જગ્યાએ સાદી સમાનતા લખી છે. સંબંધિત પદો
    સામાન્ય રીતે નામમાં જુદાં હોવાથી ત્યાં માત્ર
    $\alpha$-સામ્ય મળે છે. ઉપરનો અધૂરો વ્યાખ્યાયિતતા-
    તર્ક પૂરો કર્યા વિના આ સાંકળને સંપૂર્ણ સાબિતી
    ગણવી નહીં.}
    \item $M$ એ~$(PQ)$ સ્વરૂપનું હોય: અભ્યાસપ્રશ્ન.
  \end{enumerate}
\end{proof}

\begin{prob}
  \olref[lam][syn][alp]{thm:sub}ની સાબિતી પૂર્ણ કરો.
\end{prob}

\begin{cor}\ollabel{cor:sub}
  કોઈપણ~$M$, $R$ અને~$y$ માટે એવી જોડી~$M'$ અને~$R'$
  મળે છે કે~$M' \aeq M$, $R \aeq R'$ અને
  $\Subst{M'}{R'}{y}$ વ્યાખ્યાયિત છે. વધુમાં, બીજી
  જોડી~$M'' \aeq M$ અને~$R'' \aeq R$ હોય તથા
  $\Subst{M''}{R''}{y}$ વ્યાખ્યાયિત હોય તો
  $\Subst{M'}{R'}{y} \aeq
  \Subst{M''}{R''}{y}$.
  \sourcecorrection{OLLAM-026}{સ્થિર અંગ્રેજી મૂળના
  ઉપસિદ્ધાંતમાં બીજી જોડી માટે~$R''\aeq R$ની શરત
  ગાયબ છે અને તેની જગ્યાએ પહેલી જોડીનું સ્થાનાપન્ન
  ફરી વ્યાખ્યાયિત હોવાનું લખાયું છે. અહીં બીજા
  સ્થાનાપન્નની વ્યાખ્યાયિતતા તથા બંને પ્રતિસ્થાપકોનું
  $\alpha$-સામ્ય સ્પષ્ટ કર્યું છે.}
\end{cor}

\begin{proof}
  \olref{thm:sub}માંથી સીધું મળે છે.
\end{proof}
% END OLP-0362
\endgroup


\begingroup
\makeatletter\def\ol@sectioncs{section}\makeatother

% BEGIN OLP-0363
\olfileid{lam}{syn}{deb}

\olsection{દે બ્રુઇન સૂચકાંક}
\phantomsection\label{guunit:OLP-0363}


$\alpha$-સામ્ય ખૂબ સાહજિક છે, કારણ કે $\alpha$-સામ્યવાળાં
પદોનો ``અર્થ એક જ'' હોય છે. હકીકતમાં, લૅમ્બડા પદો માટે એવી
વાક્યરચના આપી શકાય જેમાં $\alpha$-રૂપાંતરણથી એકબીજામાં
બદલી શકાય તેવાં પદો વચ્ચે ભેદ ન રહે. એમાં સૌથી જાણીતી
રીત ચલોની જગ્યાએ તેમના \emph{દે બ્રુઇન સૂચકાંક} વાપરે છે.

$\lambd[x][M]$ લખીએ ત્યારે વિધેયનો પ્રાચલ~$x$ છે એમ
સ્પષ્ટ કહીએ છીએ, જેથી~$M$માં~$x$ વડે તે પ્રાચલનો નિર્દેશ
કરી શકીએ. પરંતુ દે બ્રુઇન સૂચકાંક પદ્ધતિમાં પ્રાચલોનાં નામ
નથી. વિધેયના મુખ્ય ભાગમાં તેમનો નિર્દેશ, પ્રાચલની ઘટના અને
તેને બાંધતા અમૂર્તીકરણ વચ્ચે અમૂર્તીકરણનાં કેટલાં સ્તર છે
તે દર્શાવતી સંખ્યા વડે થાય છે. દાખલા તરીકે,
$\lambd[x][\lambd[y][y x]]$ જુઓ: બહારનું અમૂર્તીકરણ
ચલ~$x$ને બાંધે છે; અંદરનું અમૂર્તીકરણ ચલ~$y$ને બાંધે છે.
ઉપપદ~$y x$ અંદરના અમૂર્તીકરણના વ્યાપમાં છે. $y$ અને
તેને બાંધતા~$\lambd[y]$ વચ્ચે કોઈ અમૂર્તીકરણ નથી, પરંતુ
$x$ અને તેને બાંધતા~$\lambd[x]$ વચ્ચે એક અમૂર્તીકરણ છે.
આથી~$y x$ માટે~$0\, 1$ અને આખા પદ માટે
$\lambd[][\lambd[][01]]$ લખીએ છીએ.

\begin{defn}
  દે બ્રુઇન પદો આગમનથી આ રીતે વ્યાખ્યાયિત થાય છે:
  \begin{enumerate}
  \item $n$, જ્યાં~$n$ કોઈપણ પ્રાકૃતિક સંખ્યા છે.
  \item $PQ$, જ્યાં~$P$ અને~$Q$ બંને દે બ્રુઇન પદો છે.
  \item $\lambd[][N]$, જ્યાં~$N$ દે બ્રુઇન પદ છે.
  \end{enumerate}
\end{defn}

સામાન્ય લૅમ્બડા પદોમાંથી દે બ્રુઇન સૂચકાંકવાળાં પદોમાં
ઔપચારિક રૂપાંતરણ આ પ્રમાણે છે:
\begin{defn}
  \begin{align*}
    F_\Gamma(x) &= \Gamma(x) \\
    F_\Gamma(PQ) &= F_\Gamma(P)F_\Gamma(Q) \\
    F_\Gamma(\lambd[x][N]) &= \lambd[][F_{x,\Gamma}(N)]
  \end{align*}
  અહીં~$\Gamma$ એ શૂન્યથી સૂચકાંકિત કરેલી ચલોની યાદી છે,
  અને~$\Gamma(x)$ એ~$\Gamma$માં ચલ~$x$નું સ્થાન દર્શાવે છે.
  દાખલા તરીકે, $\Gamma$ એ~$x,y,z$ હોય તો~$\Gamma(x)$ એ
  $0$ અને~$\Gamma(z)$ એ~$2$ છે.

  $x,\Gamma$ એટલે~$\Gamma$ની શરૂઆતમાં~$x$ ઉમેરીને મળતી
  યાદી. અગાઉના ઉદાહરણમાં~$\Gamma$ માટે
  $w,\Gamma$ એ~$w,x,y,z$ છે.
\end{defn}

દે બ્રુઇન પદમાંથી સામાન્ય લૅમ્બડા પદ આ રીતે ફરી મેળવી શકાય:

\begin{defn}
  \begin{align*}
    G_\Gamma(n) &= \Gamma[n] \\
    G_\Gamma(PQ) &= G_\Gamma(P) G_\Gamma(Q) \\
    G_\Gamma(\lambd[][N]) &= \lambd[x][G_{x,\Gamma}(N)]
  \end{align*}
  અહીં પણ~$\Gamma$ એ શૂન્યથી સૂચકાંકિત ચલોની યાદી છે,
  અને~$\Gamma[n]$ એ~$n$મા સ્થાને રહેલો ચલ દર્શાવે છે.
  દાખલા તરીકે, $\Gamma$ એ~$x,y,z$ હોય તો~$\Gamma[1]$
  એ~$y$ છે.

  છેલ્લા સમીકરણમાં ચલ~$x$ એવો કોઈપણ ચલ પસંદ કરાય છે જે
  $\Gamma$માં નથી.
\end{defn}

હવે સાબિતી આપ્યા વિના કેટલાંક પરિણામો નોંધીએ:

\begin{prop}
  $M \aconvone M'$ હોય અને~$\Gamma$ એ~$\FV{M}$ને સમાવી
  લેતી કોઈપણ યાદી હોય તો
  $F_\Gamma(M) \eqs F_\Gamma(M')$.
\end{prop}
% END OLP-0363
\endgroup


\begingroup
\makeatletter\def\ol@sectioncs{section}\makeatother

% BEGIN OLP-0364
\olfileid{lam}{syn}{tr}
\olsection{$\alpha$-સામ્ય વર્ગો તરીકે પદો}
\phantomsection\label{guunit:OLP-0364}


હવે પછી આપણે પદોને $\alpha$-સામ્ય સુધી લઈશું. એટલે કોઈ પદ
લખીએ ત્યારે તેનો $\alpha$-સામ્ય વર્ગ અભિપ્રેત હશે. દાખલા તરીકે,
$\lambd[a][\lambd[b][a c]]$ લખીએ ત્યારે આ પદ સાથે
$\alpha$-સામ્યવાળાં તમામ પદોનો ગણ સમજવાનો છે; તેમાં
$\lambd[a][\lambd[b][a c]]$ અને
$\lambd[b][\lambd[a][b c]]$ વગેરેનો સમાવેશ થાય છે.

અગાઉના વિભાગોમાં $N, Q$ વગેરે અક્ષરો પદ દર્શાવતા હતા;
હવે પછી તેઓ પદોના વર્ગો દર્શાવશે. આગળના અભ્યાસનો વિષય
પદોને બદલે આ વર્ગો હશે. $x, y$ વગેરે અક્ષરો હજી પણ
ચલો જ દર્શાવે છે.

વર્ગ~$M$ના કોઈપણ ઘટક માટે $\rep{M}$ સંકેત વાપરીશું.
એક કરતાં વધુ અવયવ જોઈતા હોય તો $\rep{M}[0], \rep{M}[1]$
વગેરે લખીશું.

વર્ણન સરળ બનાવવા પદો માટેના સંકેતો અહીં ફરી વાપરીશું.
વર્ગો પર નીચેની વ્યાખ્યા કરીએ:
\begin{defn}
  \begin{enumerate}
  \item $\lambd[x][N]$ એટલે $\lambd[x][\rep{N}]$ને
    સમાવતો વર્ગ.
  \item $PQ$ એટલે $\rep{P}\rep{Q}$ને સમાવતો વર્ગ.
  \end{enumerate}
\end{defn}

આ વ્યાખ્યાઓ સુવ્યાખ્યાયિત છે, કારણ કે $\alpha$-રૂપાંતરણ
રચના સાથે સુસંગત છે.

\begin{defn} \ollabel{def:fv}
  $\alpha$-સામ્ય વર્ગ~$M$ના \emph{મુક્ત ચલો}, એટલે
  $FV(M)$, એ $FV(\rep{M})$ તરીકે વ્યાખ્યાયિત થાય છે.
\end{defn}

આ વ્યાખ્યા સુવ્યાખ્યાયિત છે: \olref[alp]{thm:fv} મુજબ
$FV(\rep{M}[0]) = FV(\rep{M}[1])$.

વર્ગોમાં સ્થાનાપન્ન માટે પણ એ જ સંકેત વાપરીએ:
\begin{defn} \ollabel{def:sub}
  $M$માં $y$ની જગ્યાએ $R$નું \emph{સ્થાનાપન્ન}, એટલે
  $\Subst{M}{R}{y}$, એ એવા કોઈપણ $\rep{M}$ અને
  $\rep{R}$ માટે $\Subst{\rep{M}}{\rep{R}}{y}$ને
  સમાવતો $\alpha$-સામ્ય વર્ગ છે, જેના માટે અંદરનું
  સ્થાનાપન્ન વ્યાખ્યાયિત હોય.
  \sourcecorrection{OLLAM-027}{સ્થિર અંગ્રેજી મૂળ
  વર્ગો પરના સ્થાનાપન્નનું પરિણામ સીધું
  $\Subst{\rep{M}}{\rep{R}}{y}$ લખે છે, જે પદ છે,
  વર્ગ નથી. પરિણામ એ પદનો $\alpha$-સામ્ય વર્ગ છે
  એમ અહીં સ્પષ્ટ કર્યું છે.}
\end{defn}

આ પણ \olref[alp]{cor:sub}ના દાવા અનુસાર સુવ્યાખ્યાયિત
છે. જોકે એ ઉપસિદ્ધાંત માટેની અગાઉની સાબિતીમાં
\olref[alp]{thm:sub} અંગે નોંધાયેલો ગાળો હજી ખુલ્લો છે.

આ વ્યાખ્યાથી આપણી દલીલો ઘણી સરળ બને છે. દાખલા તરીકે:
\begin{align}
  \Subst{\lambd[x][x]}{y}{x} & \ollabel{eq:1}\\
  &= \Subst{\lambd[z][z]}{y}{x} \ollabel{eq:2}\\
  &= \lambd[z][\Subst{z}{y}{x}] \\
  &= \lambd[z][z]
\end{align}

આ સમીકરણો $\alpha$-સામ્ય વર્ગો વિશે છે: $\lambd[x][x]$
અને $\lambd[z][z]$ એક જ વર્ગમાં હોવાથી
\olref{eq:1}માંથી \olref{eq:2}માં જઈ શકાય છે. અગાઉ
સુધારેલી પદ-સ્તરની સ્થાનાપન્ન-વ્યાખ્યા મુજબ
$\Subst{\lambd[x][x]}{y}{x}$ વ્યાખ્યાયિત છે અને
તેનું મૂલ્ય $\lambd[x][x]$ છે; વર્ગ-સ્તરે તે
$\lambd[z][z]$ના વર્ગ બરાબર છે.
\sourcecorrection{OLLAM-028}{સ્થિર અંગ્રેજી મૂળ
\olref{eq:1}ના પદ-સ્તરના સ્થાનાપન્નને અવ્યાખ્યાયિત
કહે છે, પરંતુ અગાઉ સુધારેલા બદ્ધ-ચલ નિયમ પ્રમાણે તે
વ્યાખ્યાયિત અને યથાવત છે. અહીં પદની શાબ્દિક સમાનતા
અને વર્ગની સમાનતા વચ્ચેનો ભેદ બતાવ્યો છે. મૂળ
પ્રદર્શનની પહેલી પંક્તિમાં અંતે લટકતું વધારાનું
સમાનચિહ્ન પણ દૂર કર્યું છે.}

આ જ કારણથી આગળ સ્થાનાપન્ન માટે જરૂરી શરતો
સંતોષતા પ્રતિનિધિઓ પસંદ કર્યાં છે એમ માનીશું. દાખલા તરીકે,
$\Subst{\lambd[x][N]}{R}{y}$ દેખાય ત્યારે
$\lambd[x][N]$ એવો પ્રતિનિધિ પસંદ કરાયો છે કે
$x \neq y$ અને $x \notin FV(R)$.

$\lambd[x][x]$ને ``વર્ગ'' કહેવું થોડું વિચિત્ર લાગે છે.
અગાઉનાં $\lambda$-પદોથી ભેદ રાખવા આ વર્ગોને હવે પછી
$\Lambda$-પદો કહીશું; ભાગના બાકીના વિભાગોમાં ટૂંકમાં
``પદો'' પણ કહીશું.

\begin{editorial}
  પદોને આપણે હજી વિદાય આપી શકતા નથી: $\Lambda$-પદોની
  આખી વ્યાખ્યા $\lambda$-પદો પર આધારિત છે. વધુમાં,
  $\Lambda$-પદો પર વિધેયો વ્યાખ્યાયિત કરવાની સ્વતંત્ર
  રીત હજી આપી નથી. તેથી એવા વિધેયો પહેલાં
  $\lambda$-પદો પર વ્યાખ્યાયિત કરીને પછી, સ્થાનાપન્નની
  જેમ, $\Lambda$-પદો પર ``પ્રક્ષેપિત'' કરવા પડે.
  જોકે વાચક આવાં $\Lambda$-પદો પરના વિધેયો કેવી રીતે વ્યાખ્યાયિત થાય તે
  સાહજિક રીતે સમજી શકશે એમ માનીએ છીએ.
\end{editorial}
% END OLP-0364
\endgroup


\begingroup
\makeatletter\def\ol@sectioncs{section}\makeatother

% BEGIN OLP-0365
\olfileid{lam}{syn}{bet}
\olsection{$\beta$-લઘુકરણ}
\phantomsection\label{guunit:OLP-0365}

\sourcecorrection{OLLAM-029}{સ્થિર અંગ્રેજી મૂળનો
અધ્યાય-ટિપ્પણી અને \texttt{\string\olfileid} આ
વિભાગને ભૂલથી પરિચય અધ્યાયમાં મૂકે છે. તે વાક્યરચના
અધ્યાયમાંથી આયાત થાય છે; તેથી અહીં અધ્યાય-ઓળખ
\texttt{syn} કરી છે.}

$(\lambd[m][(\lambd[y][y]) m])$ જોઈએ ત્યારે તેનો
$\lambd[m][m]$ સાથે સંબંધ હોવો જોઈએ એવું સહજ લાગે છે:
બીજું પદ પહેલાને ``સરળ'' કરવાથી મળતું હોય એમ જણાય છે.
\emph{$\beta$-લઘુકરણ} આ વિચારને ઔપચારિક સ્વરૂપ આપે છે.

\begin{defn}[$\beta$-સંકોચન, $\bredone$] \ollabel{defn:betacontr}
  \emph{$\beta$-સંકોચન} ($\bredone$) પદો પરનો સૌથી
  નાનો રચના-સુસંગત સંબંધ છે જે નીચેની શરત સંતોષે છે:
  \[
    (\lambd[x][N])Q \bredone \Subst{N}{Q}{x} 
  \]
  $P \bredone Q$ હોય તો $P$નું $Q$માં
  \emph{$\beta$-સંકોચન} થાય છે એમ કહીએ. $(\lambd[x][N])Q$
  સ્વરૂપનું પદ \emph{સંકોચ્ય પદ (રેડેક્સ)} કહેવાય છે.
\end{defn}

\begin{prob} \ollabel{prob:def}
  \olref[lam][syn][alp]{defn:aconvone}માં બદ્ધ ચલનું નામ
  બદલવા માટે આપી હતી તેવી $\beta$-સંકોચનની સમકક્ષ
  આગમનાત્મક વ્યાખ્યાઓ વિગતે લખો.
\end{prob}
  
\begin{defn}[$\beta$-લઘુકરણ, $\bred$] \ollabel{defn:betared}
  \emph{$\beta$-લઘુકરણ} ($\bred$) પદો પરનો
  $\bredone$ને સમાવતો સૌથી નાનો સ્વવાચક અને
  પરંપરિત સંબંધ છે. $P \bred Q$ હોય તો $P$નું
  $Q$માં $\beta$-લઘુકરણ થાય છે એમ કહીએ.
\end{defn}

સંદર્ભ સ્પષ્ટ હોય ત્યારે $\bredone$ને બદલે $\redone$
અને $\bred$ને બદલે $\red$ લખીશું.

અનૌપચારિક રીતે, $M \bred N$ ત્યારે અને માત્ર ત્યારે જ
જ્યારે $\beta$-સંકોચનનાં શૂન્ય અથવા વધુ પગલાંથી
$M$ને $N$માં બદલી શકાય.

\begin{defn}[$\beta$-સામાન્ય]
જે પદનું આગળ $\beta$-સંકોચન થઈ ન શકે તેને
\emph{$\beta$-સામાન્ય} કહે છે.
\end{defn}

$M \bred N$ અને $N$ $\beta$-સામાન્ય હોય તો $N$ને
$M$નું \emph{સામાન્ય સ્વરૂપ} કહીએ. કોઈ પદનું સામાન્ય
સ્વરૂપ અનન્ય છે કે નહીં એવો પ્રશ્ન થાય; જવાબ હા છે,
જેમ આગળ જોઈશું.

કેટલાંક ઉદાહરણો જોઈએ.
\begin{enumerate}
\item આપણને મળે છે:
\begin{align*}
(\lambd[x][xxy]) \lambd[z][z] & \redone (\lambd[z][z])(\lambd[z][z]) y \\
& \redone (\lambd[z][z]) y \\
& \redone y
\end{align*}
\item પદને ``સરળ'' કરવાથી તે વધુ જટિલ પણ બની શકે છે:
\begin{align*}
(\lambd[x][xxy])(\lambd[x][xxy]) & \redone (\lambd[x][xxy])(\lambd[x][xxy])y \\
& \redone (\lambd[x][xxy])(\lambd[x][xxy])yy \\
& \redone \dots
\end{align*}
\item સંકોચન પછી પદ યથાવત પણ રહી શકે છે:
\[
(\lambd[x][xx])(\lambd[x][xx]) \redone (\lambd[x][xx])(\lambd[x][xx])
\]
\item વળી, કેટલાંક પદોનું એક કરતાં વધુ રીતે લઘુકરણ
  થઈ શકે છે. દાખલા તરીકે,
\[
(\lambd[x][(\lambd[y][yx]) z]) v \redone (\lambd[y][yv]) z
\]
બહારના અનુપ્રયોગનું સંકોચન કરીને મળે છે; અને
\[
(\lambd[x][(\lambd[y][yx]) z]) v \redone (\lambd[x][zx]) v
\]
અંદરના અનુપ્રયોગનું સંકોચન કરીને મળે છે. જોકે આ
કિસ્સામાં બંને પદો આગળ જઈને એ જ પદ $zv$માં લઘુકૃત
થાય છે.
\end{enumerate}

છેલ્લા ઉદાહરણમાં મળેલું સમાન અંતિમ પરિણામ સંયોગ નથી;
તે ``ચર્ચ--રોસર ગુણધર્મ'' નામનો લૅમ્બડા કલનનો ઊંડો
અને મહત્વનો ગુણધર્મ દર્શાવે છે.

\begin{digress}
  સામાન્ય રીતે પદનું $\beta$-લઘુકરણ કરવાની એક કરતાં
  વધુ રીતો હોય છે. તેથી લઘુકરણની ઘણી વ્યૂહરચનાઓ
  ઘડાઈ છે; તેમાં સૌથી સામાન્ય \emph{સહજ વ્યૂહરચના} છે.
  આ વ્યૂહરચના હંમેશા સૌથી ડાબા સંકોચ્ય પદનું સંકોચન કરે
  છે; અહીં સંકોચ્ય પદનું સ્થાન આખા પદમાં તેની શરૂઆતથી
  નક્કી થાય છે. આ વ્યૂહરચનાનો ઉપયોગી ગુણધર્મ એ છે કે
  કોઈ વ્યૂહરચનાથી પદનું સામાન્ય સ્વરૂપ મળી શકે ત્યારે
  અને માત્ર ત્યારે જ સહજ વ્યૂહરચનાથી પણ મળી શકે.
  આગળ જુદું ન કહીએ ત્યાં સુધી સહજ વ્યૂહરચના વાપરીશું.
\end{digress}

\begin{defn}[$\beta$-સામ્ય, $\equal$]
  \emph{$\beta$-સામ્ય} ($\equal$) નીચેના નિયમો વડે
  આગમનથી વ્યાખ્યાયિત સંબંધ છે:
  \begin{enumerate}
  \item $M \equal M$.
  \item જો $M \equal N$, તો $N \equal M$.
  \item જો $M \equal N$, $N \equal O$, તો $M \equal O$.
  \item જો $M \equal N$, તો $PM \equal PN$.
  \item જો $M \equal N$, તો $MQ \equal NQ$.
  \item જો $M \equal N$, તો $\lambd[x][M] \equal \lambd[x][N]$.
  \item $(\lambd[x][N])Q \equal \Subst{N}{Q}{x}$.
  \end{enumerate}
\end{defn}

પહેલા ત્રણ નિયમો આ સંબંધને સામ્ય સંબંધ બનાવે છે;
પછીના ત્રણ તેને રચના-સુસંગત બનાવે છે; અને છેલ્લો નિયમ
તેમાં $\beta$-સંકોચનનો સમાવેશ કરે છે.

અનૌપચારિક રીતે, બે પદો $\beta$-સામ્યવાળાં ત્યારે અને
માત્ર ત્યારે જ કહેવાય જ્યારે એકને $\beta$-સંકોચનનાં
અથવા $\beta$-સંકોચનના ``વિપરીત'' પગલાંનાં શૂન્ય કે
વધુ પગલાંથી બીજામાં બદલી શકાય. વિપરીત $\beta$-સંકોચન
એ રીતે વ્યાખ્યાયિત છે કે $M$નું વિપરીત $\beta$-સંકોચન કરીને
$N$ મળે ત્યારે અને માત્ર ત્યારે જ $N$નું $\beta$-સંકોચન
કરીને $M$ મળે.

ઉપરના નિયમો ઉપરાંત આપણે આ સંબંધમાં વધુ નિયમો ઉમેરીશું.
વિસ્તારેલા સામ્ય સંબંધને $\equal[X]$થી દર્શાવીશું,
જ્યાં $X$ ઉમેરેલો નિયમ છે.
% END OLP-0365
\endgroup


\begingroup
\makeatletter\def\ol@sectioncs{section}\makeatother

% BEGIN OLP-0366
\olfileid{lam}{syn}{eta}
\olsection{$\eta$-રૂપાંતરણ}
\phantomsection\label{guunit:OLP-0366}


$\lambda$-પદો પર બીજો પણ એક સંબંધ છે. \olref[fv]{sec}માં
આપણે $\lambd[x][(fx)]$નું ઉદાહરણ જોયું હતું: તે આર્ગ્યુમેન્ટ
સ્વીકારીને તેના પર~$f$ લાગુ કરે છે. બીજા શબ્દોમાં, તે~$f$
જેવું જ વિધેય છે: $\lambd[x][(fx)]N$ અને~$fN$ બંનેનું
લઘુકરણ~$fN$માં થાય છે. આ વિચાર વ્યક્ત કરવા
$\eta$-લઘુકરણ (અને $\eta$-વિસ્તરણ) વાપરીએ છીએ.

\begin{defn}[$\eta$-સંકોચન, $\eredone$]
  \ollabel{defn:beredone}
  \emph{$\eta$-સંકોચન} ($\eredone$) પદો પરનો સૌથી નાનો
  રચના-સુસંગત સંબંધ છે જે નીચેની શરત સંતોષે છે:
  \[
    \lambd[x][M x] \eredone M \text{ જો } x \notin FV(M)
  \]
\end{defn}

\begin{defn}[$\beta\eta$-લઘુકરણ, $\bered$] \ollabel{defn:bered}
  \emph{$\beta\eta$-લઘુકરણ} ($\bered$) પદો પરનો
  $\bredone$ અને $\eredone$ને સમાવતો સૌથી નાનો
  સ્વવાચક અને પરંપરિત સંબંધ છે. એટલે સ્વવાચકતા અને
  પરંપરિતતાના નિયમો ઉપરાંત નીચેના બે નિયમો છે:
  \begin{enumerate}
  \item જો $M \bredone N$, તો $M \bered N$. \ollabel{defn:bered3}
  \item જો $M \eredone N$, તો $M \bered N$. \ollabel{defn:bered4}
  \end{enumerate}
    
\end{defn}

\begin{defn}
  સામ્ય સંબંધ~$\equal$માં $\eta$-રૂપાંતરણનો નિયમ ઉમેરીએ:
  \[
  \lambd[x][f x] \equal f
  \]
  આ રીતે વિસ્તરેલા સંબંધને $\equal[\eta]$થી દર્શાવીએ.
\end{defn}

$\eta$-સામ્ય મહત્વનું છે, કારણ કે તે લૅમ્બડા પદોની
વિસ્તારાત્મકતા સાથે સંબંધિત છે:

\begin{defn}[વિસ્તારાત્મકતા]
  સામ્ય સંબંધ~$\equal$માં (\ext) નિયમ ઉમેરીએ:
  \begin{center}
  જો $Mx \equal Nx$, તો $M \equal N$; શરત છે કે $x \notin FV(MN)$.
  \end{center}
  આ રીતે વિસ્તરેલા સંબંધને $\equal[\ext]$થી દર્શાવીએ.
\end{defn}

આ નિયમનો સાર એ છે કે પદોને વિધેયો માનીએ ત્યારે એક જ
મનસ્વી આર્ગ્યુમેન્ટ પર તેમનું વર્તન સમાન હોય તો તેમને
સમાન ગણવાં.

હવે સાબિત કરીએ કે $\eta$-નિયમથી બરાબર વિસ્તારાત્મકતા
મળે છે, તેનાથી વધુ કંઈ નહીં.

\begin{thm}
  $M \equal[\ext] N$ ત્યારે અને માત્ર ત્યારે જ જ્યારે
  $M \equal[\eta] N$.
\end{thm}

\begin{proof}
  પ્રદર્શિત $\eta$-નિયમમાં $f$ ચલ છે. પરંતુ
  $x \notin FV(M)$ હોય અને~$f$ને $M$ તથા~$x$થી
  તાજો લઈએ તો રચના-સુસંગતતાથી
  $\lambd[f][\lambd[x][f x]] \equal[\eta]
  \lambd[f][f]$. બંને બાજુ~$M$ લાગુ કરી
  $\beta$-સામ્ય વાપરતાં $\lambd[x][Mx]
  \equal[\eta] M$ મળે છે. આથી આગળ સામાન્ય પદો માટે
  $\eta$-નિયમ વાપરી શકીએ.
  \sourcecorrection{OLLAM-030}{સ્થિર અંગ્રેજી મૂળના
  પ્રદર્શિત $\eta$-નિયમમાં $f$ ચલ છે, પરંતુ સાબિતી
  ગમે તે પદ~$M$ માટે એ નિયમ વાપરે છે. અહીં તાજો
  $f$, રચના-સુસંગતતા અને $\beta$-સામ્યથી સામાન્ય
  નિયમ મેળવવાનું ગાયબ પગલું ઉમેર્યું છે.}
  પહેલાં બતાવીએ કે $\equal[\eta]$ વિસ્તારાત્મકતાના
  નિયમ હેઠળ બંધ છે. તેથી~\ext{} નિયમ
  $\equal[\eta]$માં નવી જોડી ઉમેરતો નથી: નિયમ-ઉપજ
  પર આગમનથી $\equal[\ext]$નો દરેક સંબંધ
  $\equal[\eta]$માં આવે છે. તેથી
  $M \equal[\ext] N$ હોય તો $M \equal[\eta] N$.

  \ext{} નિયમ હેઠળની બંધતા માટે, એ નિયમ લાગુ કરતાં
  $Mx \equal[\ext] Nx$ એ પૂર્વધારણા હોય છે અને
  $x \notin FV(MN)$ હોય છે. આગમનની ધારણાથી
  $Mx \equal[\eta] Nx$. રચના-સુસંગતતાથી
  $\lambd[x][Mx] \equal[\eta] \lambd[x][Nx]$.
  $x$ બંને પદોમાં મુક્ત ન હોવાથી, બંને બાજુ
  $\eta$-નિયમ લગાડી પરંપરિતતાથી
  $M \equal[\eta] N$ મળે છે.
  \sourcecorrection{OLLAM-031}{સ્થિર અંગ્રેજી મૂળ
  \ext{}-નિયમની પૂર્વધારણાને તરત
  $Mx \equal[\eta] Nx$ કહે છે, જ્યારે તે શરૂઆતમાં
  $Mx \equal[\ext] Nx$ છે; પહેલાંના નિષ્કર્ષ પર
  આગમનની ધારણા લગાવવી પડે. અહીં એ પગલું અને
  $x \notin FV(MN)$ની શરત સ્પષ્ટ કર્યાં છે.
  મૂળના બીજા દિશાના તર્કમાં $\equal[ext]$ લખાયું
  છે; વ્યાખ્યાયિત સંકેત $\equal[\ext]$ છે.}

  ઊલટું, $\eta$-નિયમ $\equal[\ext]$માં સમાયેલો છે.
  જો $x \notin FV(M)$ હોય તો $\beta$-નિયમથી
  $(\lambd[x][Mx])x \equal[\ext] Mx$. વળી,
  $x \notin FV((\lambd[x][Mx])M)$; તેથી
  \ext{} નિયમથી $\lambd[x][Mx] \equal[\ext] M$.
  આ રીતે
  $\equal[\eta]$ના નિયમો પણ $\equal[\ext]$માં
  સમાયેલા છે; એટલે બંને સંબંધો સમાન છે.
\end{proof}
% END OLP-0366
\endgroup


\OLEndChapterHook


\begin{editorial}
આગળનો સંબંધિત અંશ મૂળના સક્રિય આયાતક્રમમાં નથી. તેને અહીં સ્પષ્ટ વૈકલ્પિક સંદર્ભમાં જાળવ્યો છે.
\end{editorial}
\begingroup
\makeatletter\def\ol@sectioncs{section}\makeatother

% BEGIN OLP-0651
\olfileid{lam}{int}{con}
\olsection{રૂપાંતરણ અને લઘુકરણ}
\phantomsection\label{guunit:OLP-0651}


લૅમ્બડા કલનમાં રૂપાંતરણ કે લઘુકરણના ત્રણ
પ્રકાર છે. આગળ જોઈશું તેમ, રૂપાંતરણ અને
લઘુકરણ વચ્ચેનો તફાવત એ છે કે
``બદલાવ'' એકદિશ છે કે દ્વિદિશ.

$\alpha$-રૂપાંતરણ ચલનું નામ બદલવા વિશે છે.
$\lambda x. x$ અને $\lambda y.y$ એ એક
જ વિધેય, એટલે ઓળખ-વિધેય, દર્શાવે છે એમ
માનવું સ્વાભાવિક છે. આ વિચારને નીચે મુજબ
ઔપચારિક સ્વરૂપ આપવાનું છે:
\sourcecorrection{OLMD-188}{સ્થિર મૂળ આ
દ્વિબિંદુ પછી તરત સમાપ્ત થાય છે; જણાવેલી
ઔપચારિક વ્યાખ્યા આ વિભાગમાં આપેલી નથી.
અલગ $\alpha$-રૂપાંતરણ વિભાગમાં વિગત છે,
પણ તેને અહીં લખાયેલું ગણ્યું નથી.}
% END OLP-0651
\endgroup
% END OLP-0356
\endgroup


\begingroup
\makeatletter\def\ol@sectioncs{section}\makeatother

% BEGIN OLP-0367
\olchapter{lam}{cr}{ચર્ચ--રોસર ગુણધર્મ}
\olchapterid{lam}{cr}
\phantomsection\label{guunit:OLP-0367}


\begingroup
\makeatletter\def\ol@sectioncs{section}\makeatother

% BEGIN OLP-0368
\olfileid{lam}{cr}{dap}

\olsection{વ્યાખ્યા અને ગુણધર્મો}
\phantomsection\label{guunit:OLP-0368}


આ અધ્યાયમાં ચર્ચ--રોસર ગુણધર્મનો ખ્યાલ અને તેના
કેટલાક સામાન્ય ગુણધર્મો રજૂ કરીએ છીએ.

\begin{defn}[ચર્ચ--રોસર ગુણધર્મ, CR]
  પદો પરનો સંબંધ~$\xredone$ \emph{ચર્ચ--રોસર ગુણધર્મ}
  ધરાવે છે ત્યારે અને માત્ર ત્યારે જ જ્યારે
  $M \xredone P$ અને $M \xredone Q$ હોય ત્યારે
  એવું~$N$ મળે કે $P \xredone N$ અને
  $Q \xredone N$.
\end{defn}

લૅમ્બડા કલનને સંગણનાના નમૂના તરીકે જોઈ શકાય: તેમાં
સામાન્ય સ્વરૂપનાં પદો ``મૂલ્યો'' છે અને પદો પરનો
લઘુકરણ-સંબંધ ``ગણતરીના નિયમો'' આપે છે. ચર્ચ--રોસર
ગુણધર્મ કહે છે કે ગણતરીમાં આગળ વધવાની એક કરતાં વધુ
રીતો હોય તોપણ અભિવ્યક્તિનું અંતિમ મૂલ્ય એક જ રહે છે.

પ્રાથમિક બીજગણિતનું ઉદાહરણ લઈએ. $4 \times (1+2) + 3$
ગણવાની એકથી વધુ રીતો છે. પહેલા $1+2$ને~$3$ કરીએ તો
$4 \times 3+3$ મળે; અથવા વિતરણ-નિયમથી પહેલા
$4 \times (1+2)$ને વિસ્તારી $4 \times 1+4
\times 2+3$ મળે. જોકે બંનેને આગળ ગણતાં $12+3$ જ
મળે છે.

$\xredone$ને $\beta$-લઘુકરણ માનીએ તો ચર્ચ--રોસર
ગુણધર્મનું પરિણામ છે કે પદને સામાન્ય સ્વરૂપ મળે તો તે
અનન્ય છે. માનો કે~$M$નું $P$ અને $Q$ બંનેમાં લઘુકરણ
થાય છે અને બંને સામાન્ય સ્વરૂપો છે. ચર્ચ--રોસર
ગુણધર્મથી એવું~$N$ મળે કે $P$ અને $Q$ બંનેનું
તેમાં લઘુકરણ થાય. $P$ અને $Q$ સામાન્ય સ્વરૂપમાં
હોવાથી $P$ અને $Q$નું $N$માં લઘુકરણ ફક્ત તુચ્છ,
એટલે શૂન્ય-પગલાનું જ
હોઈ શકે; તેથી $P$, $Q$ અને $N$ એક જ પદ છે. આથી
કોઈ પદના \emph{તે} સામાન્ય સ્વરૂપની વાત કરવી યોગ્ય છે.

લૅમ્બડા કલનને સંગણનાના નમૂના તરીકે જોતાં પદનું સામાન્ય
સ્વરૂપ તે પદથી શરૂ થતી ગણતરીનું ``અંતિમ પરિણામ''
ગણી શકાય. ઉપરના ઉપસિદ્ધાંત મુજબ આવું અંતિમ પરિણામ
હોય તો માત્ર એક જ હોય છે, જેમ કે $4 \times (1+2)+3$
ગણતાં અંતે માત્ર~$15$ મળે છે.

\begin{thm} \ollabel{thm:str}
  સંબંધ~$\xredone$ ચર્ચ--રોસર ગુણધર્મ ધરાવતો હોય
  અને $\xred$ એ $\xredone$ને સમાવતો સૌથી નાનો
  પરંપરિત સંબંધ હોય, તો $\xred$ પણ ચર્ચ--રોસર
  ગુણધર્મ ધરાવે છે.
\end{thm}

\begin{proof}
  માનો કે
  \begin{align*}
    M & \xredone P_1 \xredone \dots \xredone P_m \text{ અને}\\
    M & \xredone Q_1 \xredone \dots \xredone Q_n.
  \end{align*}
  ઊંચાઈ~$m + 1$ અને પહોળાઈ~$n + 1$ ધરાવતું
  પદોનું જાળક~$N$ રચીને પ્રમેય સાબિત કરીશું.
  $N_{i,j}$ એ $i$મી પંક્તિ અને $j$મા સ્તંભનું પદ છે.
  
  જાળક~$N$ એવું રચીએ કે $N_{i,j} \xredone N_{i+1,j}$
  અને $N_{i,j} \xredone N_{i,j+1}$. તેની વ્યાખ્યા:
  \begin{align*}
    N_{0,0} &= M \\
    N_{i,0} &= P_i && \text{જો } 1 \le i \le m \\
    N_{0,j} &= Q_j && \text{જો } 1 \le j \le n \\
  \intertext{અને બાકી કિસ્સામાં:}
    N_{i,j} &= R
  \end{align*}
  અહીં~$R$ એવું પદ છે કે $N_{i-1,j} \xredone R$ અને
  $N_{i,j-1} \xredone R$. $\xredone$ના ચર્ચ--રોસર
  ગુણધર્મથી આવું પદ હંમેશા મળે છે.

  હવે $N_{m,0} \xredone \dots \xredone N_{m,n}$
  અને $N_{0,n} \xredone \dots \xredone N_{m,n}$.
  અહીં $N_{m,0}$ એ $P_m$ અને $N_{0,n}$ એ $Q_n$ છે.
  $\xred$ની વ્યાખ્યાથી દાવો મળે છે.
  \sourcecorrection{OLLAM-032}{સ્થિર અંગ્રેજી મૂળમાં
  સાબિતીના અંતે $N_{m,0}$ને $P$ અને $N_{0,n}$ને
  $Q$ કહે છે, પરંતુ આ સાબિતીમાં ફક્ત~$P_i$ અને
  $Q_j$ વ્યાખ્યાયિત છે. અહીં ચોક્કસ અંતિમ પદો
  $P_m$ અને $Q_n$ લખ્યાં છે.}
\end{proof}
% END OLP-0368
\endgroup

\begingroup
\makeatletter\def\ol@sectioncs{section}\makeatother

% BEGIN OLP-0369
\olfileid{lam}{cr}{pb}

\olsection{સમાંતર $\beta$-લઘુકરણ}
\phantomsection\label{guunit:OLP-0369}


અહીં \emph{સમાંતર $\beta$-લઘુકરણ}નો ખ્યાલ રજૂ કરીને
તે ચર્ચ--રોસર ગુણધર્મ ધરાવે છે એમ સાબિત કરીશું.

\begin{defn}[સમાંતર $\beta$-લઘુકરણ, $\bredpar$] \ollabel{defn:bredpar}
  પદોનું સમાંતર લઘુકરણ ($\bredpar$) આગમનથી આ રીતે
  વ્યાખ્યાયિત થાય છે:
  \begin{enumerate}
    \item \ollabel{defn:bredpar1} $x \bredpar x$.
    \item \ollabel{defn:bredpar2} જો $N \bredpar N'$, તો
      $\lambd[x][N] \bredpar \lambd[x][N']$.
      \sourcecorrection{OLLAM-033}{સ્થિર અંગ્રેજી મૂળના
      બીજા નિયમની પૂર્વધારણા $N \xrightarrow{\beta} N'$
      છે, જ્યારે પછીની સ્વવાચકતા અને આગમનાત્મક
      સાબિતીઓ આ જ સંબંધની પૂર્વધારણા
      $N \bredpar N'$ વાપરે છે. અહીં તે સુધાર્યું છે.}
    \item \ollabel{defn:bredpar3} જો $P \bredpar P'$ અને
      $Q \bredpar Q'$, તો $PQ \bredpar P'Q'$.
    \item \ollabel{defn:bredpar4} જો $N \bredpar N'$ અને
      $Q \bredpar Q'$, તો
      $(\lambd[x][N])Q \bredpar \Subst{N'}{Q'}{x}$.
  \end{enumerate}
\end{defn}

સમાંતર $\beta$-લઘુકરણથી પદમાંનાં ગમે તેટલાં
સંકોચ્ય પદો એક જ પગલામાં સંકોચી શકાય છે. તે સામાન્ય
$\beta$-લઘુકરણથી જુદું છે: અહીં
મૂળ પદમાં પહેલેથી હાજર સંકોચ્ય પદો જ સંકોચી શકાય;
સમાંતર $\beta$-લઘુકરણથી નવાં ઊભાં થતાં સંકોચ્ય પદો એ જ
પગલામાં સંકોચી શકાતાં નથી. દાખલા તરીકે,
$(\lambd[f][fx])(\lambd[y][y])$નું એક પગલામાં
સમાંતર $\beta$-લઘુકરણ તેના પોતાના પદમાં અથવા
$(\lambd[y][y])x$માં થઈ શકે છે, પણ સીધું~$x$માં
નહીં. જોકે સામાન્ય $\beta$-લઘુકરણથી તે~$x$માં
જાય છે. નવું સંકોચ્ય પદ પહેલા સમાંતર $\beta$-લઘુકરણના
પગલા પછી જ ઊભું થાય છે; બીજા સમાંતર $\beta$-લઘુકરણના
પગલાથી~$x$ મળે છે.

\begin{thm}\ollabel{thm:refl}
  $M \bredpar M$.
\end{thm}

\begin{proof}
  અભ્યાસપ્રશ્ન.
\end{proof}

\begin{prob}
  \olref[lam][cr][pb]{thm:refl} સાબિત કરો.
\end{prob}

\begin{defn}[$\beta$-પૂર્ણ વિકાસ]\ollabel{defn:bcd}
  $M$નો \emph{$\beta$-પૂર્ણ વિકાસ}~$\bcd{M}$ આ રીતે
  આગમનથી વ્યાખ્યાયિત થાય છે:
  \begin{align}
    \bcd{x} &= x \ollabel{defn:bcd1} \\
    \bcd{(\lambd[x][N])} &= \lambd[x][\bcd{N}] \ollabel{defn:bcd2}\\
    \bcd{(PQ)} &= \bcd{P}\bcd{Q} && \text{જો $P$ એ $\lambd$-અમૂર્તીકરણ ન હોય}
    \ollabel{defn:bcd3} \\
    \bcd{((\lambd[x][N])Q)} &= \Subst{\bcd{N}}{\bcd{Q}}{x} \ollabel{defn:bcd4}
  \end{align}
\end{defn}

પદનો $\beta$-પૂર્ણ વિકાસ નામ મુજબ ``પૂર્ણ સમાંતર
લઘુકરણ'' છે. સામાન્ય સમાંતર $\beta$-લઘુકરણમાં કોઈ
સંકોચ્ય પદ ન સંકોચવાનું પસંદ કરી શકાય છે; પૂર્ણ વિકાસમાં
બધાંનું સંકોચન કરવું જ પડે. તેથી
$(\lambd[f][fx])(\lambd[y][y])$નો પૂર્ણ વિકાસ
$(\lambd[y][y])x$ છે, મૂળ પદ પોતે નહીં.

\begin{editorial}
  આ વ્યાખ્યામાં એક અડચણ છે: આપણે હજી
  ($\lambd$-)પદો પર વિધેયોનું પુનરાવર્તી વ્યાખ્યાયન
  કેવી રીતે કરવું તે રજૂ કર્યું નથી. આગળ સુધારીશું.
\end{editorial}

\begin{lem}\ollabel{lem:comp}
  $M \bredpar M'$ અને $R \bredpar R'$ હોય તો
  $\Subst{M}{R}{y} \bredpar
  \Subst{M'}{R'}{y}$.
\end{lem}

\begin{proof}
  $M \bredpar M'$ની નિષ્પત્તિ પર આગમનથી.
  \begin{enumerate}
    \item છેલ્લો નિયમ \olref{defn:bredpar1} હોય: અભ્યાસપ્રશ્ન.
    \item છેલ્લો નિયમ \olref{defn:bredpar2} હોય:
      $M$ એ $\lambd[x][N]$ અને $M'$ એ $\lambd[x][N']$,
      જ્યાં $N \bredpar N'$. બતાવવું છે કે
      $\Subst{(\lambd[x][N])}{R}{y} \bredpar
      \Subst{(\lambd[x][N'])}{R'}{y}$; એટલે કે,
      યોગ્ય તાજા બાંધનાર~$x$ પસંદ કર્યા પછી
      $\lambd[x][\Subst{N}{R}{y}] \bredpar
      \lambd[x][\Subst{N'}{R'}{y}]$.
      આગમનની ધારણા અને \olref{defn:bredpar2}થી આ મળે છે.
      \sourcecorrection{OLLAM-034}{સ્થિર અંગ્રેજી મૂળ
      આ કિસ્સાના બીજા અમૂર્તીકરણમાં
      $\Subst{N'}{R}{y}$ લખે છે, જોકે દાવામાં બીજો
      પ્રતિસ્થાપક~$R'$ છે. અહીં~$R'$ કર્યું છે;
      બાંધનાર તાજો પસંદ કરવાની જરૂર પણ સ્પષ્ટ કરી છે.}
    \item છેલ્લો નિયમ \olref{defn:bredpar3} હોય: અભ્યાસપ્રશ્ન.
    \item છેલ્લો નિયમ \olref{defn:bredpar4} હોય:
      $M$ એ $(\lambd[x][N])Q$ અને
      $M'$ એ $\Subst{N'}{Q'}{x}$. બતાવવું છે કે
      $\Subst{((\lambd[x][N])Q)}{R}{y} \bredpar
      \Subst{\Subst{N'}{Q'}{x}}{R'}{y}$.
      બાંધનાર~$x$ને $y$, $R$ અને~$R'$થી તાજો
      પસંદ કર્યા પછી દાવાને નીચેના સ્વરૂપમાં લાવવાનો રહે છે:
      $(\lambd[x][\Subst{N}{R}{y}])\Subst{Q}{R}{y} \bredpar
      \Subst{\Subst{N'}{R'}{y}}{\Subst{Q'}{R'}{y}}{x}$.
      અહીં જમણી બાજુ માટે સ્થાનાપન્ન-સંયોજનનો
      અલગ લેમ્મો જોઈએ. તે સાબિત થાય તો, આગમનની ધારણા અને
      \olref{defn:bredpar4}થી દાવો મળે છે.
      \sourcecorrection{OLLAM-035}{સ્થિર અંગ્રેજી મૂળ
      જમણી બાજુનાં બે ક્રમિક સ્થાનાપન્નોનો ક્રમ
      કોઈ દલીલ વિના ફેરવે છે. આ માટે~$x\neq y$
      અને~$x$ પ્રતિસ્થાપકમાં મુક્ત ન હોય તેવી
      તાજગી-શરતો સાથે સ્થાનાપન્ન-સંયોજનનો લેમ્મો
      સાબિત કરવો પડે. અહીં તે અનુપસ્થિત સહાયક
      લેમ્મો સ્પષ્ટ કર્યો છે; તેથી સ્થિર મૂળની આ
      સાબિતી આ બિંદુએ અધૂરી છે.}
  \end{enumerate}
\end{proof}

\begin{prob}
  \olref[lam][cr][pb]{lem:comp}ની સાબિતી પૂર્ણ કરો.
\end{prob}

\begin{lem}\ollabel{lem:cont}
  $M \bredpar M'$ હોય તો $M' \bredpar \bcd{M}$.
\end{lem}
\begin{proof}
  $M \bredpar M'$ની નિષ્પત્તિ પર આગમનથી.
  \begin{enumerate}
    \item છેલ્લો નિયમ \olref{defn:bredpar1} હોય: અભ્યાસપ્રશ્ન.
    \item છેલ્લો નિયમ \olref{defn:bredpar2} હોય:
      $M$ એ $\lambd[x][N]$ અને $M'$ એ $\lambd[x][N']$,
      જ્યાં $N \bredpar N'$. બતાવવું છે કે
      $\lambd[x][N'] \bredpar
      \bcd{(\lambd[x][N])}$, એટલે
      \olref{defn:bcd2} મુજબ
      $\lambd[x][N'] \bredpar \lambd[x][\bcd{N}]$.
      આગમનની ધારણા અને \olref{defn:bredpar2}થી આ મળે છે.
    \item છેલ્લો નિયમ \olref{defn:bredpar3} હોય:
      $M$ એ $PQ$ અને $M'$ એ $P'Q'$,
      એવા કોઈક $P$, $Q$, $P'$ અને $Q'$ માટે, જ્યાં
      $P \bredpar P'$ તથા $Q \bredpar Q'$.
      આગમનની ધારણાથી $P' \bredpar \bcd{P}$
      અને $Q' \bredpar \bcd{Q}$ છે.
      \begin{enumerate}
        \item $P$ એ $\lambd[x][N]$ હોય તો,
          એવા કોઈક $x$ અને $N$ માટે,
          $P'$ એ $\lambd[x][N']$ એવા કોઈ~$N'$ માટે,
          જ્યાં $N \bredpar N'$. આગમનની ધારણાથી
          $N' \bredpar \bcd{N}$ અને
          $Q' \bredpar \bcd{Q}$. તેથી
          \olref{defn:bredpar4} મુજબ
          $(\lambd[x][N'])Q' \bredpar
          \Subst{\bcd{N}}{\bcd{Q}}{x}$.
        \item $P$ કોઈ $\lambd$-અમૂર્તીકરણ ન હોય તો
          \olref{defn:bredpar3} મુજબ
          $P'Q' \bredpar \bcd{P}\bcd{Q}$;
          \olref{defn:bcd3} મુજબ જમણી બાજુ
          $\bcd{PQ}$ છે.
      \end{enumerate}
    \item છેલ્લો નિયમ \olref{defn:bredpar4} હોય:
      $M$ એ $(\lambd[x][N])Q$ અને
      $M'$ એ $\Subst{N'}{Q'}{x}$,
      એવા કોઈક $x$, $N$, $Q$, $N'$ અને~$Q'$ માટે, જ્યાં
      $N \bredpar N'$ અને $Q \bredpar Q'$.
      આગમનની ધારણાથી $N' \bredpar \bcd{N}$
      અને $Q' \bredpar \bcd{Q}$. તેથી
      \olref{lem:comp} મુજબ
      $\Subst{N'}{Q'}{x} \bredpar
      \Subst{\bcd{N}}{\bcd{Q}}{x}$;
      જમણી બાજુ બરાબર
      $\bcd{((\lambd[x][N])Q)}$ છે. અહીં
      \olref{lem:comp}ની અગાઉ નોંધેલી સાબિતી-ખામીનો
      આધાર હજી ખુલ્લો છે.
  \end{enumerate}
\end{proof}

\begin{prob}
  \olref[lam][cr][pb]{lem:cont}ની સાબિતી પૂર્ણ કરો.
\end{prob}

\begin{thm}\ollabel{thm:cr}
  $\bredpar$ ચર્ચ--રોસર ગુણધર્મ ધરાવે છે.
\end{thm}

\begin{proof}
  \olref{lem:cont}માંથી તરત મળે છે; એ લેમ્માની
  ઉપર નોંધેલી સાબિતી-ખામી અહીં પણ લાગુ પડે છે.
\end{proof}
% END OLP-0369
\endgroup

\begingroup
\makeatletter\def\ol@sectioncs{section}\makeatother

% BEGIN OLP-0370
\olfileid{lam}{cr}{b}

\olsection{$\beta$-લઘુકરણ}
\phantomsection\label{guunit:OLP-0370}


\begin{lem}\ollabel{lem:one-par}
  $M \bredone M'$ હોય તો $M \bredpar M'$.
\end{lem}
\begin{proof}
  $M \bredone M'$ની રચના પર આગમનથી. મૂળ સંકોચનનો
  કિસ્સામાં $M$ એ $(\lambd[x][N])Q$ અને
  $M'$ એ $\Subst{N}{Q}{x}$ છે, કોઈક
  $x$, $N$ અને~$Q$ માટે. \olref[pb]{thm:refl}
  મુજબ $N \bredpar N$ અને $Q \bredpar Q$;
  તેથી \olref[pb]{defn:bredpar}\olref[pb]{defn:bredpar4}
  મુજબ $(\lambd[x][N])Q \bredpar
  \Subst{N}{Q}{x}$. જો સંકોચન અમૂર્તીકરણની અંદર
  થયું હોય તો આગમનની ધારણા અને
  \olref[pb]{defn:bredpar2} વાપરીએ. અનુપ્રયોગની
  ડાબી અથવા જમણી બાજુ થયું હોય તો બદલાતી બાજુ
  માટે આગમનની ધારણા, બીજી માટે
  \olref[pb]{thm:refl} અને પછી
  \olref[pb]{defn:bredpar3} વાપરીએ.
  \sourcecorrection{OLLAM-036}{સ્થિર અંગ્રેજી મૂળની
  આ સાબિતી ફક્ત આખા પદના મૂળ સંકોચનનો કિસ્સો
  લે છે, જ્યારે $\bredone$ સૌથી નાનો
  રચના-સુસંગત સંબંધ છે અને સંકોચન અંદર પણ થઈ શકે.
  અહીં અમૂર્તીકરણ અને અનુપ્રયોગ હેઠળના બાકીના
  આગમનાત્મક કિસ્સા ઉમેર્યા છે.}
\end{proof}

\begin{lem}\ollabel{lem:par-red}
  $M \bredpar M'$ હોય તો $M \bred M'$.
\end{lem}

\begin{proof}
  $M \bredpar M'$ની નિષ્પત્તિ પર આગમનથી.
  \begin{enumerate}
  \item છેલ્લો નિયમ \olref[pb]{defn:bredpar1} હોય:
    $M$ અને $M'$ બંને ચલ~$x$ છે, તેથી $x \bred x$.
  \item છેલ્લો નિયમ \olref[pb]{defn:bredpar2} હોય:
    $M$ એ $\lambd[x][N]$ અને $M'$ એ
    $\lambd[x][N']$, કોઈક $x$, $N$, $N'$ માટે,
    જ્યાં $N \bredpar N'$. આગમનની ધારણાથી
    $N \bred N'$. $N \bred N'$ની જ $\bredone$-સંકોચન શ્રેણી
    અમૂર્તીકરણની અંદર
    લાગુ કરતાં $\lambd[x][N] \bred
    \lambd[x][N']$ મળે છે.
  \item છેલ્લો નિયમ \olref[pb]{defn:bredpar3} હોય:
    $M$ એ $PQ$ અને $M'$ એ $P'Q'$ છે,
    કોઈક $P$, $Q$, $P'$, $Q'$ માટે, જ્યાં
    $P \bredpar P'$ અને $Q \bredpar Q'$.
    આગમનની ધારણાથી $P \bred P'$ અને
    $Q \bred Q'$. પહેલાં $P \bred P'$,
    પછી $Q \bred Q'$ મુજબ લઘુકરણ કરીએ તો
    $PQ \bred P'Q'$ મળે છે.
  \item છેલ્લો નિયમ \olref[pb]{defn:bredpar4} હોય:
    $M$ એ $(\lambd[x][N])Q$ અને
    $M'$ એ $\Subst{N'}{Q'}{x}$, કોઈક
    $x$, $N$, $N'$, $Q$, $Q'$ માટે; અહીં
    $N \bredpar N'$ તથા $Q \bredpar Q'$.
    આગમનની ધારણાથી $Q \bred Q'$ અને
    $N \bred N'$. પહેલાં $N \bred N'$, પછી
    $Q \bred Q'$ મુજબ લઘુકરણ કરીને છેલ્લે
    $(\lambd[x][N'])Q'$નું
    $\Subst{N'}{Q'}{x}$માં સંકોચન કરીએ તો
    $(\lambd[x][N])Q \bred
    \Subst{N'}{Q'}{x}$ મળે છે.
    \sourcecorrection{OLLAM-037}{સ્થિર અંગ્રેજી મૂળ
    આ છેલ્લી કિસ્સાની પદ-યાદીમાં $N'$ની જગ્યાએ
    ફરી~$M'$ લખે છે. અહીં નિયમની પૂર્વધારણા પ્રમાણે
    $N'$ પુનઃસ્થાપિત કર્યું છે.}
  \end{enumerate}
\end{proof}

\begin{lem}\ollabel{lem:str}
  $\bred$ એ $\bredpar$ને સમાવતો સૌથી નાનો
  પરંપરિત સંબંધ છે.
\end{lem}

\begin{proof}
  $\xred$ એ $\bredpar$ને સમાવતો સૌથી નાનો
  પરંપરિત સંબંધ માનો.

  $\bred \subseteq \xred$: માનો કે $M \bred M'$.
  શૂન્ય પગલાંના કિસ્સામાં $M=M'$ છે, અને
  \olref[pb]{thm:refl}થી $M \bredpar M$, તેથી
  $M \xred M'$. બાકી $M \ident M_1
  \bredone \dots \bredone M_k \ident M'$.
  \olref{lem:one-par}થી $M \ident M_1
  \bredpar \dots \bredpar M_k \ident M'$.
  $\xred$ એ $\bredpar$ને સમાવતો પરંપરિત
  સંબંધ હોવાથી $M \xred M'$.
  \sourcecorrection{OLLAM-038}{સ્થિર અંગ્રેજી મૂળ
  $\bred$ના શૂન્ય-પગલાના સ્વવાચક કિસ્સાને
  એક-અથવા-વધુ પગલાંની શ્રેણીથી આવરી લેવાનો
  પ્રયાસ કરે છે. અહીં પહેલાં $\bredpar$ની
  સ્વવાચકતાથી તે કિસ્સો અલગ સાબિત કર્યો છે.}

  $\xred \subseteq \bred$: માનો કે $M \xred M'$,
  એટલે $M \ident M_1 \bredpar \dots
  \bredpar M_k \ident M'$. \olref{lem:par-red}થી
  $M \ident M_1 \bred \dots \bred
  M_k \ident M'$. $\bred$ પરંપરિત હોવાથી
  $M \bred M'$.
\end{proof}

\begin{thm}\ollabel{thm:cr}
  $\bred$ ચર્ચ--રોસર ગુણધર્મ ધરાવે છે.
\end{thm}

\begin{proof}
  \olref[dap]{thm:str}, \olref[pb]{thm:cr} અને
  \olref{lem:str}માંથી તરત મળે છે. જોકે
  \olref[pb]{thm:cr}ની અગાઉ નોંધેલી સાબિતી-ખામી
  હજી ખુલ્લી છે.
\end{proof}
% END OLP-0370
\endgroup

\begingroup
\makeatletter\def\ol@sectioncs{section}\makeatother

% BEGIN OLP-0371
\olfileid{lam}{cr}{pbe}

\olsection{સમાંતર $\beta\eta$-લઘુકરણ}
\phantomsection\label{guunit:OLP-0371}


આ વિભાગમાં $\beta\eta$-લઘુકરણને અનુરૂપ સમાંતર
$\beta\eta$-લઘુકરણ માટે ચર્ચ--રોસર ગુણધર્મ
સાબિત કરવાનો હેતુ છે.

\begin{defn}[સમાંતર $\beta\eta$-લઘુકરણ, $\beredpar$] \ollabel{defn:beredpar}
  પદો પરનું \emph{સમાંતર $\beta\eta$-લઘુકરણ}
  ($\beredpar$) આગમનથી આ રીતે વ્યાખ્યાયિત થાય છે:
  \begin{enumerate}
    \item \ollabel{defn:beredpar1} $x \beredpar x$.
    \item \ollabel{defn:beredpar2} જો $N \beredpar N'$,
      તો $\lambd[x][N] \beredpar \lambd[x][N']$.
      \sourcecorrection{OLLAM-039}{સ્થિર અંગ્રેજી મૂળ
      આ અમૂર્તીકરણ-નિયમમાં $N \xrightarrow{\beta} N'$
      લખે છે; પછીની સ્વવાચકતા અને આગમનાત્મક
      સાબિતીઓ માટે $N \beredpar N'$ જરૂરી છે.
      અગાઉના સમાંતર $\beta$-નિયમની જેમ અહીં સંબંધ
      સુધાર્યો છે.}
    \item \ollabel{defn:beredpar3} જો $P \beredpar P'$ અને
      $Q \beredpar Q'$, તો $PQ \beredpar P'Q'$.
    \item \ollabel{defn:beredpar4} જો $N \beredpar N'$ અને
      $Q \beredpar Q'$, તો
      $(\lambd[x][N])Q \beredpar \Subst{N'}{Q'}{x}$.
    \item \ollabel{defn:beredpar5} જો $N \beredpar N'$,
      તો $\lambd[x][Nx] \beredpar N'$, જો
      $x \notin FV(N)$.
  \end{enumerate}
\end{defn}

\begin{thm}\ollabel{thm:refl}
  $M \beredpar M$.
\end{thm}

\begin{proof}
  અભ્યાસપ્રશ્ન.
\end{proof}

\begin{prob}
  \olref[lam][cr][pbe]{thm:refl} સાબિત કરો.
\end{prob}

\begin{defn}[$\beta\eta$-પૂર્ણ વિકાસ]\ollabel{defn:becd}
  $M$નો \emph{$\beta\eta$-પૂર્ણ વિકાસ}~$\becd{M}$
  નીચે મુજબ વ્યાખ્યાયિત થાય છે:
  \begin{align}
    \becd{x} &= x \ollabel{defn:becd1} \\
    \becd{(\lambd[x][N])} &= \lambd[x][\becd{N}] \ollabel{defn:becd2}\\
    \becd{(PQ)} &= \becd{P}\becd{Q} && \text{જો $P$ એ $\lambd$-અમૂર્તીકરણ ન હોય}
    \ollabel{defn:becd3} \\
    \becd{((\lambd[x][N])Q)} &= \Subst{\becd{N}}{\becd{Q}}{x}
                             \ollabel{defn:becd4} \\
    \becd{(\lambd[x][Nx])} &= \becd{N} \ollabel{defn:becd5} & \text{જો $x
                                                      \notin FV(N)$}
  \end{align}
  પાંચમી કલમ લાગુ પડતી હોય ત્યારે તેને બીજી કલમથી
  પહેલાં લાગુ કરવાની છે; બીજી કલમ બાકી અમૂર્તીકરણો
  માટે છે. નહિ તો એક જ $\eta$-સંકોચ્ય પદને બંને
  કલમો જુદાં મૂલ્યો આપે.
  \sourcecorrection{OLLAM-040}{સ્થિર અંગ્રેજી મૂળની
  બીજી કલમ દરેક અમૂર્તીકરણ માટે છે, જ્યારે પાંચમી
  કલમ કેટલાક અમૂર્તીકરણો માટે જુદું પરિણામ આપે છે.
  આથી વ્યાખ્યા પરસ્પર આવરી લેતી અને અસંગત બને છે.
  અહીં પાંચમી કલમને પ્રાથમિકતા આપી બીજી કલમને
  બાકી કિસ્સાઓ સુધી મર્યાદિત કરી છે; આ પસંદગીની
  અનુરૂપ સાબિતી આગળ હજી જરૂરી છે.}
\end{defn}

\begin{lem}\ollabel{lem:comp}
  $M \beredpar M'$ અને $R \beredpar R'$ હોય તો
  $\Subst{M}{R}{y} \beredpar
  \Subst{M'}{R'}{y}$.
\end{lem}

\begin{proof}
  $M \beredpar M'$ની નિષ્પત્તિ પર આગમનથી.

  પહેલા ચાર કિસ્સા \olref[pb]{lem:comp} જેવા છે;
  ત્યાં નોંધાયેલો સ્થાનાપન્ન-સંયોજનનો સાબિતી-ગાળો
  અહીં પણ રહે છે. છેલ્લો નિયમ
  \olref{defn:beredpar5} હોય તો
  $M$ એ $\lambd[x][Nx]$ અને $M'$ એ $N'$, કોઈક
  $x$ અને~$N'$ માટે, જ્યાં $x \notin FV(N)$ અને
  $N \beredpar N'$. બતાવવું છે કે
  $\Subst{(\lambd[x][Nx])}{R}{y} \beredpar
  \Subst{N'}{R'}{y}$. $x$ને $y$, $R$ અને~$R'$થી
  તાજો પસંદ કરીએ તો જરૂરી દાવો
  $\lambd[x][\Subst{N}{R}{y} x] \beredpar
  \Subst{N'}{R'}{y}$ બને છે. અહીં
  $x \notin FV(\Subst{N}{R}{y})$ પણ છે.
  તેથી આગમનની ધારણા અને
  \olref{defn:beredpar}\olref{defn:beredpar5}
  લાગુ પડે છે.
  \sourcecorrection{OLLAM-041}{સ્થિર અંગ્રેજી મૂળ
  $\eta$-કિસ્સામાં સ્થાનાપન્ન પછીનો બાંધનાર~$x$
  તાજો રહે છે એમ કોઈ શરત વિના માને છે. વર્ગના
  પ્રતિનિધિમાં~$x$ને $y$, $R$ અને~$R'$થી તાજો
  પસંદ કરીને જ પાંચમો નિયમ લગાડી શકાય. પહેલા
  ચાર કિસ્સા પણ \olref[pb]{lem:comp}ના અગાઉના
  ખુલ્લા સંયોજન-ગાળાને વારસામાં લે છે.}
\end{proof}

\begin{lem}\ollabel{lem:cont}
  $M \beredpar M'$ હોય તો
  $M' \beredpar \becd{M}$.
\end{lem}

\begin{proof}
  $M \beredpar M'$ની નિષ્પત્તિ પર આગમનથી.

  સ્થિર મૂળ પહેલા ચાર કિસ્સા
  \olref[pb]{lem:cont} જેવા કહે છે. પરંતુ
  $M$ એ $\eta$-સંકોચ્ય અમૂર્તીકરણ હોય ત્યારે
  ઉપર પ્રાથમિકતા આપેલી પાંચમી કલમથી
  $\becd{M}$ નક્કી થાય છે; સામાન્ય
  અમૂર્તીકરણ-કિસ્સાની જૂની દલીલ એ પરિણામ સુધી
  પહોંચે છે એવું અહીં સાબિત થયું નથી.
  \sourcecorrection{OLLAM-042}{સ્થિર અંગ્રેજી મૂળની
  પૂર્ણ-વિકાસ વ્યાખ્યાના બીજા અને પાંચમા નિયમો
  અથડાય છે. પાંચમાને પ્રાથમિકતા આપ્યા પછી
  અમૂર્તીકરણ હેઠળના સમાંતર પગલાં માટે સંપૂર્ણ
  વિકાસ સુધીનું એક વધુ પગલું મળે છે તે અલગ
  સાબિત કરવું પડે; ``પહેલા ચાર કિસ્સા અગાઉ જેવા''
  કહેવું પૂરતું નથી. આ સાબિતી-ગાળો ખુલ્લો છે.}
  છેલ્લો નિયમ \olref{defn:beredpar5} હોય તો
  $M$ એ $\lambd[x][Nx]$ અને $M'$ એ $N'$, કોઈક
  $x$, $N$, $N'$ માટે, જ્યાં $x \notin FV(N)$ અને
  $N \beredpar N'$. બતાવવું છે કે
  $N' \beredpar \becd{(\lambd[x][Nx])}$,
  એટલે પાંચમી કલમ પ્રમાણે
  $N' \beredpar \becd{N}$. આ આગમનની ધારણાથી
  તરત મળે છે, પરંતુ ઉપરનો અમૂર્તીકરણ-કિસ્સો
  હજી ખુલ્લો રહે છે.
\end{proof}

\begin{thm}\ollabel{thm:cr}
  $\beredpar$ ચર્ચ--રોસર ગુણધર્મ ધરાવે છે.
\end{thm}

\begin{proof}
  \olref{lem:cont}માંથી આ પરિણામ મળે, જો તેની
  ઉપર નોંધેલી સાબિતી-ખામી પૂરી કરવામાં આવે.
\end{proof}
% END OLP-0371
\endgroup

\begingroup
\makeatletter\def\ol@sectioncs{section}\makeatother

% BEGIN OLP-0372
\olfileid{lam}{cr}{be}

\olsection{$\beta\eta$-લઘુકરણ}
\phantomsection\label{guunit:OLP-0372}


$\beta\eta$-લઘુકરણ ($\bered$) માટે ચર્ચ--રોસર
ગુણધર્મ લાગુ પડે છે. અહીં $\beredone$નો અર્થ
$\bredone$ અથવા $\eredone$માંથી કોઈ એકનું
એક-પગલાનું સંકોચન છે.
\sourcecorrection{OLLAM-043}{સ્થિર અંગ્રેજી મૂળ
$\beredone$ સંકેત વાપરે છે, પરંતુ આ પ્રકરણમાં
તેના સંબંધની સ્પષ્ટ વ્યાખ્યા આપતું નથી. અગાઉ
$\bered$ બંને પ્રકારના સંકોચનનું સ્વવાચક-પરંપરિત
સંવરણ છે; તેથી અહીં એક-પગલાના સંબંધને
$\bredone$ અને $\eredone$ના સંઘ તરીકે સ્પષ્ટ
કર્યો છે.}

\begin{lem}\ollabel{lem:one-par}
  $M \beredone M'$ હોય તો $M \beredpar M'$.
\end{lem}

\begin{proof}
  $M \beredone M'$ની નિષ્પત્તિ પર આગમનથી.
  મૂળ $\beta$-સંકોચન હોય તો
  \olref[pbe]{thm:refl} અને
  \olref[pbe]{defn:beredpar4} વાપરીએ; મૂળ
  $\eta$-સંકોચન, એટલે $M \eredone M'$, હોય તો
  \olref[syn][eta]{defn:beredone}ની તાજગી-શરત,
  \olref[pbe]{thm:refl} અને
  \olref[pbe]{defn:beredpar5} વાપરીએ.
  અમૂર્તીકરણ અથવા અનુપ્રયોગની અંદર સંકોચન
  થાય ત્યારે \olref[b]{lem:one-par}ની જેમ
  આગમનની ધારણા અને અનુક્રમે
  \olref[pbe]{defn:beredpar2} અથવા
  \olref[pbe]{defn:beredpar3} વાપરીએ.
  આથી દરેક એક-પગલાનું $\beta\eta$-સંકોચન
  સમાંતર પગલું છે.
  \sourcecorrection{OLLAM-044}{સ્થિર અંગ્રેજી મૂળ
  ``$M \bredone M'$ એ $\eta$-રૂપાંતરણથી'' એમ
  કહે છે; $\bredone$ તો ફક્ત $\beta$-સંકોચન
  દર્શાવે છે. અહીં $\eta$ માટે $\eredone$
  લીધો છે. મૂળની ટૂંકી દલીલમાં રચના-સુસંગત
  સંબંધના અંદરના કિસ્સા પણ ગાયબ હતા; તે
  આગમનથી ઉમેર્યા છે.}
\end{proof}

\begin{lem}\ollabel{lem:par-red}
  $M \beredpar M'$ હોય તો $M \bered M'$.
\end{lem}

\begin{proof}
  $M \beredpar M'$ની નિષ્પત્તિ પર આગમનથી.
  પહેલા ચાર નિયમો માટે
  \olref[b]{lem:par-red} જેવી દલીલ છે.

  છેલ્લો નિયમ \olref[pbe]{defn:beredpar5} હોય તો
  $M$ એ $\lambd[x][Nx]$ અને $M'$ એ $N'$, કોઈક
  $x$, $N$, $N'$ માટે, જ્યાં
  $x \notin FV(N)$ અને $N \beredpar N'$.
  પહેલાં $\eta$-રૂપાંતરણથી
  $\lambd[x][Nx]$નું $N$માં લઘુકરણ થાય.
  પછી આગમનની ધારણાથી $N \bered N'$ માટેની
  $\beredone$-પગલાંની શ્રેણી લાગુ કરીએ;
  તેથી $\lambd[x][Nx] \bered N'$.
\end{proof}

\begin{lem}\ollabel{lem:str}
  $\bered$ એ $\beredpar$ને સમાવતો સૌથી નાનો
  પરંપરિત સંબંધ છે.
\end{lem}

\begin{proof}
  \olref[b]{lem:str}ની દલીલ જેવી જ; શૂન્ય પગલાં
  માટે \olref[pbe]{thm:refl} વપરાય છે.
\end{proof}

\begin{thm}\ollabel{thm:cr}
  $\bered$ ચર્ચ--રોસર ગુણધર્મ ધરાવે છે.
\end{thm}

\begin{proof}
  \olref[dap]{thm:str}, \olref[pbe]{thm:cr} અને
  \olref{lem:str}થી દાવો મળે છે, જો
  \olref[pbe]{thm:cr}ની અગાઉ નોંધેલી સાબિતી-ખામી
  પૂરી કરવામાં આવે.
\end{proof}
% END OLP-0372
\endgroup


\OLEndChapterHook
% END OLP-0367
\endgroup


\begingroup
\makeatletter\def\ol@sectioncs{section}\makeatother

% BEGIN OLP-0373
\olchapter{lam}{ldf}{લૅમ્બડા દ્વારા વ્યાખ્યાયિતતા}
\olchapterid{lam}{ldf}
\phantomsection\label{guunit:OLP-0373}

\sourcecorrection{OLLAM-045}{સ્થિર અંગ્રેજી મૂળના
અધ્યાય-ચાલકમાં ઓળખ~\texttt{rep} છે, પરંતુ આ
અધ્યાયના મોટા ભાગના વિભાગોની ઓળખ~\texttt{ldf}
છે. આંતરિક સંદર્ભો એકરૂપ રહે તે માટે અહીં
અધ્યાયની ઓળખ~\texttt{ldf} લીધી છે; જુદી
ઓળખવાળા વિભાગો તેમના અનુવાદ સમયે અલગ
સુધારવાના છે.}

\begin{editorial}
  આ અધ્યાય પ્રાયોગિક છે. તેમાં વધુ સમજૂતી અને
  ઉદાહરણો જોઈએ. સામગ્રીને વ્યાખ્યાઓ અને
  સાબિતીવાળાં પ્રસ્તાવોમાં વધુ સારી રીતે ગોઠવવી
  જોઈએ.
\end{editorial}

\begingroup
\makeatletter\def\ol@sectioncs{section}\makeatother

% BEGIN OLP-0374
\olfileid{lam}{ldf}{int}
\olsection{પરિચય}
\phantomsection\label{guunit:OLP-0374}


પ્રથમ નજરે લૅમ્બડા કલન માત્ર વિધેયો અને તેમનો બીજા
પદો પર અનુપ્રયોગ દર્શાવતી અભિવ્યક્તિઓનું ખૂબ અમૂર્ત
કલન લાગે છે. તેની વાક્યરચનામાં આ વિધેયો કયા પ્રકારની
વસ્તુઓ પરનાં છે તેનો કોઈ નિર્દેશ નથી; ખાસ કરીને
પ્રાકૃતિક સંખ્યાઓનો ઉલ્લેખ નથી. તેની મૂળ ક્રિયાઓ—
અનુપ્રયોગ અને લૅમ્બડા અમૂર્તીકરણ—ફક્ત પ્રાકૃતિક
સંખ્યાઓ પરનાં નહીં, કોઈપણ વિધેયને લાગુ પડે છે.

તેમ છતાં થોડી સૂઝ વાપરીને લૅમ્બડા કલનમાં
અંકગણિતીય વિધેયો, એટલે પ્રાકૃતિક સંખ્યાઓ પરનાં
વિધેયો, વ્યાખ્યાયિત કરી શકાય છે. તે માટે દરેક
પ્રાકૃતિક સંખ્યા~$n \in \Nat$ને ખાસ
$\lambd$-પદ~$\num n$ આપીએ છીએ; તેને~$n$નો
\emph{ચર્ચ અંક} કહે છે. ચર્ચ અંકોનું નામ એલોન્ઝો
ચર્ચ પરથી પડ્યું છે.

\begin{defn}
  $n \in \Nat$ હોય તો તેને અનુરૂપ \emph{ચર્ચ અંક}
  $\num{n}$ એ~$n$નું નિરૂપણ છે:
  \[
    \num{n} \ident \lambd[fx][f^n(x)]
  \]
  અહીં $f^n(x)$ એટલે~$f$ને~$x$ પર $n$ વાર
  લાગુ કરવાનું પરિણામ. દાખલા તરીકે,
  $\num{0}$ એ $\lambd[fx][x]$ છે અને
  $\num{3}$ એ $\lambd[fx][f(f(f\,x))]$ છે.
\end{defn}
  
ચર્ચ અંક~$\num n$ એ લૅમ્બડા પદ તરીકે સંકેતિત
છે. તે બે આર્ગ્યુમેન્ટ~$f$ અને~$x$ સ્વીકારી
$f^n(x)$ પાછું આપતું વિધેય દર્શાવે છે. ચર્ચ
અંકો સ્પષ્ટપણે સામાન્ય સ્વરૂપમાં છે.

લૅમ્બડા કલનમાં પ્રાકૃતિક સંખ્યાઓનું નિરૂપણ ત્યારે જ
ઉપયોગી બને જો તેમની સાથે ગણતરી કરી શકીએ. ચર્ચ અંક
સાથે ગણતરી એટલે $\lambd$-પદ~$F$ને એ અંક પર
લાગુ કરીને સંયુક્ત પદ~$F\, \num n$નું સામાન્ય
સ્વરૂપમાં લઘુકરણ કરવું. આ હંમેશા સામાન્ય સ્વરૂપમાં
લઘુકૃત થાય અને તે સ્વરૂપ હંમેશા ચર્ચ અંક~$\num m$
હોય તો ગણતરીનું પરિણામ સંખ્યા~$m$ ગણી શકાય.
પછી~$F$ વિધેય $f\colon \Nat \to \Nat$ને
વ્યાખ્યાયિત કરે છે એમ કહી શકીએ; ચોક્કસ રીતે,
$f(n) = m$ ત્યારે અને માત્ર ત્યારે જ જ્યારે
$F\, \num n \red \num m$. ચર્ચ--રોસર ગુણધર્મને
કારણે સામાન્ય સ્વરૂપ અસ્તિત્વમાં હોય તો અનન્ય છે.
તેથી $F\, \num n \red \num m$ હોય તો
$F \, \num n$નું લઘુકરણ બીજા કોઈ સામાન્ય
સ્વરૂપમાં, ખાસ કરીને બીજા ચર્ચ અંકમાં, થઈ શકે નહીં.

ઊલટું, વિધેય $f\colon \Nat \to \Nat$ આપેલું હોય
તો આ રીતે વિધેય~$f$ને વ્યાખ્યાયિત કરતું પદ~$F$ છે કે
નહીં તે પૂછી શકાય. હોય તો~$F$ વિધેય~$f$ને
\emph{લૅમ્બડા દ્વારા વ્યાખ્યાયિત} કરે છે અને~$f$
લૅમ્બડા દ્વારા વ્યાખ્યાયિત છે એમ કહીએ. આ વિચારને
અનેક-સ્થાનીય અને આંશિક વિધેયો સુધી વિસ્તારી શકાય.

\begin{defn}
માનો કે $f\colon \Nat^k \pto \Nat$. લૅમ્બડા
પદ~$F$ એ~$f$ને \emph{લૅમ્બડા દ્વારા વ્યાખ્યાયિત}
કરે છે એમ કહીએ, જો બધા
$n_0$, \dots,~$n_{k-1}$ માટે,
\[
F \, \num{n_0} \, \num{n_1} \dots \num{n_{k-1}} \red \num{f(n_0, n_1, \dots,
  n_{k-1})}
\]
જ્યારે $f(n_0, \dots, n_{k-1})$ વ્યાખ્યાયિત
હોય; અને બાકી કિસ્સામાં
$F \, \num{n_0} \, \num{n_1} \dots \num{n_{k-1}}$
ને કોઈ સામાન્ય સ્વરૂપ ન હોય.
\sourcecorrection{OLLAM-046}{સ્થિર અંગ્રેજી મૂળમાં
અહીં $f\colon \Nat^k \to \Nat$ લખાયું છે,
પરંતુ તરત જ $f$ અવ્યાખ્યાયિત હોય તેવો કિસ્સો
આપ્યો છે. આંશિક વિધેય માટે યોગ્ય સંકેત
$f\colon \Nat^k \pto \Nat$ અહીં વાપર્યો છે.}
\end{defn}

અચલ વિધેયો સરળ ઉદાહરણ છે. પદ
$C_k \ident \lambd[x][\num{k}]$ એ
$c_k\colon \Nat \to \Nat$ને લૅમ્બડા દ્વારા વ્યાખ્યાયિત
કરે છે, જ્યાં $c_k(n) = k$. ખરેખર, દરેક~$n$
માટે $C_k \, \num n \ident
(\lambd[x][\num{k}])\num n \redone \num{k}$.
ઓળખ-વિધેય $\lambd[x][x]$ વડે
લૅમ્બડા દ્વારા વ્યાખ્યાયિત છે. વધુ જટિલ વિધેયો વ્યાખ્યાયિત
કરવા માટે ઘણી વાર વધુ સૂઝ જોઈએ. તેથી દરેક
સંગણનીય વિધેય લૅમ્બડા દ્વારા વ્યાખ્યાયિત છે એ કદાચ
આશ્ચર્યજનક લાગે. ઊલટું પણ સાચું છે: વિધેય
લૅમ્બડા દ્વારા વ્યાખ્યાયિત હોય તો તે સંગણનીય છે.
\sourcecorrection{OLLAM-047}{સ્થિર અંગ્રેજી મૂળમાં
અચલ વિધેયને~$c_k$ નામ આપ્યા પછી શરત
$c(n)=k$ લખે છે. અહીં એ જ વિધેયનું નામ
$c_k(n)=k$ રાખ્યું છે.}
% END OLP-0374
\endgroup

\begingroup
\makeatletter\def\ol@sectioncs{section}\makeatother

% BEGIN OLP-0375
\olfileid{lam}{ldf}{arf}
\olsection{લૅમ્બડા દ્વારા વ્યાખ્યાયિત અંકગણિતીય વિધેયો}
\phantomsection\label{guunit:OLP-0375}

\sourcecorrection{OLLAM-048}{સ્થિર અંગ્રેજી મૂળમાં આ વિભાગની
અધ્યાય-ઓળખ~\texttt{rep} છે, જ્યારે અધ્યાયની સુધારેલી
ઓળખ~\texttt{ldf} છે. અધ્યાયની અંદરના સંદર્ભો
એકરૂપ રહે તે માટે અહીં~\texttt{ldf} વાપર્યું છે.}

\begin{prop}
  \ollabel{prop:succ-ld}
  અનુગામી વિધેય~$\Succ$ લૅમ્બડા દ્વારા વ્યાખ્યાયિત છે.
\end{prop}

\begin{proof}
અનુગામી વિધેયને લૅમ્બડા દ્વારા વ્યાખ્યાયિત કરતું એક પદ આ છે:
\begin{align*}
  \fn{Succ} & \ident \lambd[a][\lambd[fx][f ({a} f x)]].
  \intertext{આપણી રૂઢિ અનુસાર આ નીચેના પદનું ટૂંકું સ્વરૂપ છે:}
  \fn{Succ} & \ident \lambd[a][\lambd[f][\lambd[x][(f (({a} f) x))]]].
  \intertext{$\fn{Succ}$ આર્ગ્યુમેન્ટ તરીકે સંખ્યા~$a$ લેતું વિધેય છે
    અને તેનું મૂલ્ય બીજું વિધેય $\lambd[fx][f ({a}
      f x)]$ છે. એ વિધેય પોતે ચર્ચ અંક નથી. જોકે,
    આર્ગ્યુમેન્ટ~$a$ ચર્ચ અંક હોય તો પરિણામ ચર્ચ અંકમાં
    લઘુકૃત થાય છે. જુઓ:}
   (\lambd[a][\lambd[fx][f ({a} f x)]])\,\num{n} & \redone
   \lambd[fx][f ({\num{n}} f x)].
   \intertext{અંદરનું પદ $\num{n}fx$ લઘુકરણીય છે, કારણ કે
     $\num{n}$ એ $\lambd[fx][f^nx]$ છે. તેથી
     $\num{n}fx \red f^nx$ અને આખા પદ માટે મળે છે:}
   \fn{Succ}\, \num n & \red \lambd[fx][f(f^n(x))],
\end{align*}
એટલે કે, $\num{n+1}$.
\sourcecorrection{OLLAM-049}{સ્થિર અંગ્રેજી મૂળ
$\num{n}fx$થી~$f^nx$ સુધી એક જ $\redone$
પગલું બતાવે છે. ચર્ચ અંકના $f$ અને~$x$ બંને
આર્ગ્યુમેન્ટ લાગુ કરવા બે બીટા-પગલાં જોઈએ;
અહીં બહુ-પગલાનું~$\red$ લખ્યું છે.}
\end{proof}

\begin{ex}
  $\fn{Succ}$ને~$\num{0}$, એટલે કે
  $\lambd[fx][x]$, પર લાગુ કરીએ તો શું થાય તે જોઈએ.
  પદોને પૂરેપૂરાં લખીએ:
  \begin{align*}
    \fn{Succ}\, \num{0} & \ident (\lambd[a][\lambd[f][\lambd[x][(f (({a} f) x))]]])(\lambd[f][\lambd[x][x]])\\
    & \redone \lambd[f][\lambd[x][(f (({(\lambd[f][\lambd[x][x]])} f) x))]] \\
    & \redone \lambd[f][\lambd[x][(f ({(\lambd[x][x])} x))]] \\
    & \redone \lambd[f][\lambd[x][(f x)]] \ident \num{1}\\
  \end{align*}
\end{ex}

\begin{prob}
  પદ
  \begin{align*}
    \fn{Succ}' & \ident \lambd[{n}][\lambd[fx][{n} f (f x)]]
  \end{align*}
  અનુગામી વિધેયને લૅમ્બડા દ્વારા વ્યાખ્યાયિત કરે છે. કેમ, તે સમજાવો.
\end{prob}

\begin{prop}
  \ollabel{prop:add-ld}
  સરવાળાનું વિધેય~$\Add$ લૅમ્બડા દ્વારા વ્યાખ્યાયિત છે.
\end{prop}

\begin{proof}
નીચેનાં પદો સરવાળાને લૅમ્બડા દ્વારા વ્યાખ્યાયિત કરે છે:
\begin{align*}
  \fn{Add} & \ident \lambd[{a}{b}][\lambd[fx][{a} f ({b} f x)]]
\intertext{અથવા વૈકલ્પિક રીતે,}
  \fn{Add}' & \ident \lambd[{a}{b}][{a}\, \fn{Succ} \, {b}].
\intertext{પહેલું પદ આ રીતે કામ કરે છે: $\fn{Add}$ પહેલાં
  બે સંખ્યાઓ ${a}$ અને ${b}$ લે છે. પરિણામે મળતું
  વિધેય $f$ અને~$x$ લે છે તથા $af(bfx)$ આપે છે.
  $a$ અને~$b$ ચર્ચ અંકો $\num{n}$ અને~$\num{m}$
  હોય તો આ $f^{n+m}(x)$માં લઘુકૃત થાય છે;
  એ $f^{n}(f^{m}(x))$ જેટલું જ છે. પગલે પગલે:}
  (\lambd[{a}{b}][\lambd[fx][{a} f ({b} f x)]])\num n\,\num m & \red
  \lambd[fx][\num{n}\, f (\num {m}\, f x)] \\
  & \red \lambd[fx][\num{n}\, f (f^m x)] \\
  & \red \lambd[fx][f^n (f^m x)] \ident \num {n+m}.
\intertext{સરવાળાનું બીજું નિરૂપણ $\fn{Add'}$ જુદી રીતે
  કામ કરે છે. બે ચર્ચ અંકો $\num{n}$ અને
  $\num{m}$ પર લાગુ કરીએ તો}
\fn{Add}' \num n \,\num m
& \red \num{n}\, \fn{Succ}\, \num{m}.
\intertext{પરંતુ $\num{n} f x$ હંમેશા $f^n(x)$માં
  લઘુકૃત થાય છે. તેથી}
  \num{n}\,  \fn{Succ}\, \num{m} & \red \fn{Succ}^n(\num{m}).
\end{align*}
અને $\fn{Succ}$ અનુગામી વિધેયને લૅમ્બડા દ્વારા વ્યાખ્યાયિત
કરે છે તથા~$m$ પર તેને~$n$ વાર લાગુ કરીએ તો
$n+m$ મળે છે; તેથી આખું પદ~$\num{n+m}$માં
લઘુકૃત થાય છે.
\sourcecorrection{OLLAM-050}{સ્થિર અંગ્રેજી મૂળમાં
સરવાળાના પ્રદર્શનનાં અનેક $\redone$ ચિહ્નો એક-પગલાનું
લઘુકરણ બતાવે છે, જોકે બે આર્ગ્યુમેન્ટ લગાડવા
અને ચર્ચ અંકનું લઘુકરણ કરવા અનેક બીટા-પગલાં
જોઈએ. અહીં એ સ્થળોએ બહુ-પગલાનું~$\red$ વાપર્યું છે.}
\end{proof}

\begin{prop}
  \ollabel{prop:mult-ld}
  ગુણાકાર આ પદ વડે લૅમ્બડા દ્વારા વ્યાખ્યાયિત છે:
  \[
  \fn{Mult} \ident \lambd[ab][\lambd[fx][a (b f) x]]
  \]
\end{prop}

\begin{proof}
આ પદ કેવી રીતે કામ કરે છે તે જોવા $\fn{Mult}$ને ચર્ચ
અંકો $\num{n}$ અને $\num{m}$ પર લાગુ કરીએ:
$\fn{Mult} \, \num{n} \, \num{m}$નું લઘુકરણ
$\lambd[fx][\num{n}(\num{m}\, f)x]$માં થાય છે.
$\num{m} f$ એવું વિધેય દર્શાવે છે જે પોતાના
આર્ગ્યુમેન્ટ પર~$f$ને $m$ વાર લાગુ કરે છે. તેથી
$\num{n} (\num{m} f) x$ એ ``$f$ને $m$
વાર લાગુ કરો'' એવા વિધેયને જ~$x$ પર $n$ વાર
લાગુ કરે છે. બીજા શબ્દોમાં, $f$ને~$x$ પર
$n\cdot m$ વાર લાગુ કરીએ છીએ. પરિણામે મળતું
સામાન્ય સ્વરૂપ ચર્ચ અંક~$\num{nm}$ જ છે.
\end{proof}

\begin{editorial}
$\eta$-લઘુકરણ વડે આ પદને હજી વધુ સરળ કરી શકાય:
\[
  \fn{Mult} \ident \lambd[ab][\lambd[f][a (b f)]].
\]

પણ એ પહેલાં $\eta$-લઘુકરણ સમજાવવું પડે.
\end{editorial}

\begin{prob}
ગુણાકાર આ પદ વડે લૅમ્બડા દ્વારા વ્યાખ્યાયિત કરી શકાય છે:
\[
  \fn{Mult}' \ident \lambd[ab][b (\fn{Add}\, a) \num{0}].
\]
તે કેમ કામ કરે છે તે સમજાવો.
\sourcecorrection{OLLAM-051}{સ્થિર અંગ્રેજી મૂળમાં
$\fn{Mult}'$ના મુખ્ય ભાગમાં
$a(\fn{Add}\,a)\num{0}$ છે. એ~$a$ને
$a$ વાર ઉમેરે છે અને બીજો આર્ગ્યુમેન્ટ~$b$
વાપરતું નથી; એટલે સામાન્ય રીતે $ab$ને બદલે
$a^2$ આપે છે. અહીં~$b(\fn{Add}\,a)\num{0}$
રાખ્યું છે, જેથી~$a$ને $b$ વાર ઉમેરીને~$ab$ મળે.}
\end{prob}

ઘાતાંકને $\lambd$-પદ તરીકે વ્યાખ્યાયિત કરવો
આશ્ચર્યજનક રીતે સરળ છે:
\[
  \fn{Exp} \ident \lambd[be][e b].
\]
પહેલો આર્ગ્યુમેન્ટ~$b$ આધાર છે અને બીજો~$e$
ઘાતાંક છે. સંખ્યાઓના આપણા સંકેતન પ્રમાણે
$e f$ એ $f^e$ છે એમ સહજ રીતે સમજાય છે.
આ સમજવું મુશ્કેલ લાગે તો પુનરાવર્તિત ગુણાકાર
વડે પણ ઘાતાંક વ્યાખ્યાયિત કરી શકાય:
\[
  \fn{Exp}' \ident \lambd[be][e (\fn{Mult}\, b) \num{1}].
\]

ચર્ચ અંક પર પૂર્વગામી અને બાદબાકી વ્યાખ્યાયિત
કરવાનું આપણી ધારણા કરતાં ઓછું સરળ છે:
તે માટે જોડીઓનું સંકેતન જોઈએ.
% END OLP-0375
\endgroup

\begingroup
\makeatletter\def\ol@sectioncs{section}\makeatother

% BEGIN OLP-0376
\olfileid{lam}{ldf}{pai}
\olsection{જોડીઓ અને પૂર્વગામી}
\phantomsection\label{guunit:OLP-0376}


\begin{defn}
$M$ અને~$N$ની જોડી ($\tuple{M,N}$ તરીકે લખાતી)
આ રીતે વ્યાખ્યાયિત થાય છે:
\[
\tuple{M,N} \ident \lambd[f][fMN].
\]
\end{defn}
  
સહજ રીતે, આ એવું વિધેય છે જે બીજું વિધેય
સ્વીકારે છે અને તેને જોડીના બંને ઘટકો પર લાગુ
કરે છે. આ વિચાર પરથી બે પદો લઈને તેમની
જોડી આપતું રચનાપદ મળે છે:
\[
  \fn{Pair} \ident \lambd[mn][\lambd[f][fmn]]
\]
જોડી આપેલી હોય ત્યારે તેના ઘટકો પાછા પણ
મેળવવા છે. તેના માટે બે પ્રવેશ-વિધેયો જોઈએ;
તે જોડી આર્ગ્યુમેન્ટ તરીકે લે છે અને અનુક્રમે
પહેલો કે બીજો ઘટક આપે છે:
\begin{align*}
  \fn{Fst} & \ident \lambd[p][p(\lambd[mn][m])]\\
  \fn{Snd} & \ident \lambd[p][p(\lambd[mn][n])]
\end{align*}

\begin{prob}
  પ્રવેશ-વિધેયો $\fn{Fst}$ અને $\fn{Snd}$
  કેમ કામ કરે છે તે સમજાવો.
\end{prob}

હવે જોડીઓ વડે પૂર્વગામી વિધેયને લૅમ્બડા દ્વારા વ્યાખ્યાયિત
કરી શકીએ:
\[
  \fn{Pred} \ident \lambd[n][\fn{Fst}(n (\lambd[p][\tuple{\fn{Snd}\, {p}, \fn{Succ}(\fn{Snd}\, {p})}]) \tuple{\num 0, \num 0})]
\]
યાદ રાખો કે $\num n\, f x$નું લઘુકરણ
$f^{n}(x)$માં થાય છે. અહીં~$f$ જોડી~$p$
સ્વીકારી જોડી~$p$નો બીજો ઘટક તથા એ બીજા ઘટકનો
અનુગામી ધરાવતી નવી જોડી આપતું વિધેય છે;
$x$ એ જોડી $\tuple{\num 0,\num 0}$ છે.
તેથી $n=0$ માટે પરિણામ $\tuple{\num 0,\num 0}$
અને અન્ય કિસ્સામાં $\tuple{\num{n-1}, \num n}$
મળે છે. પછી~$\fn{Pred}$ એ પરિણામનો
પહેલો ઘટક આપે છે.
\sourcecorrection{OLLAM-052}{સ્થિર અંગ્રેજી મૂળ
ગદ્યમાં શરૂઆતની અને શૂન્ય માટેની જોડીને
$\tuple{0,0}$ લખે છે, જ્યારે પદની વ્યાખ્યામાં
ચર્ચ અંકોની જોડી~$\tuple{\num 0,\num 0}$
છે. અહીં પદ-સ્તરે ચર્ચ અંક જ લખ્યા છે.}

બાદબાકી~$a$ પર $\fn{Pred}$ને $b$ વાર
લાગુ કરીને વ્યાખ્યાયિત કરી શકાય:
\[
\fn{Sub} \ident \lambd[ab][b \fn{Pred}\, a].
\]
% END OLP-0376
\endgroup

\begingroup
\makeatletter\def\ol@sectioncs{section}\makeatother

% BEGIN OLP-0377
\olfileid{lam}{ldf}{tvr}
\olsection{સત્યમૂલ્યો અને સંબંધો}
\phantomsection\label{guunit:OLP-0377}


શુદ્ધ લૅમ્બડા કલનમાં સત્યમૂલ્યો આ રીતે
સંકેતિત કરી શકાય:
\begin{align*}
  \fn{true} & \ident \lambd[x][\lambd[y][x]]\\
  \fn{false} & \ident \lambd[x][\lambd[y][y]]
\end{align*}

સત્યમૂલ્યોનું નિરૂપણ \emph{પસંદગી-વિધેયો} તરીકે
થાય છે, એટલે બે આર્ગ્યુમેન્ટ લઈને તેમાંથી એક
પાછું આપતાં વિધેયો તરીકે. સત્યમૂલ્ય
$\fn{true}$ પહેલો આર્ગ્યુમેન્ટ પસંદ કરે છે અને
$\fn{false}$ બીજો. દાખલા તરીકે,
$\fn{true}\, M N$ હંમેશા~$M$માં લઘુકૃત થાય છે,
જ્યારે $\fn{false}\, M N$ હંમેશા~$N$માં થાય છે.

\begin{defn}
સંબંધ $R \subseteq \Nat^k$ને લૅમ્બડા દ્વારા વ્યાખ્યાયિત
કહીએ, જો એવું પદ~$R$ હોય કે
\begin{align*}
  R\, \num{n_1} \dots \num{n_k} & \bred \fn{true}
  \intertext{જ્યારે $R(n_1, \dots, n_k)$ સાચો હોય, અને}
  R\, \num{n_1} \dots \num{n_k} & \bred \fn{false}
\end{align*}
બાકી કિસ્સામાં.
\sourcecorrection{OLLAM-053}{સ્થિર અંગ્રેજી મૂળમાં
સંબંધનું ક્ષેત્ર~$\Nat^n$ છે, પરંતુ તેની દલીલો
$n_1,\dots,n_k$ અને અનુપ્રયોગમાં પણ~$k$
ચર્ચ અંકો છે. અરીટી એકરૂપ રાખવા અહીં
$R\subseteq\Nat^k$ લખ્યું છે.}
\end{defn}

દાખલા તરીકે, જે~$0$ માટે અને ફક્ત~$0$ માટે જ સાચો છે તે સંબંધ
$\fn{IsZero} = \{0\}$ આ પદ વડે
લૅમ્બડા દ્વારા વ્યાખ્યાયિત છે:
\[
  \fn{IsZero} \ident \lambd[n][n (\lambd[x][\fn{false}])\, \fn{true}].
\]
તે કેવી રીતે કામ કરે છે? ચર્ચ અંકો પુનરાવર્તકો
તરીકે વ્યાખ્યાયિત છે: બીજો આર્ગ્યુમેન્ટ લઈને
તેના પર પહેલા આર્ગ્યુમેન્ટને~$n$ વાર લાગુ કરે છે.
અહીં શરૂઆતનું મૂલ્ય~$\fn{true}$ છે, અને
પુનરાવર્તનના દરેક પગલે અગાઉના પગલાનું પરિણામ
ગમે તે હોય, $\fn{false}$ મળે છે. આ પગલું
શરૂઆતના મૂલ્ય પર $n$ વાર લાગુ થાય છે. એટલે
પગલું એકેય વાર ન લાગ્યું હોય, એટલે કે $n = 0$
હોય, ત્યારે અને માત્ર ત્યારે જ પરિણામ
$\fn{true}$ મળે છે.

સત્યમૂલ્યોના આ નિરૂપણ પરથી બીજા સત્ય-વિધેયો
પણ વ્યાખ્યાયિત કરી શકાય. અહીં નિષેધ અને
સંયોજનનાં બે નિરૂપણ છે:
\begin{align*}
  \fn{Not} & \ident \lambd[x][x\, \fn{false}\, \fn{true}]\\
  \fn{And} & \ident \lambd[x][\lambd[y][xy \,\fn{false}]]
\end{align*}
$\fn{Not}$ એક આર્ગ્યુમેન્ટ લે છે અને તે
$\fn{false}$ હોય તો $\fn{true}$, તથા
$\fn{true}$ હોય તો $\fn{false}$ આપે છે.
$\fn{And}$ બે સત્યમૂલ્યો લે છે અને બંને
$\fn{true}$ હોય ત્યારે અને માત્ર ત્યારે જ
$\fn{true}$ આપવું જોઈએ. ઉપર જણાવ્યા પ્રમાણે
સત્યમૂલ્યો પસંદગી-વિધેયો છે: $x$ સત્યમૂલ્ય હોય
અને તેને બે આર્ગ્યુમેન્ટ લાગુ કરીએ તો~$x$
$\fn{true}$ હોય ત્યારે પહેલો, અન્યથા બીજો
આર્ગ્યુમેન્ટ મળે. હવે $\fn{And}$ પોતાના
આર્ગ્યુમેન્ટો $x$ અને~$y$ લઈને પહેલા
આર્ગ્યુમેન્ટ~$x$ને $y$ અને $\fn{false}$
સોંપે છે. $x$ સત્યમૂલ્ય હોય એમ માનીએ તો
$x$ એ $\fn{true}$ હોય ત્યારે પરિણામ~$y$,
અને $x$ એ $\fn{false}$ હોય ત્યારે પરિણામ
$\fn{false}$ મળે છે. આ જ જોઈએ છે.

અહીં આપણે માની લીધું છે કે $\fn{And}$ને
આર્ગ્યુમેન્ટ તરીકે માત્ર સત્યમૂલ્યો જ અપાય છે.
તેને બીજાં પદો આપીએ તો પરિણામ, એટલે સામાન્ય
સ્વરૂપ અસ્તિત્વમાં હોય તો તે, સત્યમૂલ્ય ન પણ હોય.

\begin{prob}
  સત્યમૂલ્યોના $\lambd$-પદ તરીકેના આ સંકેતન
  વડે સમાવેશી તથા અપવર્જી વિયોજનના સત્ય-વિધેયો
  દર્શાવતાં $\fn{Or}$ અને $\fn{Xor}$ વ્યાખ્યાયિત કરો.
\end{prob}
% END OLP-0377
\endgroup

\begingroup
\makeatletter\def\ol@sectioncs{section}\makeatother

% BEGIN OLP-0378
\olfileid{lam}{ldf}{prf}
\olsection{આદિમ પુનરાવર્તી વિધેયો લૅમ્બડા દ્વારા વ્યાખ્યાયિત છે}
\phantomsection\label{guunit:OLP-0378}


યાદ કરો કે આદિમ પુનરાવર્તી વિધેયો એ છે, જે મૂળભૂત
વિધેયો $\Zero$, $\Succ$ અને $\Proj{n}{i}$માંથી
સંયોજન તથા આદિમ પુનરાવર્તન દ્વારા વ્યાખ્યાયિત કરી શકાય.

\begin{lem}
  \ollabel{lem:basic}
  મૂળભૂત આદિમ પુનરાવર્તી વિધેયો $\Zero$, $\Succ$ અને
  પ્રક્ષેપો~$\Proj{n}{i}$ લૅમ્બડા દ્વારા વ્યાખ્યાયિત છે.
\end{lem}

\begin{proof}
નીચેનાં પદો તેમને લૅમ્બડા દ્વારા વ્યાખ્યાયિત કરે છે:
\begin{align*}
  \fn{Zero} & \ident \lambd[a][\lambd[fx][x]]\\
  \fn{Succ} & \ident \lambd[a][\lambd[fx][f (a f x)]] \\
  \fn{Proj}^n_i & \ident \lambd[x_0\dots x_{n-1}][x_i]
\end{align*}
\end{proof}

\begin{lem}
  \ollabel{lem:comp} ધારો કે $k$-સ્થાનીય વિધેય $f$ અને
  $n$-સ્થાનીય વિધેયો $g_0, \dots, g_{k-1}$ અનુક્રમે
  $F$, $G_0$, \dots, $G_{k-1}$ પદો દ્વારા
  લૅમ્બડા દ્વારા વ્યાખ્યાયિત છે, અને $h$ તેમની
  સંયોજનથી વ્યાખ્યાયિત થાય છે. તો $h$ પણ
  લૅમ્બડા દ્વારા વ્યાખ્યાયિત છે.
  \sourcecorrection{OLLAM-054}{સ્થિર અંગ્રેજી મૂળમાં
  $k$ વિધેયો $g_0,\dots,g_{k-1}$ હોવા છતાં
  તેમને દર્શાવતાં પદોનો ક્રમ $G_k$ સુધી જાય છે;
  અહીં છેલ્લું પદ $G_{k-1}$ કર્યું છે. મૂળના
  નિષ્કર્ષમાં ``$H$ is lambda definable'' છે,
  પરંતુ $H$ નિરૂપણ પદ છે; વ્યાખ્યાયિત થતું
  વિધેય $h$ છે.}
\end{lem}

\begin{proof}
  વિધેય~$h$ને આ પદ લૅમ્બડા દ્વારા વ્યાખ્યાયિત કરે છે:
  \[
  H \ident \lambd[x_0 \dots x_{n-1}][F\, (G_0 x_0 \dots
    x_{n-1}) \dots (G_{k-1} x_0 \dots x_{n-1})]
  \]
  આ હકીકતની ચકાસણી કસરત તરીકે છોડી છે.
\end{proof}

\begin{prob}
  $H\num{n_0}\dots\num{n_{n-1}} \red \num{h(n_0, \dots,
    n_{n-1})}$ બતાવીને
  \olref[lam][ldf][prf]{lem:comp}નો પુરાવો પૂર્ણ કરો.
\end{prob}

નોંધો કે \olref{lem:comp}માં $f$ અને $g_0$,
\dots,~$g_{k-1}$ આદિમ પુનરાવર્તી હોવાં જરૂરી
નથી; તેઓ સર્વત્ર વ્યાખ્યાયિત અને લૅમ્બડા દ્વારા વ્યાખ્યાયિત
હોય એટલું જ જરૂરી છે.

\begin{lem}\ollabel{lem:prim}
  ધારો કે $f$ એ $n$-સ્થાનીય વિધેય અને~$g$ એ
  $n+2$-સ્થાનીય વિધેય છે; તેઓ અનુક્રમે
  $F$ અને~$G$ પદો દ્વારા લૅમ્બડા દ્વારા વ્યાખ્યાયિત છે,
  અને $h$ એ $f$ તથા $g$માંથી આદિમ
  પુનરાવર્તન દ્વારા વ્યાખ્યાયિત છે. તો $h$ પણ
  લૅમ્બડા દ્વારા વ્યાખ્યાયિત છે.
\end{lem}

\begin{proof}
  યાદ કરો કે $h$ની વ્યાખ્યા આ છે:
  \begin{align*}
    h(x_1, \dots, x_n, 0) &= f(x_1, \dots, x_n)\\
    h(x_1, \dots, x_n, y+1) & = g(x_1, \dots, x_n, y, h(x_1, \dots, x_n, y)).
  \end{align*}
  \sourcecorrection{OLLAM-055}{સ્થિર અંગ્રેજી મૂળના
  બીજા પુનરાવર્તન સમીકરણમાં જમણી બાજુનું બહારનું
  વિધેય ભૂલથી $h$ છે. પગલું-વિધેય $g$ હોવું
  જોઈએ; પછીના પુરાવામાં પણ $g$ જ વપરાય છે.}
  અનૌપચારિક રીતે, આદિમ પુનરાવર્તનની વ્યાખ્યા
  $f(x_1,\dots,x_n)$થી શરૂ કરીને પગલું-વિધેય~$g$
  $y$ વાર લાગુ કરે છે. આ ચર્ચ અંકોની પુનરાવર્તક
  તરીકેની વ્યાખ્યાની યાદ અપાવે છે.

  સરળતા માટે, એક જ વધારાની દલીલ~$x$ હોય તે
  કિસ્સાની વ્યાખ્યા અને સાબિતી આપીએ છીએ.
  વિધેય~$h$ને આ પદ લૅમ્બડા દ્વારા વ્યાખ્યાયિત કરે છે:
  \begin{align*}
    H \ident & \lambd[x][\lambd[y][\fn{Snd} (y
        D \tuple{\num{0}, F x})]]
    \intertext{જ્યાં }
    D \ident & \lambd[p][\tuple{\fn{Succ} (\fn{Fst}\, p), (G x
          (\fn{Fst}\, p) (\fn{Snd}\, p))}]
  \end{align*}
  \sourcecorrection{OLLAM-056}{અહીં $D$માં $x$
  મુક્ત રહે છે; $H$ના બહારના $\lambd[x]$ની અંદર
  તેનો અર્થ દરેક નક્કી કરેલા $x$ માટેનું પગલું-પદ
  છે. પછીના પુરાવાનું $D_n$ એ $x$ના સ્થાને
  $\num n$ મૂકેલું આ જ પદ છે. આ અવલંબન
  મૂળના સંકેતનમાં સ્પષ્ટ નહોતું.}
  આપણે રાખેલી પુનરાવર્તનની સ્થિતિ એક જોડી છે:
  તેનો પહેલો ઘટક હાલનું~$y$ છે અને બીજો~$h$નું
  અનુરૂપ મૂલ્ય છે. દરેક પગલે $y$ અને~$h$નાં
  નવા મૂલ્યોની જોડી બનાવીએ છીએ; પુનરાવર્તન પૂરું
  થયા પછી તેનો બીજો ઘટક પાછો આપીએ છીએ;
  તે જ $h$નું અંતિમ મૂલ્ય છે.
  હવે આ ખરેખર આદિમ પુનરાવર્તનનું નિરૂપણ છે
  તે સાબિત કરીએ.

  આપણે દરેક $n$ અને $m$ માટે
  $H\,\num{n}\,\num{m} \red \num{h(n,m)}$
  સાબિત કરવું છે. પહેલાં, $D_n \ident
  \Subst{D}{\num{n}}{x}$ હોય તો
  $D_n^m \tuple{\num{0}, F\, \num{n}} \red
  \tuple{\num{m}, \num{h(n, m)}}$ છે એમ
  $m$ પર ગાણિતિક અનુમાનથી બતાવીએ.

  $m=0$ હોય તો
  $D_n^0 \tuple{\num{0}, F\, \num{n}} \red
  \tuple{\num{0}, \num{h(n, 0)}}$ બતાવવું છે.
  પરંતુ $D_n^0 \tuple{\num{0}, F\, \num{n}}$
  એટલે $\tuple{\num{0}, F\, \num{n}}$ જ.
  $F$ એ~$f$ને લૅમ્બડા દ્વારા વ્યાખ્યાયિત કરતું હોવાથી
  આ $\tuple{\num{0}, \num{f(n)}}$માં
  લઘુકૃત થાય છે; અને $f(n)=h(n,0)$ હોવાથી
  તે $\tuple{\num{0}, \num{h(n,0)}}$ છે.

  હવે ધારો કે
  $D_n^m \tuple{\num{0}, F\, \num{n}} \red
  \tuple{\num{m}, \num{h(n, m)}}$.
  આપણે $D_n^{m+1} \tuple{\num{0}, F\, \num{n}} \red
  \tuple{\num{m+1}, \num{h(n, m+1)}}$
  બતાવવું છે.
  \begin{align*}
    D_n^{m+1} \tuple{\num{0}, F\, \num{n}}
    & \ident D_n(D_n^{m} \tuple{\num{0}, F\, \num{n}})\\
    & \red D_n\,\tuple{\num{m}, \num{h(n, m)}} \text{\qquad (અનુમાનથી)}\\
    & \ident (\lambd[p][
      \tuple{
        \fn{Succ} (\fn{Fst}\, p), (G \, \num{n} (\fn{Fst}\, p) (\fn{Snd}\, p))
    }]) \tuple{\num{m}, \num{h(n,m)}}\\
    & \redone
    \tuple{\fn{Succ} (\fn{Fst}\, \tuple{\num{m}, \num{h(n,m)}}),\\
      & \qquad (G \, \num{n} (\fn{Fst}\, \tuple{\num{m}, \num{h(n,m)}}) (\fn{Snd}\, \tuple{\num{m}, \num{h(n,m)}}))}\\
    & \red
    \tuple{\fn{Succ}\, \num{m}, (G \, \num{n} \,\num{m} \, \num{h(n,m)})}\\
    & \red \tuple{\num{m+1}, \num{g(n, m, h(n, m))}}
  \end{align*}
  $g(n,m,h(n,m))=h(n,m+1)$ હોવાથી ઇચ્છિત
  પરિણામ મળે છે.

  અંતે, આ ગણતરી જુઓ:
  \begin{align*}
    H\,\num n\,\num m &\ident \lambd[x][\lambd[y][\fn{Snd} (y
        (\lambd[p].\tuple{\fn{Succ} (\fn{Fst}\, p), (G\, x\,
          (\fn{Fst}\, p)\, (\fn{Snd}\, p))}) \tuple{\num{0}, F x})]]\\
    & \qquad\,\num n\, \num m\\
    & \red \fn{Snd} (\num m \,
    \underbrace{(\lambd[p].\tuple{\fn{Succ} (\fn{Fst}\, p), (G \,\num n\,
      (\fn{Fst}\, p) (\fn{Snd}\, p))})}_{D_n} \tuple{\num{0}, F \num n})\\
    & \ident \fn{Snd} (\num m \, D_n\, \tuple{\num{0}, F \num n})\\
    & \red \fn{Snd} \, (D_n^m \tuple{\num{0}, F \num n})\\
    & \red \fn{Snd} \,\tuple{\num{m}, \num{h(n, m)}} \\
    & \red \num{h(n,m)}.
  \end{align*}
\end{proof}

\begin{prop}
  દરેક આદિમ પુનરાવર્તી વિધેય લૅમ્બડા દ્વારા વ્યાખ્યાયિત છે.
\end{prop}

\begin{proof}
  \olref{lem:basic} મુજબ બધાં મૂળભૂત વિધેયો
  લૅમ્બડા દ્વારા વ્યાખ્યાયિત છે. \olref{lem:comp}
  અને \olref{lem:prim} મુજબ
  લૅમ્બડા દ્વારા વ્યાખ્યાયિત વિધેયો સંયોજન તથા
  આદિમ પુનરાવર્તન હેઠળ બંધ છે.
\end{proof}
% END OLP-0378
\endgroup

\begingroup
\makeatletter\def\ol@sectioncs{section}\makeatother

% BEGIN OLP-0379
\olfileid{lam}{ldf}{fp}

\olsection{સ્થિરબિંદુઓ}
\phantomsection\label{guunit:OLP-0379}


ધારો કે આપણે ક્રમગુણિત વિધેયને પુનરાવર્તન દ્વારા એવા
પદ~$\fn{Fac}$ તરીકે વ્યાખ્યાયિત કરવા માગીએ છીએ, જેમાં નીચેનો ગુણધર્મ હોય:
\[
\fn{Fac} \ident \lambd[n][\fn{IsZero}\, n\, \num 1 (\fn{Mult}\, n (\fn{Fac}(\fn{Pred} \, n)))]
\]
એટલે $n=0$ હોય ત્યારે $n$નું ક્રમગુણિત~$1$ છે અને
અન્યથા તે $n$ ગુણ્યા $n-1$નું ક્રમગુણિત છે. અલબત્ત, આ રીતે
$\fn{Fac}$ પદ વ્યાખ્યાયિત કરી શકાતું નથી, કારણ કે જમણી બાજુએ
$\fn{Fac}$ પોતે જ આવે છે. પોતાનો જ સંદર્ભ લેતી આવી
પુનરાવર્તી વ્યાખ્યાઓ લૅમ્બડા કલનનો ભાગ નથી. દાખલા તરીકે,
\[
\fn{Mult} \ident \lambd[ab][a (\fn{Add}\, a) 0]
\]
વડે પદ વ્યાખ્યાયિત કરીએ ત્યારે જમણી બાજુએ $\fn{Add}$ જેવાં
પહેલેથી વ્યાખ્યાયિત પદો જ આવે છે. $\fn{Add}$ને તેનું વ્યાખ્યાયક
પદ મૂકીને હંમેશા દૂર કરી શકીએ છીએ. આમ $\fn{Mult}$નું
શુદ્ધ લૅમ્બડા પદ મળે; જો $\fn{Add}$ પોતે વ્યાખ્યાયિત પદો
ધરાવતું હોય (જેમ કે $\fn{Add}'$ ધરાવે છે), તો આ પ્રક્રિયા
ચાલુ રાખીને છેવટે શુદ્ધ લૅમ્બડા પદ સુધી પહોંચી શકીએ.

પરંતુ ઉપરની $\fn{Fac}$ જેવી પુનરાવર્તી વ્યાખ્યામાં એવું થતું
નથી. જમણી બાજુએ આવતા $\fn{Fac}$ના સ્થાને
$\fn{Fac}$ની પોતાની વ્યાખ્યા મૂકીએ તો મળે છે:
\begin{align*}
  \fn{Fac} & \ident \lambd[n][\fn{IsZero}\, n\, \num{1}\,
    (\fn{Mult}\, n
  ((\lambd[n][\fn{IsZero} \, n \, \num 1\, (\fn{Mult}\, n\, (\fn{Fac}
    (\fn{Pred}\, n)))]) (\fn{Pred}\, n)))]
\end{align*}
\sourcecorrection{OLLAM-057}{સ્થિર અંગ્રેજી મૂળના આ
અવેજીકરણમાં $\fn{IsZero}$ના ત્રીજા આર્ગ્યુમેન્ટ પહેલાં
લૅમ્બડા પદ બંધ થઈ જાય છે અને $\fn{Mult}$ બહાર રહી જાય છે.
અહીં એ શાખાને પદની અંદર જ રાખી છે, જેથી ઉપરની
પુનરાવર્તી સમીકરણમાં કરેલો અવેજ સાચો રહે.}
જમણી બાજુમાંથી $\fn{Fac}$ હજુ દૂર થયું નથી. આ પ્રક્રિયા
પુનરાવર્તિત કરીએ તો વ્યાખ્યા લાંબી થતી જાય છે, પરંતુ
શુદ્ધ લૅમ્બડા પદ કદી મળતું નથી. તેથી ક્રમગુણિતને
(અથવા સામાન્ય રીતે પુનરાવર્તી વિધેયોને) આ રીતે વ્યાખ્યાયિત
કરવું શક્ય નથી.

તેમ છતાં પુનરાવર્તી વ્યાખ્યા આપણને કંઈક કહે છે: જો
$f$ ક્રમગુણિત વિધેયનું નિરૂપણ કરતું પદ હોય, તો
\[
\fn{Fac}' \ident \lambd[g][\lambd[n][\fn{IsZero} \, n \, \num 1\, (\fn{Mult} \, n\, (g (\fn{Pred} n)))]]
\]
પદને $f$ પર લાગુ કરતાં મળતું $\fn{Fac}'\,f$ પણ ક્રમગુણિત
વિધેયનું નિરૂપણ કરે છે. એટલે $\fn{Fac}'$ને વિધેય સ્વીકારીને
વિધેય પરત કરતું વિધેય ગણીએ તો, $f$ ક્રમગુણિત વિધેયનું
નિરૂપણ કરતું હોય ત્યારે $\fn{Fac'}\, f$ એ જ વિધેયનું
નિરૂપણ કરે છે. જે પદ~$f$ માટે વધુમાં $\fn{Fac}' \, f
\equal[\beta] f$ હોય, તેને $\fn{Fac}'$નું
\emph{સ્થિરબિંદુ} કહે છે. વિધેયોના સ્તરે ક્રમગુણિત
$\fn{Fac}'$ હેઠળ અચળ રહે છે; આવું બીટા-સ્થિરબિંદુ
આપતું ચોક્કસ પદ~$f$ નીચે બનાવશું.
\sourcecorrection{OLLAM-059}{સ્થિર અંગ્રેજી મૂળમાં
બંને પદો એક જ ક્રમગુણિત વિધેયનું નિરૂપણ કરે છે
એ હકીકતથી પદોની $\beta$-સમતા તરફ સીધી છલાંગ છે.
સમાન વિધેયનું નિરૂપણ માત્ર $\beta$-સમતા આપતું નથી.
અહીં વિધેયનું અચળપણું અને પદની વધારાની
$\beta$-સ્થિરબિંદુ શરત અલગ દર્શાવી છે.}

લૅમ્બડા કલનમાં આપેલા પદનાં સ્થિરબિંદુઓ ગણતાં પદો છે;
તેમનો ઉપયોગ કરીને $\fn{Fac}'$ જેવા પદમાંથી ક્રમગુણિતની
વ્યાખ્યા મેળવી શકાય છે.

\begin{defn}
  \ollabel{defn:Turing-Y}
  \emph{Y-સંયોજક} આ પદ છે:
  \[
  Y \ident (\lambd[ux][x(uux)])(\lambd[ux][x(uux)]).
  \]
\end{defn}

\begin{thm}
  $Y$નો એવો ગુણધર્મ છે કે કોઈ પણ પદ~$g$ માટે
  $Yg \red g(Yg)$ છે. તેથી
  $Yg$ હંમેશા $g$નું સ્થિરબિંદુ છે.
\end{thm}

\begin{proof}
  $(\lambd[ux][x(uux)])$ને સંક્ષેપમાં $U$ કહીએ; ત્યારે
  $Y \ident UU$. હવે
  \begin{align*}
    Y g &\ident (\lambd[ux][x(uux)])U\,g \\
    &\red (\lambd[x][x(UUx)])g\\
    & \red g(UUg) \ident g(Yg).
  \end{align*}
  $g(Yg)$ અને $Yg$ બંને $g(Yg)$માં લઘુકૃત થાય છે,
  તેથી $g(Yg) \equal[\beta] Yg$; આમ $Yg$ એ $g$નું
  સ્થિરબિંદુ છે.
\end{proof}

અલબત્ત, $Yg$ લઘુકરણીય પદ હોવાથી લઘુકરણ અનંત સુધી
ચાલુ રહી શકે છે:
\begin{align*}
  Y g &\red g (Y g) \\
    &\red g (g (Y g)) \\
    &\red g(g (g (Y g)))\\
    &\ldots
\end{align*}
તેથી $Yg$ને જાણે $g$ને પોતે જ અનંત વખત લાગુ કરવાથી
મળતું પદ ગણી શકીએ. તેના પર $g$ એક વધુ વખત લાગુ
કરીએ તો, કહીએ તો, કંઈ વધારાનું થતું નથી: $Yg$ પર
અનંત વખત લાગુ થયેલા $g$ પર ફરી $g$ લાગુ કરીએ તોય
અનંત વખતનો જ અનુપ્રયોગ રહે છે. એટલે $Yg$ પર
$g$નો અનંત અનુપ્રયોગ જ મળે છે.

ધ્યાન આપો કે $Yg$થી શરૂ થતો ઉપરનો $\beta$-લઘુકરણક્રમ
અનંત છે. તેથી $Yg$ને કોઈ પદ પર લાગુ કરીને $(Yg)N$
વિચારીએ તો આ પદ પણ અનંત જુદાં પદોમાં લઘુકૃત થઈ શકે:
$(g(Yg))N$, $(g(g(Yg)))N$, \dots. તેમ છતાં લઘુકરણનો
કોઈ \emph{બીજો} ક્રમ સામાન્ય સ્વરૂપે સમાપ્ત થાય એવું શક્ય છે.

દાખલા તરીકે ક્રમગુણિત લો. $\fn{Fac}$ને $Y\,\fn{Fac}'$
તરીકે, એટલે $\fn{Fac'}$ના સ્થિરબિંદુ તરીકે, વ્યાખ્યાયિત
કરો. ત્યારે:
\begin{align*}
  \fn{Fac}\,\num 3
  &\red Y\, \fn{Fac}' \, \num 3 \\
  &\red \fn{Fac}' (Y \,\fn{Fac}') \, \num 3 \\
  &\ident (\lambd[x][\lambd[n][\fn{IsZero} \, n\, \num 1\,
      (\fn{Mult}\, n\, (x (\fn{Pred}\, n)))]]) \, \fn{Fac} \, \num 3\\
  & \red \fn{IsZero} \, \num 3\, \num 1\,
      (\fn{Mult}\, \num 3\, (\fn{Fac} (\fn{Pred}\, \num 3))) \\
  &\red  \fn{Mult}\, \num 3 \, (\fn{Fac} \, \num 2).
  \intertext{તે જ રીતે,}
  \fn{Fac}\,\num 2
  &\red  \fn{Mult}\, \num 2 \, (\fn{Fac} \, \num 1) \\
  \fn{Fac}\,\num 1
  &\red  \fn{Mult}\, \num 1 \, (\fn{Fac} \, \num 0)
  \intertext{પણ}
  \fn{Fac}\,\num 0
  &\red \fn{Fac}' (Y \,\fn{Fac}') \, \num 0 \\
  &\ident (\lambd[x][\lambd[n][\fn{IsZero} \, n\, \num 1\,
      (\fn{Mult}\, n\, (x (\fn{Pred}\, n)))]]) \, \fn{Fac} \, \num 0\\
  & \red \fn{IsZero} \, \num 0\, \num 1\,
      (\fn{Mult}\, \num 0\, (\fn{Fac} (\fn{Pred}\, \num 0))).\\
  &\red \num 1.
  \intertext{આથી એકંદરે}
  \fn{Fac}\,\num 3
  &\red  \fn{Mult}\, \num 3 \,
  (\fn{Mult} \, \num 2\, (\fn{Mult} \, \num 1\, \num 1)).
\end{align*}

$\fn{Fac'}$ માટે જે થયું તે કોઈ પણ પુનરાવર્તી વ્યાખ્યા
માટે થાય છે. ધારો કે આપણી પાસે નીચેનું પુનરાવર્તી
સમીકરણ છે:
\begin{align*}
g\,x_1\dots x_n & \equal[\beta] N
\intertext{અહીં $N$માં $g$ અને $x_1$, \dots,~$x_n$ આવી શકે છે. ત્યારે
હંમેશા એવું પદ $G \ident (Y \lambd[g][\lambd[x_1\dots x_n][N]])$
હોય છે કે}
G\,x_1 \dots x_n & \equal[\beta] \Subst{N}{G}{g}.
\intertext{કારણ કે સ્થિરબિંદુ પ્રમેયથી,}
G \ident (Y \lambd[g][\lambd[x_1\dots x_n][N]]) & \red \lambd[g][\lambd[x_1\dots x_n][N]](Y \lambd[g][\lambd[x_1\dots x_n][N]])\\
& \ident (\lambd[g][\lambd[x_1\dots x_n][N]])\,G
\intertext{અને તેથી}
G\,x_1 \dots x_n
& \red (\lambd[g][\lambd[x_1\dots x_n][N]])\,G\,x_1\dots x_n\\
& \red (\lambd[x_1\dots x_n][\Subst{N}{G}{g}])\,x_1\dots x_n\\
  & \red \Subst{N}{G}{g}.
\end{align*}

\olref{defn:Turing-Y}નો $Y$-સંયોજક એલન ટ્યુરિંગે આપ્યો હતો.
અલોન્ઝો ચર્ચે બીજો પ્રકાર સૂચવ્યો હતો, જેને આપણે~$Y_C$
કહીશું:
\[
Y_C \ident \lambd[g][(\lambd[x][g(xx)])(\lambd[x][g(xx)])].
\]
ચર્ચનો સંયોજક ટ્યુરિંગના સંયોજક કરતાં આ અર્થમાં થોડો
નબળો છે: $Y_Cg \equal[\beta] g(Y_Cg)$ છે, પરંતુ
$Y_Cg \bred g(Y_Cg)$ નથી.
\sourcecorrection{OLLAM-058}{સ્થિર અંગ્રેજી મૂળમાં
ચર્ચના સંયોજકના આ સરખામણીવાક્યમાં ભૂલથી
$Yg$ લખાયું છે. તે તો અગાઉના પ્રમેય મુજબ
$g(Yg)$માં લઘુકૃત થાય છે. અહીં ઉદ્દેશિત
$Y_Cg$ મૂક્યું છે.}
$V$ને $\lambd[x][g(xx)]$ પદ ગણો; ત્યારે
$Y_C \ident \lambd[g][VV]$. હવે
\begin{align*}
  VV & \ident (\lambd[x][g(xx)])V \red g(VV)
  \text{ અને તેથી}\\
  Y_C g & \ident (\lambd[g][VV])g \red VV \red g(VV), \text{ પરંતુ}\\
  g(Y_C g) & \ident g((\lambd[g][VV])g) \red g(VV).
\end{align*}
અન્ય શબ્દોમાં, $Y_Cg$ અને $g(Y_Cg)$ બંને એક જ પદ
$g(VV)$માં લઘુકૃત થાય છે; એટલે $Y_Cg \equal[\beta]
g(Y_Cg)$. ઉપયોગોમાં ઘણી વાર આટલું પૂરતું છે.
% END OLP-0379
\endgroup

\begingroup
\makeatletter\def\ol@sectioncs{section}\makeatother

% BEGIN OLP-0380
\olfileid{lam}{ldf}{min}
\olsection{લઘુત્તમીકરણ}
\phantomsection\label{guunit:OLP-0380}


સામાન્ય પુનરાવર્તી વિધેયો એવાં વિધેયો છે, જે મૂળભૂત
વિધેયો $\Zero$, $\Succ$, $\Proj{n}{i}$માંથી સંયોજન,
આદિમ પુનરાવર્તન અને નિયમિત લઘુત્તમીકરણ વડે મેળવી
શકાય છે. બધા સામાન્ય પુનરાવર્તી વિધેયો
લૅમ્બડા દ્વારા વ્યાખ્યાયિત છે એવું બતાવવા, લૅમ્બડા દ્વારા વ્યાખ્યાયિત વિધેય પર નિયમિત લઘુત્તમીકરણ લાગુ કરતાં
મળતું કોઈ પણ વિધેય પોતે લૅમ્બડા દ્વારા વ્યાખ્યાયિત છે,
એ બતાવવું પડશે.

\begin{lem}
  \ollabel{lem:min} જો $f(x_1, \dots, x_k, y)$
  નિયમિત અને લૅમ્બડા દ્વારા વ્યાખ્યાયિત હોય, તો
  \[
  g(x_1, \dots, x_k) = \umin{y}{f(x_1,\dots,x_k, y) = 0}
  \]
  વડે વ્યાખ્યાયિત $g$ પણ લૅમ્બડા દ્વારા વ્યાખ્યાયિત છે.
\end{lem}

\begin{proof}
  ધારો કે લૅમ્બડા પદ~$F$ નિયમિત વિધેય
  $f(\vec x, y)$ને $\lambda$-વ્યાખ્યાયિત કરે છે.
  $g$ને લૅમ્બડા દ્વારા વ્યાખ્યાયિત કરવા આપણે શોધ-વિધેય
  અને સ્થિરબિંદુ સંયોજક વાપરીશું:
  \begin{align*}
    \fn{Search} & \ident \lambd[g][\lambd[f\,\vec{x}\,y][
        \fn{IsZero} (f\, \vec{x}\, y)\, y\, (g\, \vec{x} (\fn{Succ}\, y))]]\\
    H & \ident \lambd[\vec x][(Y \, \fn{Search}) F\, \vec{x}\, \num{0}],
  \end{align*}
  જ્યાં $Y$ કોઈ પણ સ્થિરબિંદુ સંયોજક છે.
  \sourcecorrection{OLLAM-060}{સ્થિર અંગ્રેજી મૂળની
  લઘુત્તમીકરણ-પ્રમેયમાં મેળવવાનું વિધેય~$g$ છે,
  પરંતુ પુરાવાની શરૂઆત અને અંતે ભૂલથી~$h$ આવે છે.
  અહીં વિધેય માટે~$g$ રાખ્યું છે; તેનું નિરૂપણ કરતું
  લૅમ્બડા પદ~$H$ મૂળ મુજબ જ રાખ્યું છે.}
  \sourcecorrection{OLLAM-061}{સ્થિર અંગ્રેજી મૂળમાં
  $\fn{Search}$ના પદમાં $(g\,\vec{x}(\fn{Succ}\,y))$
  આહ્વાનનું અંતિમ કૌંસ ગાયબ છે. અહીં તે ઉમેર્યું છે,
  જેથી $\fn{IsZero}$ની ત્રીજી શાખા યોગ્ય રીતે બંધ થાય.}
  અનૌપચારિક રીતે, $\fn{Search}$ પોતાનો સંદર્ભ લેતું
  વિધેય છે: $y$થી શરૂ કરીને $f\,\vec x\,y$
  શૂન્ય છે કે નહીં તે તપાસે છે; હોય તો $y$ પરત કરે
  છે, નહીંતર પોતાને $\fn{Succ}\,y$ સાથે ફરી
  બોલાવે છે. તેથી $(Y \, \fn{Search}) F
  \num{n_1}\dots\num{n_k}\,\num{0}$ એવું નાનામાં
  નાનું~$m$ પરત કરે છે કે જેના માટે $f(n_1,
  \dots, n_k, m) = 0$.

  ખાસ કરીને ધ્યાન આપો:
  \begin{align*}
    (Y \, \fn{Search}) F \num{n_1}
    \dots\num{n_k}\, \num{m} & \red \num{m}
    \intertext{જો $f(n_1, \dots,
      n_k, m) = 0$ હોય; અથવા}
    & \red (Y \, \fn{Search}) F\, \num{n_1} \dots\num{n_k}\, \num{m+1}
    \intertext{અન્યથા. $f$ નિયમિત હોવાથી કોઈક $y$ માટે
      $f(n_1, \dots, n_k, y) = 0$ થાય છે, તેથી}
    (Y \, \fn{Search}) F \num{n_1} \dots\num{n_k}\,\num{0}
    & \red \num{g(n_1, \dots, n_k)}.
    \end{align*}
\end{proof}


\begin{prop}
  દરેક સામાન્ય પુનરાવર્તી વિધેય લૅમ્બડા દ્વારા વ્યાખ્યાયિત છે.
\end{prop}

\begin{proof}
 \olref[prf]{lem:basic} મુજબ બધા મૂળભૂત વિધેયો
 લૅમ્બડા દ્વારા વ્યાખ્યાયિત છે; અને \olref[prf]{lem:comp},
 \olref[prf]{lem:prim} તથા \olref{lem:min} મુજબ
 લૅમ્બડા દ્વારા વ્યાખ્યાયિત વિધેયોનો ગણ સંયોજન,
 આદિમ પુનરાવર્તન અને નિયમિત લઘુત્તમીકરણ હેઠળ
 બંધ છે.
\end{proof}
% END OLP-0380
\endgroup

\begingroup
\makeatletter\def\ol@sectioncs{section}\makeatother

% BEGIN OLP-0381
\olfileid{lam}{ldf}{par}
\olsection{આંશિક પુનરાવર્તી વિધેયો લૅમ્બડા દ્વારા વ્યાખ્યાયિત છે}
\phantomsection\label{guunit:OLP-0381}


આંશિક પુનરાવર્તી વિધેયો મૂળભૂત વિધેયોમાંથી
સંયોજન, આદિમ પુનરાવર્તન અને અનિયત મર્યાદા
સુધીની શોધક્રિયા દ્વારા મળે છે. સામાન્ય પુનરાવર્તી
વિધેયોથી તેમનો તફાવત એ છે કે અનિયત મર્યાદા સુધીની
શોધમાં વપરાતાં વિધેયો નિયમિત હોવાં જરૂરી નથી.
નિયમિતતાની શરત ન હોવાથી લઘુત્તમીકરણ દ્વારા
વ્યાખ્યાયિત વિધેય ક્યારેક અવ્યાખ્યાયિત રહી શકે છે.

પ્રથમ નજરે એવું લાગી શકે કે (સંપૂર્ણ) સામાન્ય પુનરાવર્તી
વિધેયો બધા લૅમ્બડા દ્વારા વ્યાખ્યાયિત છે એવું બતાવતી એ જ
રીતથી બધા આંશિક પુનરાવર્તી વિધેયો પણ
લૅમ્બડા દ્વારા વ્યાખ્યાયિત છે એવું સાબિત થઈ જશે.
દાખલા તરીકે, $f$ અને $g$ને અનુક્રમે $F$ અને~$G$
પદો લૅમ્બડા દ્વારા વ્યાખ્યાયિત કરતાં હોય, તો $f$નું $g$ સાથેનું
સંયોજન $\lambd[x][F (G x)]$ વડે લૅમ્બડા દ્વારા વ્યાખ્યાયિત
થાય છે. પરંતુ વિધેયો આંશિક હોય ત્યારે આમાં
મુશ્કેલી આવે છે. $g(x)$ અવ્યાખ્યાયિત હોય ત્યારે
$G x$નું સામાન્ય સ્વરૂપ હોતું નથી. મોટાભાગના
કિસ્સામાં $F (G x)$નું પણ સામાન્ય સ્વરૂપ હોતું
નથી, અને આપણને એવું જ જોઈએ છે. પરંતુ
$F$ એ $\lambd[x][\lambd[y][y]]$ હોય તે કિસ્સો
વિચારો: ત્યારે $F (G x)$નું સામાન્ય સ્વરૂપ
$\lambd[y][y]$ છે.

આ મુશ્કેલી દૂર કરી શકાય છે. બધા આંશિક પુનરાવર્તી
વિધેયોને એ રીતે લૅમ્બડા દ્વારા વ્યાખ્યાયિત કરી શકાય છે કે
અવ્યાખ્યાયિત મૂલ્યોનું નિરૂપણ સામાન્ય સ્વરૂપ
વિનાનાં પદો કરે. જોકે, આવી રીતો સામાન્ય પુનરાવર્તી
વિધેયો માટે આપણે લીધેલી રીત કરતાં થોડી વધુ
જટિલ અને ઓછી સહજ છે. અહીં પ્રમેય પુરાવા વિના
નોંધીએ છીએ:

\begin{thm}
  બધા આંશિક પુનરાવર્તી વિધેયો લૅમ્બડા દ્વારા વ્યાખ્યાયિત છે.
\end{thm}
% END OLP-0381
\endgroup

\begingroup
\makeatletter\def\ol@sectioncs{section}\makeatother

% BEGIN OLP-0382
\olfileid{lam}{ldf}{ldr}
\olsection{લૅમ્બડા દ્વારા વ્યાખ્યાયિત વિધેયો પુનરાવર્તી છે}
\phantomsection\label{guunit:OLP-0382}

\sourcecorrection{OLLAM-062}{સ્થિર અંગ્રેજી મૂળમાં આ
વિભાગની અધ્યાય-ઓળખ~\texttt{dfl} છે, જ્યારે
અધ્યાયની ઓળખ~\texttt{ldf} અને અન્ય વિભાગોની
ઓળખ પણ~\texttt{ldf} છે. આંતરિક સંદર્ભો માટે
અહીં~\texttt{ldf} કર્યું છે.}

બધા આંશિક પુનરાવર્તી વિધેયો લૅમ્બડા દ્વારા વ્યાખ્યાયિત
છે એટલું જ નહીં, તેનો વિપરીત દાવો પણ સાચો છે.
એટલે બધા લૅમ્બડા દ્વારા વ્યાખ્યાયિત વિધેયો આંશિક
પુનરાવર્તી છે.

\begin{thm}
  \ollabel{thm:lambda-computable} આંશિક વિધેય~$f$
  લૅમ્બડા દ્વારા વ્યાખ્યાયિત હોય તો તે આંશિક પુનરાવર્તી છે.
\end{thm}

\begin{proof}
  પુરાવાની માત્ર રૂપરેખા આપીશું. પહેલા આપણે
  $\lambd$-પદોનું અંકગણિતીકરણ કરીએ છીએ: અનુક્રમોને
  અભાજ્ય સંખ્યાઓની ઘાતોના ગુણાકાર તરીકે સંકેતિત કરવાની પ્રચલિત રીતથી
  $\lambd$-પદોને વ્યવસ્થિત રીતે ગોડેલ સંખ્યાઓ આપીએ
  છીએ. પછી આંશિક પુનરાવર્તી વિધેય
  $\fn{normalize}(t)$ વ્યાખ્યાયિત કરીએ છીએ. તે
  લૅમ્બડા પદની ગોડેલ સંખ્યા~$t$ને આર્ગ્યુમેન્ટ
  તરીકે લે છે; પદનું સામાન્ય સ્વરૂપ હોય તો તેની
  ગોડેલ સંખ્યા પરત કરે છે અને ન હોય તો
  અવ્યાખ્યાયિત રહે છે. ત્યાર પછી $\fn{toChurch}$
  અને $\fn{fromChurch}$ એવાં બે આંશિક પુનરાવર્તી
  વિધેયો વ્યાખ્યાયિત કરીએ છીએ, જે પ્રાકૃતિક સંખ્યાઓને
  અનુરૂપ ચર્ચ અંકોની ગોડેલ સંખ્યાઓમાં અને તેમાંથી
  પાછી ફેરવે છે.

  આ પુનરાવર્તી વિધેયો વડે $f$ને આંશિક પુનરાવર્તી
  વિધેય તરીકે વ્યાખ્યાયિત કરી શકીએ છીએ. કોઈ
  $\lambd$-પદ~$F$ વિધેય~$f$ને લૅમ્બડા દ્વારા વ્યાખ્યાયિત
  કરે છે. $f(n_1, \dots, n_k)$ ગણવા પહેલા
  $\fn{toChurch}(n_i)$ વડે અનુરૂપ ચર્ચ અંકોની
  ગોડેલ સંખ્યાઓ મેળવો. તેમને $\Gn{F}$ સાથે જોડીને
  $F \num{n_1}\dots\num{n_k}$ પદની ગોડેલ સંખ્યા
  બનાવો. હવે આ ગોડેલ સંખ્યા પર $\fn{normalize}$
  લાગુ કરો. $f(n_1, \dots, n_k)$ વ્યાખ્યાયિત હોય
  તો $F \num{n_1}\dots\num{n_k}$નું સામાન્ય સ્વરૂપ
  હોય છે (અને તે ચર્ચ અંક જ હોવો જોઈએ); અન્યથા
  તેનું સામાન્ય સ્વરૂપ હોતું નથી (એટલે
  \[\fn{normalize}(\Gn{F\num{n_1}\dots\num{n_k}})\]
  અવ્યાખ્યાયિત છે). છેવટે સામાન્ય સ્વરૂપે આવેલા
  પદની ગોડેલ સંખ્યા પર $\fn{fromChurch}$ લાગુ કરો.
\end{proof}
% END OLP-0382
\endgroup


\OLEndChapterHook


\begin{editorial}
આગળનો સંબંધિત અંશ મૂળના સક્રિય આયાતક્રમમાં નથી. તેને અહીં સ્પષ્ટ વૈકલ્પિક સંદર્ભમાં જાળવ્યો છે.
\end{editorial}
\begingroup
\makeatletter\def\ol@sectioncs{section}\makeatother

% BEGIN OLP-0650
\olfileid{lam}{ldf}{lst}
\olsection{યાદીઓ}
\phantomsection\label{guunit:OLP-0650}


યાદી $[M_1, M_2, \ldots, M_n]$ આ પ્રમાણે
વ્યાખ્યાયિત કરી શકાય છે:
\[
  [M_1, M_2, \ldots, M_n] \ident \lambd[sz][s M_1 (s M_2 (\ldots (s M_n z)))]
\]
અનૌપચારિક રીતે, યાદીનું નિરૂપણ કરતું પદ
તે યાદીનું ``સંચયક'' વિધેય છે. તે બે પદો
લે છે. પહેલું પદ નવો ઘટક અને અત્યાર સુધીનું
સંચિત મૂલ્ય લઈને નવું સંચિત મૂલ્ય આપતું
વિધેય છે; બીજું પદ શરૂઆતનું સંચિત મૂલ્ય છે.
આ રીતે તે એક પછી એક ઘટક સંચિત કરીને અંતે
પરિણામ આપે છે.
\sourcecorrection{OLMD-187}{મૂળ ગદ્યમાં
પહેલા વિધેયને સંચિત મૂલ્ય અને નવો ઘટક
એવા ક્રમમાં આર્ગ્યુમેન્ટ મળતાં કહે છે.
પ્રદર્શિત પદ $s\,M_i\,a$માં નવો ઘટક
પહેલો અને સંચિત મૂલ્ય બીજું છે; અહીં
ગદ્યનો ક્રમ પદ મુજબ રાખ્યો છે.}

\begin{digress}
  વિધેયાત્મક પ્રોગ્રામિંગથી પરિચિત હો
  તો આ વ્યાખ્યા \emph{જમણી ફોલ્ડ} જેવી
  લાગશે: યાદી તેના ઘટકોના જમણી ફોલ્ડ
  વિધેય તરીકે વ્યાખ્યાયિત થાય છે.
\end{digress}

આ સંકેતનના આધારે કેટલાંક ઉપયોગી વિધેયો
સરળતાથી વ્યાખ્યાયિત કરી શકાય છે:
\begin{align*}
  \fn{Sum} & \ident
  \lambd[l][l \, (\lambd[xa][\fn{Add} \, x \, a])\,  \num{0}]\\
  \fn{Len} & \ident
  \lambd[l][l \, (\lambd[xa][\fn{Add} \, a \, \num{1}]) \num{0}]
\end{align*}

$\fn{Sum}$ ચર્ચ અંકોની યાદીનો સરવાળો
કાઢે છે. યાદી પર સંચયન કરતાં તેનું
શરૂઆતનું મૂલ્ય $\num{0}$ છે અને દરેક
ઘટક માટે નવું મૂલ્ય $\fn{Add} \, x\, a$
છે; અહીં $x$ હાલનો ઘટક અને $a$ સંચિત
મૂલ્ય છે. આખરે બધા ઘટકોનો સરવાળો મળે છે.

\begin{prob}
  $\fn{Len}$ શું કરે છે? સમજાવો.
\end{prob}

\begin{prob}
  યાદીનો ક્રમ ઊલટાવતું વિધેય વ્યાખ્યાયિત કરો.
\end{prob}
% END OLP-0650
\endgroup
% END OLP-0373
\endgroup


\OLEndPartHook
% END OLP-0341
